
\Data\ID{1003021501}
\Updated{2022-08-25 10:08:20}
\Level{A}
\SubArea{Circles}
\LANGen{
\Question{
Given a regular octadecagon \( A_1A_2\dots A_{18} \), determine the measure of the angle \( A_2A_8A_3 \).}
{\( 10^{\circ} \)}{\( 20^{\circ} \)}{\( 30^{\circ} \)}{\( 40^{\circ} \)}{}{}}
\LANGpl{
\Question{
Dany jest osiemnastokąt foremny  \( A_1A_2\dots A_{18} \), oblicz miarę kąta \( A_2A_8A_3 \).}
{\( 10^{\circ} \)}{\( 20^{\circ} \)}{\( 30^{\circ} \)}{\( 40^{\circ} \)}{}{}}
\LANGcs{
\Question{
V pravidelném osmnáctiúhelníku \( A_1A_2\dots A_{18} \) určete velikost úhlu \( A_2A_8A_3 \).}
{\( 10^{\circ} \)}{\( 20^{\circ} \)}{\( 30^{\circ} \)}{\( 40^{\circ} \)}{}{}}
\LANGsk{
\Question{
V pravidelnom osemnásťuholníku  \( A_1A_2\dots A_{18} \) určte veľkosť uhla  \( A_2A_8A_3 \).}
{\( 10^{\circ} \)}{\( 20^{\circ} \)}{\( 30^{\circ} \)}{\( 40^{\circ} \)}{}{}}
\LANGes{
\Question{
Dado el octodecágono regular \( A_1A_2\dots A_{18} \), calcula la medida de los ángulos \( A_2A_8A_3 \).}
{\( 10^{\circ} \)}{\( 20^{\circ} \)}{\( 30^{\circ} \)}{\( 40^{\circ} \)}{}{}}
\End

\Data\ID{1003021508}
\Updated{2022-08-25 10:08:20}
\Level{A}
\SubArea{Circles}
\LANGen{
\Question{
A triangle is inscribed in a circle. Its vertices divide the circle into three arcs whose lengths are in the ratio \( 2:4:9 \). Determine the measures of the interior angles of the triangle.}
{\( 24^{\circ};\ 48^{\circ};\ 108^{\circ} \)}{\( 30^{\circ};\ 40^{\circ};\ 110^{\circ} \)}{\( 48^{\circ};\ 15^{\circ};\ 117^{\circ} \)}{\( 15^{\circ};\ 60^{\circ};\ 105^{\circ} \)}{}{}}
\LANGpl{
\Question{
Trójkąt jest wpisany w okrąg. Jego wierzchołki dzielą okrąg na trzy łuki o długości \( 2:4:9 \).  Oblicz miarę kątów wewnętrznych trójkąta.}
{\( 24^{\circ};\ 48^{\circ};\ 108^{\circ} \)}{\( 30^{\circ};\ 40^{\circ};\ 110^{\circ} \)}{\( 48^{\circ};\ 15^{\circ};\ 117^{\circ} \)}{\( 15^{\circ};\ 60^{\circ};\ 105^{\circ} \)}{}{}}
\LANGcs{
\Question{
Do kružnice je vepsaný trojúhelník tak, že jeho vrcholy dělí kruh na tři oblouky, jejichž délky jsou v poměru \( 2:4:9 \). Určete velikost vnitřních úhlů trojúhelníku.}
{\( 24^{\circ};\ 48^{\circ};\ 108^{\circ} \)}{\( 30^{\circ};\ 40^{\circ};\ 110^{\circ} \)}{\( 48^{\circ};\ 15^{\circ};\ 117^{\circ} \)}{\( 15^{\circ};\ 60^{\circ};\ 105^{\circ} \)}{}{}}
\LANGsk{
\Question{
Do kružnice je vpísaný trojuholník tak, že jeho vrcholy delia kružnicu na oblúky. Dĺžky oblúkov sú v pomere  \( 2:4:9 \). Nájdite veľkosť  vnútorných uhlov trojuholníka.}
{\( 24^{\circ};\ 48^{\circ};\ 108^{\circ} \)}{\( 30^{\circ};\ 40^{\circ};\ 110^{\circ} \)}{\( 48^{\circ};\ 15^{\circ};\ 117^{\circ} \)}{\( 15^{\circ};\ 60^{\circ};\ 105^{\circ} \)}{}{}}
\LANGes{
\Question{
Un triángulo se inscribe en una circunferencia. Sus vértices dividen la circunferencia en tres arcos cuyas longitudes están en proporción \( 2:4:9 \). Calcula la medida de los ángulos interiores del triángulo.}
{\( 24^{\circ};\ 48^{\circ};\ 108^{\circ} \)}{\( 30^{\circ};\ 40^{\circ};\ 110^{\circ} \)}{\( 48^{\circ};\ 15^{\circ};\ 117^{\circ} \)}{\( 15^{\circ};\ 60^{\circ};\ 105^{\circ} \)}{}{}}
\End

\Data\ID{1103021502}
\Updated{2022-04-23 05:04:54}
\Level{A}
\SubArea{Circles}
\IMG{1103021502.pdf}
\LANGen{
\Question{
What is the measure of the angle contained by two line segments: the first  joining  numbers \( 8 \) and \( 11 \), and the second joining numbers \( 11 \) and \( 2 \), on a clock face?  (See the picture.)}
{\( 90^{\circ} \)}{\( 100^{\circ} \)}{\( 80^{\circ} \)}{\( 70^{\circ} \)}{}{}}
\LANGpl{
\Question{
Podaj miarę kąta pomiędzy dwoma odcinkami: pierwszy łączy liczby  \( 8 \) i \( 11 \), drugi łączy liczby  \( 11 \) i \( 2 \), na tarczy zegara.  (Spójrz na rysunek.)}
{\( 90^{\circ} \)}{\( 100^{\circ} \)}{\( 80^{\circ} \)}{\( 70^{\circ} \)}{}{}}
\LANGcs{
\Question{
Vypočítejte velikost úhlu, který na hodinách svírají spojnice bodů označených čísly \( 8 \), \( 11 \) a  \( 11 \), \( 2 \).  (viz obrázek)}
{\( 90^{\circ} \)}{\( 100^{\circ} \)}{\( 80^{\circ} \)}{\( 70^{\circ} \)}{}{}}
\LANGsk{
\Question{
Vypočítajte veľkosť uhla, ktorý na hodinovom ciferníku zvierajú spojnice bodov označené číslami  \( 8 \), \( 11 \) a  \( 11 \), \( 2 \).  (Pozri obrázok.)}
{\( 90^{\circ} \)}{\( 100^{\circ} \)}{\( 80^{\circ} \)}{\( 70^{\circ} \)}{}{}}
\LANGes{
\Question{
¿Cuál es la medida del ángulo formado por dos segmentos de un reloj si el primer segmento tiene como origen y extremo \( 8 \) y \( 11 \) y el segundo \( 11 \) y \( 2 \) ?  (Mira la imagen).}
{\( 90^{\circ} \)}{\( 100^{\circ} \)}{\( 80^{\circ} \)}{\( 70^{\circ} \)}{}{}}
\End

\Data\ID{1103021503}
\Updated{2022-04-23 05:04:54}
\Level{A}
\SubArea{Circles}
\IMG{1103021503.pdf}
\LANGen{
\Question{
Determine the measure of the angle contained by two line segments: the first joining numbers \( 7 \) and \( 1 \), and the second joining numbers \( 1 \) and \( 4 \), on a clock face. (See the picture.)}
{\( 45^{\circ} \)}{\( 60^{\circ} \)}{\( 30^{\circ} \)}{\( 90^{\circ} \)}{}{}}
\LANGpl{
\Question{
Podaj miarę kąta pomiędzy dwoma odcinkami: pierwszy łączy liczby \( 7 \) i \( 1 \), a drugi łączy liczby \( 1 \) i \( 4 \), na tarczy zegara. (Spójrz na rysunek.)}
{\( 45^{\circ} \)}{\( 60^{\circ} \)}{\( 30^{\circ} \)}{\( 90^{\circ} \)}{}{}}
\LANGcs{
\Question{
Vypočítejte velikost úhlu, který na hodinách svírají spojnice bodů označené číslicemi  \( 7 \), \( 1 \) a \( 1 \), \( 4 \). (viz obrázek)}
{\( 45^{\circ} \)}{\( 60^{\circ} \)}{\( 30^{\circ} \)}{\( 90^{\circ} \)}{}{}}
\LANGsk{
\Question{
Vypočítajte veľkosť uhla, ktorý na hodinovom ciferníku zvierajú spojnice bodov označené číslami  \( 7 \), \( 1 \) a \( 1 \), \( 4 \). (Pozri obrázok.)}
{\( 45^{\circ} \)}{\( 60^{\circ} \)}{\( 30^{\circ} \)}{\( 90^{\circ} \)}{}{}}
\LANGes{
\Question{
¿Cuál es la medida del ángulo formado por dos segmentos de un reloj si el primer segmento tiene como origen y extremo \( 7 \) y \( 1 \) y el segundo \( 1 \) y \( 4 \) ?  (Mira la imagen).}
{\( 45^{\circ} \)}{\( 60^{\circ} \)}{\( 30^{\circ} \)}{\( 90^{\circ} \)}{}{}}
\End

\Data\ID{1103021504}
\Updated{2022-04-23 05:04:54}
\Level{A}
\SubArea{Circles}
\IMG{1103021504.pdf}
\LANGen{
\Question{
Give the measure of the angle contained by two line segments with the endpoints at the  numbers \( 7 \), \( 8 \) and \( 8 \), \( 10 \) on a clock face. (See the picture.)}
{\( 135^{\circ} \)}{\( 90^{\circ} \)}{\( 45^{\circ} \)}{\( 120^{\circ} \)}{}{}}
\LANGpl{
\Question{
Podaj miarę kąta pomiędzy odcinkami kończącymi się na liczbach  \( 7 \), \( 8 \) i  \( 8 \), \( 10 \) na tarczy zegara. (Spójrz na rysunek.)}
{\( 135^{\circ} \)}{\( 90^{\circ} \)}{\( 45^{\circ} \)}{\( 120^{\circ} \)}{}{}}
\LANGcs{
\Question{
Vypočítejte velikost úhlu, který na hodinách svírají spojnice bodů označené číslicemi  \( 7 \), \( 8 \) a \( 8 \), \( 10 \). (viz obrázek)}
{\( 135^{\circ} \)}{\( 90^{\circ} \)}{\( 45^{\circ} \)}{\( 120^{\circ} \)}{}{}}
\LANGsk{
\Question{
Vypočítajte veľkosť uhla, ktorý na hodinovom ciferníku zvierajú spojnice bodov označené číslami  \( 7 \), \( 8 \) a \( 8 \), \( 10 \). (Pozri obrázok.)}
{\( 135^{\circ} \)}{\( 90^{\circ} \)}{\( 45^{\circ} \)}{\( 120^{\circ} \)}{}{}}
\LANGes{
\Question{
¿Cuál es la medida del ángulo formado por dos segmentos de un reloj si el primer segmento tiene como origen y extremo \( 7 \) y \( 8 \) y el segundo \( 8 \) y \( 10 \) ?  (Mira la imagen).}
{\( 135^{\circ} \)}{\( 90^{\circ} \)}{\( 45^{\circ} \)}{\( 120^{\circ} \)}{}{}}
\End

\Data\ID{1103021505}
\Updated{2022-04-23 05:04:54}
\Level{A}
\SubArea{Circles}
\IMG{1103021505.pdf}
\LANGen{
\Question{
Give the measure of the angle between two line segments: the first connecting  numbers \( 7 \) and \( 11 \), the second connecting numbers \( 3 \) and \( 10 \) on a clock face. (See the picture.)}
{\( 75^{\circ} \)}{\( 60^{\circ} \)}{\( 45^{\circ} \)}{\( 80^{\circ} \)}{}{}}
\LANGpl{
\Question{
Podaj miarę kąta pomiędzy odcinkami: pierwszy łączy liczby \( 7 \) i \( 11 \), drugi łączy liczby \( 3 \) i \( 10 \) na tarczy zegara. (Spójrz na rysunek.)}
{\( 75^{\circ} \)}{\( 60^{\circ} \)}{\( 45^{\circ} \)}{\( 80^{\circ} \)}{}{}}
\LANGcs{
\Question{
Vypočítejte velikost úhlu, který na hodinovém ciferníku svírají spojnice bodů označených číslicemi  \( 7 \), \( 11 \) a \( 3 \), \( 10 \). (viz obrázek)}
{\( 75^{\circ} \)}{\( 60^{\circ} \)}{\( 45^{\circ} \)}{\( 80^{\circ} \)}{}{}}
\LANGsk{
\Question{
Vypočítajte veľkosť uhla, ktorý na hodinovom ciferníku zvierajú spojnice bodov označené číslami  \( 7 \), \( 11 \) a \( 3 \), \( 10 \). (Pozri obrázok.)}
{\( 75^{\circ} \)}{\( 60^{\circ} \)}{\( 45^{\circ} \)}{\( 80^{\circ} \)}{}{}}
\LANGes{
\Question{
¿Cuál es la medida del ángulo formado por dos segmentos de un reloj si el primer segmento tiene como origen y extremo \( 7 \) y \( 11 \) y el segundo \( 3 \) y \( 10 \) ?  (Mira la imagen).}
{\( 75^{\circ} \)}{\( 60^{\circ} \)}{\( 45^{\circ} \)}{\( 80^{\circ} \)}{}{}}
\End

\Data\ID{1103021506}
\Updated{2022-04-23 05:04:54}
\Level{A}
\SubArea{Circles}
\IMG{1103021506.pdf}
\LANGen{
\Question{
Points \( A \) and \( B \) divide the circle \( k \) into two arcs whose lengths are in the ratio \( 5:13 \). Point \( C \) is an interior point of the  longer arc. What is the degree measure of the angle \( ACB \)?}
{\( 50^{\circ} \)}{\( 40^{\circ} \)}{\( 100^{\circ} \)}{\( 20^{\circ} \)}{}{}}
\LANGpl{
\Question{
Punkty  \( A \) i \( B \) dzielą okrąg  \( k \) na dwa łuki o długości  w stosunku  \( 5:13 \). Punkt \( C \) to wewnętrzny punkt dłuższego łuku. Jaka jest miara kąta  \( ACB \)?}
{\( 50^{\circ} \)}{\( 40^{\circ} \)}{\( 100^{\circ} \)}{\( 20^{\circ} \)}{}{}}
\LANGcs{
\Question{
Body \( A \) a \( B \) rozdělují kružnici  \( k \) na dva oblouky, jejichž délky jsou v poměru  \( 5:13 \). Bod \( C \) je vnitřním bodem delšího oblouku. Jakou velikost má úhel \( ACB \)?}
{\( 50^{\circ} \)}{\( 40^{\circ} \)}{\( 100^{\circ} \)}{\( 20^{\circ} \)}{}{}}
\LANGsk{
\Question{
Body \( A \) a \( B \) rozdeľujú kružnicu  \( k \) na dva oblúky, ktorých dĺžky sú v pomere   \( 5:13 \). Bod \( C \) je vnútorným bodom dlhšieho oblúka. Akú veľkosť má uhol  \( ACB \)?}
{\( 50^{\circ} \)}{\( 40^{\circ} \)}{\( 100^{\circ} \)}{\( 20^{\circ} \)}{}{}}
\LANGes{
\Question{
Los puntos \( A \) y \( B \) dividen a la circunferencia \( k \) en dos arcos cuyas longitudes están en proporción \( 5:13 \). El punto \( C \) es un punto interior del arco más largo. Calcula la medida del ángulo \( ACB \).}
{\( 50^{\circ} \)}{\( 40^{\circ} \)}{\( 100^{\circ} \)}{\( 20^{\circ} \)}{}{}}
\End

\Data\ID{1103021507}
\Updated{2020-07-15 03:07:21}
\Level{A}
\SubArea{Circles}
\IMG{1103021507.pdf}
\LANGen{
\Question{
Line segment \( AB \) is a diameter of the circle \( k \). The lengths of the arcs \( AD \) and \( DB \) are in the ratio \( 7:3 \). Find the measure of the angle \( ACD \). (See the picture.)}
{\( 63^{\circ} \)}{\( 27^{\circ} \)}{\( 126^{\circ} \)}{\( 24^{\circ} \)}{}{}}
\LANGpl{
\Question{
Odcinek \( AB \) to średnica okręgu \( k \). Długość łuków \( AD \) i \( DB \) są w stosunku \( 7:3 \). Oblicz miarę kąta  \( ACD \). (Spójrz na rysunek.)}
{\( 63^{\circ} \)}{\( 27^{\circ} \)}{\( 126^{\circ} \)}{\( 24^{\circ} \)}{}{}}
\LANGcs{
\Question{
Úsečka \( AB \) je průměrem kružnice \( k \). Délky oblouků \( AD \) a \( DB \) jsou v poměru \( 7:3 \). Zjistěte velikost úhlu \( ACD \). (viz obrázek)}
{\( 63^{\circ} \)}{\( 27^{\circ} \)}{\( 126^{\circ} \)}{\( 24^{\circ} \)}{}{}}
\LANGsk{
\Question{
Úsečka \( AB \) je priemerom kružnice  \( k \). Pomer dĺžok oblúkov \( AD \) a \( DB \) je v pomere \( 7:3 \). Nájdite veľkosť uhla \( ACD \). (Pozrite obrázok.)}
{\( 63^{\circ} \)}{\( 27^{\circ} \)}{\( 126^{\circ} \)}{\( 24^{\circ} \)}{}{}}
\LANGes{
\Question{
El segmento \( AB \) es el diámetro de la circunferencia \( k \). Las longitudes de los arcos \( AD \) y \( DB \) están en proporción \( 7:3 \). Calcula la medida del ángulo \( ACD \).}
{\( 63^{\circ} \)}{\( 27^{\circ} \)}{\( 126^{\circ} \)}{\( 24^{\circ} \)}{}{}}
\End

\Data\ID{1103021509}
\Updated{2022-04-23 05:04:54}
\Level{A}
\SubArea{Circles}
\IMG{1103021509.pdf}
\LANGen{
\Question{
A regular dodecagon \( ABCDEFGHIJKL \) is inscribed in a circle. Find the measures of all interior angles of the quadrilateral \( ABHJ \). (See the picture.)}
{\( \alpha=120^{\circ};\ \beta=75^{\circ};\ \gamma=60^{\circ};\ \delta=105^{\circ} \)}{\( \alpha=105^{\circ};\ \beta=60^{\circ};\ \gamma=75^{\circ};\ \delta=120^{\circ} \)}{\( \alpha=120^{\circ};\ \beta=30^{\circ};\ \gamma=60^{\circ};\ \delta=105^{\circ} \)}{\( \alpha=105^{\circ};\ \beta=75^{\circ};\ \gamma=75^{\circ};\ \delta=105^{\circ} \)}{}{}}
\LANGpl{
\Question{
Dwunastokąt foremny \( ABCDEFGHIJKL \) jest wpisany w okrąg. Podaj miary kątów wewnętrznych czworokąta  \( ABHJ \). (Spójrz na rysunek.)}
{\( \alpha=120^{\circ};\ \beta=75^{\circ};\ \gamma=60^{\circ};\ \delta=105^{\circ} \)}{\( \alpha=105^{\circ};\ \beta=60^{\circ};\ \gamma=75^{\circ};\ \delta=120^{\circ} \)}{\( \alpha=120^{\circ};\ \beta=30^{\circ};\ \gamma=60^{\circ};\ \delta=105^{\circ} \)}{\( \alpha=105^{\circ};\ \beta=75^{\circ};\ \gamma=75^{\circ};\ \delta=105^{\circ} \)}{}{}}
\LANGcs{
\Question{
Do kružnice je vepsaný pravidelný dvanáctiúhelník \( ABCDEFGHIJKL \). Vypočítejte velikost vnitřních úhlů tětivového čtyřúhelníku \( ABHJ \). (viz obrázek)}
{\( \alpha=120^{\circ};\ \beta=75^{\circ};\ \gamma=60^{\circ};\ \delta=105^{\circ} \)}{\( \alpha=105^{\circ};\ \beta=60^{\circ};\ \gamma=75^{\circ};\ \delta=120^{\circ} \)}{\( \alpha=120^{\circ};\ \beta=30^{\circ};\ \gamma=60^{\circ};\ \delta=105^{\circ} \)}{\( \alpha=105^{\circ};\ \beta=75^{\circ};\ \gamma=75^{\circ};\ \delta=105^{\circ} \)}{}{}}
\LANGsk{
\Question{
Do kružnice je vpísaný pravidelný 12 uholník  \( ABCDEFGHIJKL \). Vypočítajte veľkosti vnútorných uhlov tetivového štvoruholníka  \( ABHJ \). (Pozrite obrázok.)}
{\( \alpha=120^{\circ};\ \beta=75^{\circ};\ \gamma=60^{\circ};\ \delta=105^{\circ} \)}{\( \alpha=105^{\circ};\ \beta=60^{\circ};\ \gamma=75^{\circ};\ \delta=120^{\circ} \)}{\( \alpha=120^{\circ};\ \beta=30^{\circ};\ \gamma=60^{\circ};\ \delta=105^{\circ} \)}{\( \alpha=105^{\circ};\ \beta=75^{\circ};\ \gamma=75^{\circ};\ \delta=105^{\circ} \)}{}{}}
\LANGes{
\Question{
Un dodecágono regular \( ABCDEFGHIJKL \) se inscribe en una circunferencia. Calcula la medida de todos los ángulos interiores del cuadrilátero \( ABHJ \).}
{\( \alpha=120^{\circ};\ \beta=75^{\circ};\ \gamma=60^{\circ};\ \delta=105^{\circ} \)}{\( \alpha=105^{\circ};\ \beta=60^{\circ};\ \gamma=75^{\circ};\ \delta=120^{\circ} \)}{\( \alpha=120^{\circ};\ \beta=30^{\circ};\ \gamma=60^{\circ};\ \delta=105^{\circ} \)}{\( \alpha=105^{\circ};\ \beta=75^{\circ};\ \gamma=75^{\circ};\ \delta=105^{\circ} \)}{}{}}
\End

\Data\ID{1103021510}
\Updated{2022-04-23 05:04:54}
\Level{A}
\SubArea{Circles}
\IMG{1103021510.pdf}
\LANGen{
\Question{
A regular nonagon \( ABCDEFGHI \) is inscribed in a circle. Calculate the measures of all interior angles of the quadrilateral \( ABEH \). (See the picture.)}
{\( \alpha=120^{\circ};\ \beta=100^{\circ};\ \gamma=60^{\circ};\ \delta=80^{\circ} \)}{\( \alpha=100^{\circ};\ \beta=120^{\circ};\ \gamma=60^{\circ};\ \delta=80^{\circ} \)}{\( \alpha=100^{\circ};\ \beta=100^{\circ};\ \gamma=80^{\circ};\ \delta=60^{\circ} \)}{\( \alpha=110^{\circ};\ \beta=130^{\circ};\ \gamma=70^{\circ};\ \delta=50^{\circ} \)}{}{}}
\LANGpl{
\Question{
Dziewięciokąt foremny \( ABCDEFGHI \) jest wpisany w okrąg. Oblicz miary wszystkich kątów wewnętrznych czworokąta \( ABEH \). (Spójrz na rysunek.)}
{\( \alpha=120^{\circ};\ \beta=100^{\circ};\ \gamma=60^{\circ};\ \delta=80^{\circ} \)}{\( \alpha=100^{\circ};\ \beta=120^{\circ};\ \gamma=60^{\circ};\ \delta=80^{\circ} \)}{\( \alpha=100^{\circ};\ \beta=100^{\circ};\ \gamma=80^{\circ};\ \delta=60^{\circ} \)}{\( \alpha=110^{\circ};\ \beta=130^{\circ};\ \gamma=70^{\circ};\ \delta=50^{\circ} \)}{}{}}
\LANGcs{
\Question{
Do kružnice je vepsaný pravidelný devítiúhelník \( ABCDEFGHI \). Vypočítejte velikost vnitřních úhlů tětivového čtyřúhelníku \( ABEH \). (viz obrázek)}
{\( \alpha=120^{\circ};\ \beta=100^{\circ};\ \gamma=60^{\circ};\ \delta=80^{\circ} \)}{\( \alpha=100^{\circ};\ \beta=120^{\circ};\ \gamma=60^{\circ};\ \delta=80^{\circ} \)}{\( \alpha=100^{\circ};\ \beta=100^{\circ};\ \gamma=80^{\circ};\ \delta=60^{\circ} \)}{\( \alpha=110^{\circ};\ \beta=130^{\circ};\ \gamma=70^{\circ};\ \delta=50^{\circ} \)}{}{}}
\LANGsk{
\Question{
Do kružnice je vpísaný pravidelný 9 uholník  \( ABCDEFGHI \). Vypočítajte veľkosti vnútorných uhlov tetivového štvoruholníka \( ABEH \). (Pozrite obrázok.)}
{\( \alpha=120^{\circ};\ \beta=100^{\circ};\ \gamma=60^{\circ};\ \delta=80^{\circ} \)}{\( \alpha=100^{\circ};\ \beta=120^{\circ};\ \gamma=60^{\circ};\ \delta=80^{\circ} \)}{\( \alpha=100^{\circ};\ \beta=100^{\circ};\ \gamma=80^{\circ};\ \delta=60^{\circ} \)}{\( \alpha=110^{\circ};\ \beta=130^{\circ};\ \gamma=70^{\circ};\ \delta=50^{\circ} \)}{}{}}
\LANGes{
\Question{
Un eneágono regular \( ABCDEFGHI \) se inscribe en una circunferencia. Calcula la medida de todos los ángulos interiores del cuadrilátero \( ABEH \).}
{\( \alpha=120^{\circ};\ \beta=100^{\circ};\ \gamma=60^{\circ};\ \delta=80^{\circ} \)}{\( \alpha=100^{\circ};\ \beta=120^{\circ};\ \gamma=60^{\circ};\ \delta=80^{\circ} \)}{\( \alpha=100^{\circ};\ \beta=100^{\circ};\ \gamma=80^{\circ};\ \delta=60^{\circ} \)}{\( \alpha=110^{\circ};\ \beta=130^{\circ};\ \gamma=70^{\circ};\ \delta=50^{\circ} \)}{}{}}
\End

\Data\ID{1103021512}
\Updated{2022-04-23 05:04:54}
\Level{A}
\SubArea{Circles}
\IMG{1103021512.pdf}
\LANGen{
\Question{
In the triangle \( ABC \), \( a=10\,\mathrm{cm} \), \( b=8\,\mathrm{cm} \), \( c=12\,\mathrm{cm} \). Point \( D \) is the foot of the altitude from the vertex \( C \)  (see the picture.) What is the radius of the circumcircle of the triangle \( DBC \)?}
{\( 5\,\mathrm{cm} \)}{\( 4\,\mathrm{cm} \)}{\( 6\,\mathrm{cm} \)}{\( 8\,\mathrm{cm} \)}{}{}}
\LANGpl{
\Question{
W trójkącie \( ABC \), \( a=10\,\mathrm{cm} \), \( b=8\,\mathrm{cm} \), \( c=12\,\mathrm{cm} \). Punkt \( D \) to punkt początkowy wysokości z wierzchołka  \( C \). (Spójrz na rysunek.)  Ile wynosi promień okręgu opisanego na trójkącie \( DBC \)?}
{\( 5\,\mathrm{cm} \)}{\( 4\,\mathrm{cm} \)}{\( 6\,\mathrm{cm} \)}{\( 8\,\mathrm{cm} \)}{}{}}
\LANGcs{
\Question{
V trojúhelníku \( ABC \), \( a=10\,\mathrm{cm} \), \( b=8\,\mathrm{cm} \), \( c=12\,\mathrm{cm} \). Bod \( D \) je patou výšky na stranu \( c \) (viz obrázek). Jaký poloměr má kružnice opsaná trojúhelníku \( DBC \)?}
{\( 5\,\mathrm{cm} \)}{\( 4\,\mathrm{cm} \)}{\( 6\,\mathrm{cm} \)}{\( 8\,\mathrm{cm} \)}{}{}}
\LANGsk{
\Question{
V trojuholníku \( ABC \), \( a=10\,\mathrm{cm} \), \( b=8\,\mathrm{cm} \), \( c=12\,\mathrm{cm} \). Bod \( D \) je päta výšky na stranu \( c \). (Pozrite obrázok.) Aký polomer má kružnica opísaná trojuholníku  \( DBC \)?}
{\( 5\,\mathrm{cm} \)}{\( 4\,\mathrm{cm} \)}{\( 6\,\mathrm{cm} \)}{\( 8\,\mathrm{cm} \)}{}{}}
\LANGes{
\Question{
Dado el triángulo \( ABC \), \( a=10\,\mathrm{cm} \), \( b=8\,\mathrm{cm} \), \( c=12\,\mathrm{cm} \). El punto \( D \) es el pie de la altura desde el vértice \( C \). (Mira la imagen). ¿Cuánto mide el radio de la circunferencia circunscrita del triángulo \( DBC \)?}
{\( 5\,\mathrm{cm} \)}{\( 4\,\mathrm{cm} \)}{\( 6\,\mathrm{cm} \)}{\( 8\,\mathrm{cm} \)}{}{}}
\End

\Data\ID{2000005901}
\Updated{2022-05-16 11:05:07}
\Level{A}
\SubArea{Circles}
\IMG{2000005901-figure0.pdf}
\LANGen{
\Question{
What is the measure of the angle made by two line-segments marked in the clock face (see the picture)? One segment is connecting points \(10\) and \(2\), second segment is connecting points \(5\) and \(2\).}
{\(75^{\circ}\)}{\(65^{\circ}\)}{\(85^{\circ}\)}{\(55^{\circ}\)}{}{}}
\LANGpl{
\Question{
Jaka jest miara kąta wyznaczana przez dwa odcinki linii zaznaczone na tarczy zegara (patrz rysunek)? Jeden odcinek łączy punkty \(10\) i \(2\), drugi odcinek łączy punkty \(5\) i \(2\).}
{\(75^{\circ}\)}{\(65^{\circ}\)}{\(85^{\circ}\)}{\(55^{\circ}\)}{}{}}
\LANGcs{
\Question{
Jaká je velikost úhlu, který svírají dvě úsečky vyznačené na ciferníku hodin na obrázku? Jedna úsečka spojuje čísla \(10\) a \(2\), druhá úsečka spojuje čísla \(5\) a \(2\).}
{\(75^{\circ}\)}{\(65^{\circ}\)}{\(85^{\circ}\)}{\(55^{\circ}\)}{}{}}
\LANGsk{
\Question{
Aká je veľkosť uhla, ktorý na ciferníku hodín zvierajú spojnice bodov označené číslami  \(10\), \(2\) a \(5\), \(2\), (pozri obrázok)?}
{\(75^{\circ}\)}{\(65^{\circ}\)}{\(85^{\circ}\)}{\(55^{\circ}\)}{}{}}
\LANGes{
\Question{
¿Cuál es la medida del ángulo formado por dos segmentos: el primer segmento marcado por \(10\) y \(2\), y el segundo marcado por \(5\) y \(2\) en un reloj?}
{\(75^{\circ}\)}{\(65^{\circ}\)}{\(85^{\circ}\)}{\(55^{\circ}\)}{}{}}
\End

\Data\ID{2000005902}
\Updated{2022-05-16 11:05:07}
\Level{A}
\SubArea{Circles}
\IMG{2000005901-figure1.pdf}
\LANGen{
\Question{
What is the measure of the angle made by two line-segments marked in the clock face (see the picture)? One segment is connecting points \(1\) and \(3\), second segment is connecting points \(5\) and \(3\).}
{\( 120^{\circ}\)}{\( 135^{\circ}\)}{\( 60^{\circ}\)}{\( 240^{\circ}\)}{}{}}
\LANGpl{
\Question{
Jaka jest miara kąta wyznaczana przez dwa odcinki linii zaznaczone na tarczy zegara (patrz rysunek)? Jeden segment łączy punkty \(1\) i \(3\), drugi segment łączy punkty \(5\) i \(3\).}
{\( 120^{\circ}\)}{\( 135^{\circ}\)}{\( 60^{\circ}\)}{\( 240^{\circ}\)}{}{}}
\LANGcs{
\Question{
Jaká je velikost úhlu, který svírají dvě úsečky znázorněné na ciferníku na obrázku? Jedna úsečka spojuje číslice \(1\) a \(3\), druhá úsečka spojuje číslice \(5\) a \(3\).}
{\( 120^{\circ}\)}{\( 135^{\circ}\)}{\( 60^{\circ}\)}{\( 240^{\circ}\)}{}{}}
\LANGsk{
\Question{
Vypočítajte veľkosť uhla, ktorý na ciferníku hodín zvierajú spojnice bodov označené číslami  \(1\), \(3\) a \(5\), \(3\), (pozri obrázok).}
{\( 120^{\circ}\)}{\( 135^{\circ}\)}{\( 60^{\circ}\)}{\( 240^{\circ}\)}{}{}}
\LANGes{
\Question{
¿Cuál es la medida del ángulo formado por dos segmentos: el primer segmento marcado por \(1\) y \(3\), y el segundo marcado por \(5\) y \(3\) en un reloj?}
{\( 120^{\circ}\)}{\( 135^{\circ}\)}{\( 60^{\circ}\)}{\( 240^{\circ}\)}{}{}}
\End

\Data\ID{2000005903}
\Updated{2022-05-16 11:05:07}
\Level{A}
\SubArea{Circles}
\IMG{2000005901-figure2.pdf}
\LANGen{
\Question{
What is the measure of the angle made by two line-segments marked in the clock face (see the picture)? One segment is connecting points \(7\) and \(1\), second segment is connecting points \(5\) and \(10\).}
{\(105^{\circ}\)}{\(120^{\circ}\)}{\(115^{\circ}\)}{\(75^{\circ}\)}{}{}}
\LANGpl{
\Question{
Jaka jest miara kąta wyznaczana przez dwa odcinki linii zaznaczone na tarczy zegara (patrz rysunek)? Jeden odcinek łączy punkty \(7\) i \(1\), drugi odcinek łączy punkty \(5\) i \(10\).}
{\(105^{\circ}\)}{\(120^{\circ}\)}{\(115^{\circ}\)}{\(75^{\circ}\)}{}{}}
\LANGcs{
\Question{
Jaká je velikost úhlu, který svírají dvě úsečky vyznačené na ciferníku na obrázku? Jedna úsečka spojuje číslice \(7\) a \(1\), druhá úsečka spojuje číslice \(5\) a \(10\).}
{\(105^{\circ}\)}{\(120^{\circ}\)}{\(115^{\circ}\)}{\(75^{\circ}\)}{}{}}
\LANGsk{
\Question{
Aká je veľkosť uhla, ktorý na ciferníku hodín zvierajú spojnice bodov označené číslami  \(7\), \(1\) a \(5\), \(10\), (pozri obrázok)?}
{\(105^{\circ}\)}{\(120^{\circ}\)}{\(115^{\circ}\)}{\(75^{\circ}\)}{}{}}
\LANGes{
\Question{
¿Cuál es la medida del ángulo formado por dos segmentos: el primer segmento marcado por \(7\) y \(1\), y el segundo marcado por \(5\) y \(10\) en un reloj?}
{\(105^{\circ}\)}{\(120^{\circ}\)}{\(115^{\circ}\)}{\(75^{\circ}\)}{}{}}
\End

\Data\ID{2010012801}
\Updated{2022-08-25 10:08:09}
\Level{A}
\SubArea{Circles}
\LANGen{
\Question{
A triangle is inscribed in a circle. Its vertices divide the circle into three arcs whose lengths are in the ratio \( 3:4:5 \). Determine the measures of the interior angles of the triangle.}
{\( 45^{\circ};\ 60^{\circ};\ 75^{\circ} \)}{\( 20^{\circ};\ 60^{\circ};\ 100^{\circ} \)}{\( 20^{\circ};\ 40^{\circ};\ 120^{\circ} \)}{\( 50^{\circ};\ 60^{\circ};\ 70^{\circ} \)}{}{}}
\LANGpl{
\Question{
Trójkąt wpisano w okrąg. Jego wierzchołki dzielą okrąg na trzy łuki, których długości są w stosunku \( 3:4:5 \). Wyznacz miary kątów wewnętrznych trójkąta.}
{\( 45^{\circ};\ 60^{\circ};\ 75^{\circ} \)}{\( 20^{\circ};\ 60^{\circ};\ 100^{\circ} \)}{\( 20^{\circ};\ 40^{\circ};\ 120^{\circ} \)}{\( 50^{\circ};\ 60^{\circ};\ 70^{\circ} \)}{}{}}
\LANGcs{
\Question{
Do kružnice je vepsaný trojúhelník. Jeho vrcholy rozdělují kružnici na tři oblouky, jejichž délky jsou v poměru \( 3:4:5 \). Vypočtěte velikost vnitřních úhlů trojúhelníku.}
{\( 45^{\circ};\ 60^{\circ};\ 75^{\circ} \)}{\( 20^{\circ};\ 60^{\circ};\ 100^{\circ} \)}{\( 20^{\circ};\ 40^{\circ};\ 120^{\circ} \)}{\( 50^{\circ};\ 60^{\circ};\ 70^{\circ} \)}{}{}}
\LANGsk{
\Question{
Do kruhu je vpísaný trojuholník. Jeho vrcholy rozdeľujú kruh na tri oblúky, ktorých dĺžky sú v pomere \( 3:4:5 \). Vypočítajte rozmery vnútorných uhlov trojuholníka.}
{\( 45^{\circ};\ 60^{\circ};\ 75^{\circ} \)}{\( 20^{\circ};\ 60^{\circ};\ 100^{\circ} \)}{\( 20^{\circ};\ 40^{\circ};\ 120^{\circ} \)}{\( 50^{\circ};\ 60^{\circ};\ 70^{\circ} \)}{}{}}
\LANGes{
\Question{
Un triángulo se inscribe en una circunferencia. Sus vértices dividen la circunferencia en tres arcos cuyas longitudes están en proporción \( 3:4:5 \). Calcula la medida de los ángulos interiores del triángulo.}
{\( 45^{\circ};\ 60^{\circ};\ 75^{\circ} \)}{\( 20^{\circ};\ 60^{\circ};\ 100^{\circ} \)}{\( 20^{\circ};\ 40^{\circ};\ 120^{\circ} \)}{\( 50^{\circ};\ 60^{\circ};\ 70^{\circ} \)}{}{}}
\End

\Data\ID{2010012802}
\Updated{2022-08-25 10:08:09}
\Level{A}
\SubArea{Circles}
\IMG{2010012801-figure0.pdf}
\LANGen{
\Question{
Determine the measure of the angle contained by two line segments: the first joining numbers \( 7 \) and \( 9 \), and the second joining numbers \( 7 \) and \( 2 \), on a clock face. (See the picture.)}
{\( 75^{\circ}\)}{\( 30^{\circ}\)}{\( 55^{\circ}\)}{\( 60^{\circ}\)}{}{}}
\LANGpl{
\Question{
Określ miarę kąta utworzonego przez dwa odcinki: pierwszy łączący liczby \( 7 \) i \( 9 \) i drugi łączący liczby \( 7 \) i \( 2 \) na tarczy zegara. (Patrz rysunek.)}
{\( 75^{\circ}\)}{\( 30^{\circ}\)}{\( 55^{\circ}\)}{\( 60^{\circ}\)}{}{}}
\LANGcs{
\Question{
Určete velikost úhlu, který na hodinovém ciferníku svírají dvě úsečky. První z nich spojuje čísla \( 7 \) a \( 9 \) a druhá čísla \( 7 \) a \( 2 \) (Viz obrázek.).}
{\( 75^{\circ}\)}{\( 30^{\circ}\)}{\( 55^{\circ}\)}{\( 60^{\circ}\)}{}{}}
\LANGsk{
\Question{
Vypočítajte veľkosť uhla, ktorý na hodinovom ciferníku zvierajú spojnice bodov označené číslami  \( 7 \),  \( 9 \), a  \( 7 \), \( 2 \), (Pozri obrázok.)}
{\( 75^{\circ}\)}{\( 30^{\circ}\)}{\( 55^{\circ}\)}{\( 60^{\circ}\)}{}{}}
\LANGes{
\Question{
¿Cuál es la medida del ángulo formado por dos segmentos de un reloj si el primer segmento tiene como origen y extremo \( 7 \) y \( 9 \) y el segundo \( 7 \) y \( 2 \)? (Mira la imagen).}
{\( 75^{\circ}\)}{\( 30^{\circ}\)}{\( 55^{\circ}\)}{\( 60^{\circ}\)}{}{}}
\End

\Data\ID{2010012803}
\Updated{2022-08-25 10:08:09}
\Level{A}
\SubArea{Circles}
\IMG{2010012801-figure1.pdf}
\LANGen{
\Question{
What is the measure of the angle contained by two line segments: the first  joining  numbers \( 6 \) and \( 10 \), and the second joining numbers \( 6 \) and \( 4 \), on a clock face?  (See the picture.)}
{\( 90^{\circ}\)}{\( 60^{\circ}\)}{\( 85^{\circ}\)}{\( 95^{\circ}\)}{}{}}
\LANGpl{
\Question{
Jaka jest miara kąta utworzonego przez dwa odcinki: pierwszy łączący liczby \( 6 \) i \( 10 \) i drugi łączący liczby \( 6 \) i \( 4 \) na tarczy zegara? (Zobacz zdjęcie.)}
{\( 90^{\circ}\)}{\( 60^{\circ}\)}{\( 85^{\circ}\)}{\( 95^{\circ}\)}{}{}}
\LANGcs{
\Question{
Určete velikost úhlu, který na hodinovém ciferníku svírají dvě úsečky. První z nich spojuje čísla \( 6 \) a \( 10 \), druhá čísla \( 6 \) a \( 4 \) (Viz obrázek.).}
{\( 90^{\circ}\)}{\( 60^{\circ}\)}{\( 85^{\circ}\)}{\( 95^{\circ}\)}{}{}}
\LANGsk{
\Question{
Vypočítajte veľkosť uhla, ktorý na hodinovom ciferníku zvierajú spojnice bodov označené číslami  \( 6 \),  \( 10 \) a \( 6 \),  \( 4 \).  (Pozri obrázok.)}
{\( 90^{\circ}\)}{\( 60^{\circ}\)}{\( 85^{\circ}\)}{\( 95^{\circ}\)}{}{}}
\LANGes{
\Question{
¿Cuál es la medida del ángulo formado por dos segmentos de un reloj si el primer segmento tiene como origen y extremo \( 6 \) y \( 10 \) y el segundo  \( 6 \) y \( 4 \)?  (Mira la imagen).}
{\( 90^{\circ}\)}{\( 60^{\circ}\)}{\( 85^{\circ}\)}{\( 95^{\circ}\)}{}{}}
\End

\Data\ID{2010012804}
\Updated{2022-08-25 10:08:12}
\Level{A}
\SubArea{Circles}
\IMG{2010012801-figure2.pdf}
\LANGen{
\Question{
Give the measure of the angle contained by two line segments with the endpoints at the numbers \( 6 \), \( 5 \) and \( 6 \), \( 9 \) on a clock face. (See the picture.)}
{\( 120^{\circ}\)}{\( 135^{\circ}\)}{\( 100^{\circ}\)}{\( 60^{\circ}\)}{}{}}
\LANGpl{
\Question{
Podaj miarę kąta utworzonego przez dwa odcinki o końcach w liczbach \( 6 \), \( 5 \) i \( 6 \), \( 9 \) na tarczy zegara. (Patrz rysunek.)}
{\( 120^{\circ}\)}{\( 135^{\circ}\)}{\( 100^{\circ}\)}{\( 60^{\circ}\)}{}{}}
\LANGcs{
\Question{
Určete velikost úhlu, který na hodinách svírají spojnice bodů označené číslicemi \( 6 \), \( 5 \) a \( 6 \), \( 9 \) (Viz obrázek.).}
{\( 120^{\circ}\)}{\( 135^{\circ}\)}{\( 100^{\circ}\)}{\( 60^{\circ}\)}{}{}}
\LANGsk{
\Question{
Vypočítajte veľkosť uhla, ktorý na hodinovom ciferníku zvierajú spojnice bodov označené číslami   \( 6 \), \( 5 \) a \( 6 \), \( 9 \). (Pozri obrázok.)}
{\( 120^{\circ}\)}{\( 135^{\circ}\)}{\( 100^{\circ}\)}{\( 60^{\circ}\)}{}{}}
\LANGes{
\Question{
¿Cuál es la medida del ángulo formado por dos segmentos de un reloj si el primer segmento tiene como origen y extremo \( 6 \) y \( 5 \) y  el segundo \( 6 \) y \( 9 \)? (Mira la imagen.)}
{\( 120^{\circ}\)}{\( 135^{\circ}\)}{\( 100^{\circ}\)}{\( 60^{\circ}\)}{}{}}
\End

\Data\ID{2010012805}
\Updated{2022-08-25 10:08:12}
\Level{A}
\SubArea{Circles}
\IMG{2010012801-figure3_0.pdf}
\LANGen{
\Question{
Give the measure of the angle between two line segments: the first connecting  numbers \( 6 \) and \( 11 \), the second connecting numbers \(  8\) and \( 2\) on a clock face. (See the picture.)}
{\( 75^{\circ}\)}{\( 30^{\circ}\)}{\( 60^{\circ}\)}{\( 45^{\circ}\)}{}{}}
\LANGpl{
\Question{
Podaj miarę kąta między dwoma odcinkami: pierwszy łączący liczby \( 6 \) i \( 11 \), drugi łączący liczby \( 8\) i \( 2\) na tarczy zegara. (Patrz rysunek.)}
{\( 75^{\circ}\)}{\( 30^{\circ}\)}{\( 60^{\circ}\)}{\( 45^{\circ}\)}{}{}}
\LANGcs{
\Question{
Určete velikost úhlu, který na hodinovém ciferníku svírají dvě úsečky. První z nich spojuje čísla \( 6 \) a \( 11 \), druhá čísla \(  8\) a \( 2\) (Viz obrázek.).}
{\( 75^{\circ}\)}{\( 30^{\circ}\)}{\( 60^{\circ}\)}{\( 45^{\circ}\)}{}{}}
\LANGsk{
\Question{
Vypočítajte veľkosť uhla, ktorý na hodinovom ciferníku zvierajú spojnice bodov označené číslami  \( 6 \),  \( 11 \) a \(  8\),  \( 2\). (Pozri obrázok.)}
{\( 75^{\circ}\)}{\( 30^{\circ}\)}{\( 60^{\circ}\)}{\( 45^{\circ}\)}{}{}}
\LANGes{
\Question{
¿Cuál es la medida del ángulo formado por dos segmentos de un reloj si el primer segmento tiene como origen y extremo \( 6 \) y \( 11 \) y el segundo \(  8\) y \( 2\)? (Mira la imagen.)}
{\( 75^{\circ}\)}{\( 30^{\circ}\)}{\( 60^{\circ}\)}{\( 45^{\circ}\)}{}{}}
\End

\Data\ID{2010012806}
\Updated{2022-08-25 10:08:12}
\Level{A}
\SubArea{Circles}
\IMG{2010012801-figure4.pdf}
\LANGen{
\Question{
Points \( A \) and \( B \) divide the circle \( k \) into two arcs whose lengths are in the ratio \( 3:12 \). Point \( C \) is an interior point of the longer arc. What is the degree measure of the angle \( ACB \)?}
{\( 36^{\circ}\)}{\( 72^{\circ}\)}{\( 24^{\circ}\)}{\( 45^{\circ}\)}{}{}}
\LANGpl{
\Question{
Punkty \( A \) i \( B \) dzielą okrąg \( k \) na dwa łuki, których stosunek długości wynosi \( 3:12 \). Punkt \( C \) jest wewnętrznym punktem dłuższego łuku. Jaka jest miara kąta \( ACB \)?}
{\( 36^{\circ}\)}{\( 72^{\circ}\)}{\( 24^{\circ}\)}{\( 45^{\circ}\)}{}{}}
\LANGcs{
\Question{
Body \( A \) a \( B \) rozdělují kružnici \( k \) na dva oblouky, jejichž délky jsou v poměru \( 3:12 \). Bod \( C \) je vnitřním bodem delšího z obou oblouků. Jaká je velikost úhlu \( ACB \)?}
{\( 36^{\circ}\)}{\( 72^{\circ}\)}{\( 24^{\circ}\)}{\( 45^{\circ}\)}{}{}}
\LANGsk{
\Question{
Body \( A \) a \( B \) rozdeľujú kružnicu \( k \) na dva oblúky, ktorých dĺžky sú v pomere \( 3:12 \). Bod \( C \) je vnútorným bodom dlhšieho oblúka. Akú veľkosť má uhol \( ACB \)?}
{\( 36^{\circ}\)}{\( 72^{\circ}\)}{\( 24^{\circ}\)}{\( 45^{\circ}\)}{}{}}
\LANGes{
\Question{
Los puntos \( A \) y \( B \) dividen a la circunferencia \( k \) en dos arcos cuyas longitudes están en proporción \( 3:12 \). El punto \( C \) es un punto interior del arco más largo. Calcula la medida del ángulo \( ACB \).}
{\( 36^{\circ}\)}{\( 72^{\circ}\)}{\( 24^{\circ}\)}{\( 45^{\circ}\)}{}{}}
\End

\Data\ID{2010012807}
\Updated{2022-08-25 10:08:12}
\Level{A}
\SubArea{Circles}
\IMG{2010012801-figure5.pdf}
\LANGen{
\Question{
A regular dodecagon \( ABCDEFGHIJKL \) is inscribed in a circle. Find the measures of all interior angles of the quadrilateral \( BFIL \). (See the picture.)}
{\( \alpha=90^{\circ};\ \beta=75^{\circ};\ \gamma=90^{\circ};\ \delta=105^{\circ} \)}{\( \alpha=90^{\circ};\ \beta=60^{\circ};\ \gamma=80^{\circ};\ \delta=130^{\circ} \)}{\( \alpha=80^{\circ};\ \beta=75^{\circ};\ \gamma=90^{\circ};\ \delta=115^{\circ} \)}{\( \alpha=90^{\circ};\ \beta=105^{\circ};\ \gamma=90^{\circ};\ \delta=105^{\circ} \)}{}{}}
\LANGpl{
\Question{
W okrąg wpisany jest dwunastokąt foremny \( ABCDEFGHIJKL \). Znajdź miary wszystkich kątów wewnętrznych czworokąta \( BFIL \). (Patrz rysunek.)}
{\( \alpha=90^{\circ};\ \beta=75^{\circ};\ \gamma=90^{\circ};\ \delta=105^{\circ} \)}{\( \alpha=90^{\circ};\ \beta=60^{\circ};\ \gamma=80^{\circ};\ \delta=130^{\circ} \)}{\( \alpha=80^{\circ};\ \beta=75^{\circ};\ \gamma=90^{\circ};\ \delta=115^{\circ} \)}{\( \alpha=90^{\circ};\ \beta=105^{\circ};\ \gamma=90^{\circ};\ \delta=105^{\circ} \)}{}{}}
\LANGcs{
\Question{
Do kružnice je vepsaný pravidelný dvanáctiúhelník \( ABCDEFGHIJKL \). Vypočtěte velikost vnitřních úhlů tětivového čtyřúhelníku \( BFIL \) (Viz obrázek.).}
{\( \alpha=90^{\circ};\ \beta=75^{\circ};\ \gamma=90^{\circ};\ \delta=105^{\circ} \)}{\( \alpha=90^{\circ};\ \beta=60^{\circ};\ \gamma=80^{\circ};\ \delta=130^{\circ} \)}{\( \alpha=80^{\circ};\ \beta=75^{\circ};\ \gamma=90^{\circ};\ \delta=115^{\circ} \)}{\( \alpha=90^{\circ};\ \beta=105^{\circ};\ \gamma=90^{\circ};\ \delta=105^{\circ} \)}{}{}}
\LANGsk{
\Question{
Do kružnice je vpísaný pravidelný 12 uholník \( ABCDEFGHIJKL \).  Vypočítajte veľkosti vnútorných uhlov tetivového štvoruholníka \( BFIL \). (Pozri obrázok.)}
{\( \alpha=90^{\circ};\ \beta=75^{\circ};\ \gamma=90^{\circ};\ \delta=105^{\circ} \)}{\( \alpha=90^{\circ};\ \beta=60^{\circ};\ \gamma=80^{\circ};\ \delta=130^{\circ} \)}{\( \alpha=80^{\circ};\ \beta=75^{\circ};\ \gamma=90^{\circ};\ \delta=115^{\circ} \)}{\( \alpha=90^{\circ};\ \beta=105^{\circ};\ \gamma=90^{\circ};\ \delta=105^{\circ} \)}{}{}}
\LANGes{
\Question{
Un dodecágono regular \( ABCDEFGHIJKL \) se inscribe en una circunferencia. Calcula la medida de todos los ángulos interiores del cuadrilátero \( BFIL \).}
{\( \alpha=90^{\circ};\ \beta=75^{\circ};\ \gamma=90^{\circ};\ \delta=105^{\circ} \)}{\( \alpha=90^{\circ};\ \beta=60^{\circ};\ \gamma=80^{\circ};\ \delta=130^{\circ} \)}{\( \alpha=80^{\circ};\ \beta=75^{\circ};\ \gamma=90^{\circ};\ \delta=115^{\circ} \)}{\( \alpha=90^{\circ};\ \beta=105^{\circ};\ \gamma=90^{\circ};\ \delta=105^{\circ} \)}{}{}}
\End

\Data\ID{2010012808}
\Updated{2022-08-25 10:08:12}
\Level{A}
\SubArea{Circles}
\IMG{2010012801-figure6.pdf}
\LANGen{
\Question{
A regular nonagon \( ABCDEFGHI \) is inscribed in a circle. Calculate the measures of all interior angles of the quadrilateral \( BDGI \). (See the picture.)}
{\( \alpha=100^{\circ};\ \beta=80^{\circ};\ \gamma=80^{\circ};\ \delta=100^{\circ} \)}{\( \alpha=110^{\circ};\ \beta=80^{\circ};\ \gamma=80^{\circ};\ \delta=90^{\circ} \)}{\( \alpha=110^{\circ};\ \beta=70^{\circ};\ \gamma=70^{\circ};\ \delta=110^{\circ} \)}{\( \alpha=120^{\circ};\ \beta=80^{\circ};\ \gamma=80^{\circ};\ \delta=120^{\circ} \)}{}{}}
\LANGpl{
\Question{
Dziewięcikąt foremny \( ABCDEFGHI \) jest wpisany w okrąg. Oblicz miary wszystkich kątów wewnętrznych czworokąta \(BDGI \). (Patrz rysunek.)}
{\( \alpha=100^{\circ};\ \beta=80^{\circ};\ \gamma=80^{\circ};\ \delta=100^{\circ} \)}{\( \alpha=110^{\circ};\ \beta=80^{\circ};\ \gamma=80^{\circ};\ \delta=90^{\circ} \)}{\( \alpha=110^{\circ};\ \beta=70^{\circ};\ \gamma=70^{\circ};\ \delta=110^{\circ} \)}{\( \alpha=120^{\circ};\ \beta=80^{\circ};\ \gamma=80^{\circ};\ \delta=120^{\circ} \)}{}{}}
\LANGcs{
\Question{
Do kružnice je vepsaný pravidelný devítiúhelník \( ABCDEFGHI \). Vypočtěte velikost všech vnitřních úhlů tětivového čtyřúhelníku \( BDGI \) (Viz obrázek.).}
{\( \alpha=100^{\circ};\ \beta=80^{\circ};\ \gamma=80^{\circ};\ \delta=100^{\circ} \)}{\( \alpha=110^{\circ};\ \beta=80^{\circ};\ \gamma=80^{\circ};\ \delta=90^{\circ} \)}{\( \alpha=110^{\circ};\ \beta=70^{\circ};\ \gamma=70^{\circ};\ \delta=110^{\circ} \)}{\( \alpha=120^{\circ};\ \beta=80^{\circ};\ \gamma=80^{\circ};\ \delta=120^{\circ} \)}{}{}}
\LANGsk{
\Question{
Do kružnice je vpísaný pravidelný 9 uholník  \( ABCDEFGHI \). Vypočítajte veľkosti vnútorných uhlov tetivového štvoruholníka \( BDGI \). (Pozri obrázok.)}
{\( \alpha=100^{\circ};\ \beta=80^{\circ};\ \gamma=80^{\circ};\ \delta=100^{\circ} \)}{\( \alpha=110^{\circ};\ \beta=80^{\circ};\ \gamma=80^{\circ};\ \delta=90^{\circ} \)}{\( \alpha=110^{\circ};\ \beta=70^{\circ};\ \gamma=70^{\circ};\ \delta=110^{\circ} \)}{\( \alpha=120^{\circ};\ \beta=80^{\circ};\ \gamma=80^{\circ};\ \delta=120^{\circ} \)}{}{}}
\LANGes{
\Question{
Un eneágono regular \( ABCDEFGHI \) se inscribe en una circunferencia. Calcula la medida de todos los ángulos interiores del cuadrilátero \( BDGI \).}
{\( \alpha=100^{\circ};\ \beta=80^{\circ};\ \gamma=80^{\circ};\ \delta=100^{\circ} \)}{\( \alpha=110^{\circ};\ \beta=80^{\circ};\ \gamma=80^{\circ};\ \delta=90^{\circ} \)}{\( \alpha=110^{\circ};\ \beta=70^{\circ};\ \gamma=70^{\circ};\ \delta=110^{\circ} \)}{\( \alpha=120^{\circ};\ \beta=80^{\circ};\ \gamma=80^{\circ};\ \delta=120^{\circ} \)}{}{}}
\End

\Data\ID{2010012810}
\Updated{2022-08-25 10:08:12}
\Level{A}
\SubArea{Circles}
\IMG{2010012801-figure8_0.pdf}
\LANGen{
\Question{
In the triangle \( KLM \), \( k=10\,\mathrm{cm} \), \( l=8\,\mathrm{cm} \), \( m=12\,\mathrm{cm} \). Point \( N \) is the foot of the altitude from the vertex \( K \) (See the picture.) What is the radius of the circumcircle of the triangle \( KLN \)?}
{\( 6\,\mathrm{cm} \)}{\( 5\,\mathrm{cm} \)}{\( 7\,\mathrm{cm} \)}{\( 8\,\mathrm{cm} \)}{}{}}
\LANGpl{
\Question{
W trójkącie \( KLM \), \( k=10\,\mathrm{cm} \), \( l=8\,\mathrm{cm} \), \( m=12\,\mathrm{cm } \). Punkt \( N \) to podstawa wysokości z wierzchołka \( K \) (Patrz rysunek.). Jaki jest promień okręgu opisanego w trójkącie \( KLN \)?}
{\( 6\,\mathrm{cm} \)}{\( 5\,\mathrm{cm} \)}{\( 7\,\mathrm{cm} \)}{\( 8\,\mathrm{cm} \)}{}{}}
\LANGcs{
\Question{
Je dán trojúhelník \( KLM \), \( k=10\,\mathrm{cm} \), \( l=8\,\mathrm{cm} \), \( m=12\,\mathrm{cm} \). Bod \( N \) je patou výšky vedené z bodu \( K \) (Viz obrázek.) Jaký je poloměr kružnice opsané trojúhelníku \( KLN \)?}
{\( 6\,\mathrm{cm} \)}{\( 5\,\mathrm{cm} \)}{\( 7\,\mathrm{cm} \)}{\( 8\,\mathrm{cm} \)}{}{}}
\LANGsk{
\Question{
V trojuholníku \( KLM \), \( k=10\,\mathrm{cm} \), \( l=8\,\mathrm{cm} \), \( m=12\,\mathrm{cm} \). Bod \( N \) je päta výšky na stranu \( k \). (Pozrite obrázok.) Aký polomer má kružnica opísaná trojuholníku \( KLN \)?}
{\( 6\,\mathrm{cm} \)}{\( 5\,\mathrm{cm} \)}{\( 7\,\mathrm{cm} \)}{\( 8\,\mathrm{cm} \)}{}{}}
\LANGes{
\Question{
Dado el triángulo \( KLM \), \( k=10\,\mathrm{cm} \), \( l=8\,\mathrm{cm} \), \( m=12\,\mathrm{cm} \). El punto \( N \) es el pie de la altura desde el vértice\( K \) (Mira la imagen.). ¿Cuánto mide el radio de la circunferencia circunscrita del triángulo \( KLN \)?}
{\( 6\,\mathrm{cm} \)}{\( 5\,\mathrm{cm} \)}{\( 7\,\mathrm{cm} \)}{\( 8\,\mathrm{cm} \)}{}{}}
\End

\Data\ID{1003021603}
\Updated{2022-08-25 10:08:20}
\Level{B}
\SubArea{Circles}
\IMG{1103021601-figure0.pdf}
\LANGen{
\Question{
Which of the following  formulas expresses the area of a regular decagon inscribed in a circle of radius \( r \)? (See the picture.)}
{\( 10r^2\sin18^{\circ}\cos18^{\circ} \)}{\( 10r^2\sin36^{\circ}\cos36^{\circ} \)}{\( 5r^2\sin36^{\circ} \)}{\( 5r^2\sin18^{\circ} \)}{}{}}
\LANGpl{
\Question{
Które z podanych wyrażeń określa obraz dziesięciokąta foremnego wpisanego w okrąg o promieniu \( r \) (patrz rysunek)?}
{\( 10r^2\sin18^{\circ}\cos18^{\circ} \)}{\( 10r^2\sin36^{\circ}\cos36^{\circ} \)}{\( 5r^2\sin36^{\circ} \)}{\( 5r^2\sin18^{\circ} \)}{}{}}
\LANGcs{
\Question{
Která z uvedených rovnic vyjadřuje obsah pravidelného desetiúhelníku vepsaného do kružnice o poloměru \( r \) (viz obrázek)?}
{\( 10r^2\sin18^{\circ}\cos18^{\circ} \)}{\( 10r^2\sin36^{\circ}\cos36^{\circ} \)}{\( 5r^2\sin36^{\circ} \)}{\( 5r^2\sin18^{\circ} \)}{}{}}
\LANGsk{
\Question{
Ktorý z uvedených vzorcov vyjadruje obsah pravidelného desaťuholníka vpísaného do kružnice  s polomerom  \( r \)? (Pozri obrázok.)}
{\( 10r^2\sin18^{\circ}\cos18^{\circ} \)}{\( 10r^2\sin36^{\circ}\cos36^{\circ} \)}{\( 5r^2\sin36^{\circ} \)}{\( 5r^2\sin18^{\circ} \)}{}{}}
\LANGes{
\Question{
¿Cuál de las siguientes fórmulas expresa el área de un decágono inscrito en una circunferencia con  radio \( r \)?}
{\( 10r^2\sin18^{\circ}\cos18^{\circ} \)}{\( 10r^2\sin36^{\circ}\cos36^{\circ} \)}{\( 5r^2\sin36^{\circ} \)}{\( 5r^2\sin18^{\circ} \)}{}{}}
\End

\Data\ID{1003021607}
\Updated{2022-08-25 10:08:20}
\Level{B}
\SubArea{Circles}
\LANGen{
\Question{
Consider a right-angled triangle \( ABC \) with the right angle at the vertex \( C \). Calculate the measure of the angle \( CAB \), if the side \( b=9\,\mathrm{cm} \) and the radius of the circumscribed circle \( r=6\,\mathrm{cm} \). Round the result to one decimal place.}
{\( 41.4^{\circ} \)}{\( 48.6^{\circ} \)}{\( 36.9^{\circ} \)}{\( 48.2^{\circ} \)}{}{}}
\LANGpl{
\Question{
Dany jest trójkąt prostokątny \( ABC \) o kącie prostym na wierzchołku  \( C \). Oblicz miarę kąta  \( CAB \), jeśli bok \( b=9\,\mathrm{cm} \) oraz promień okręgu opisanego na trójkącie  \( r=6\,\mathrm{cm} \). Zaokrągli wynik do jednego miejsca po przecinku.}
{\( 41{,}4^{\circ} \)}{\( 48{,}6^{\circ} \)}{\( 36{,}9^{\circ} \)}{\( 48{,}2^{\circ} \)}{}{}}
\LANGcs{
\Question{
Je dán pravoúhlý trojúhelník \( ABC \) s pravým úhlem u vrcholu \( C \). Vypočítejte velikost úhlu \( CAB \), pokud strana \( b=9\,\mathrm{cm} \) a poloměr kružnice opsané je \( r=6\,\mathrm{cm} \). Výsledek zaokrouhlete na jedno desetinné místo.}
{\( 41{,}4^{\circ} \)}{\( 48{,}6^{\circ} \)}{\( 36{,}9^{\circ} \)}{\( 48{,}2^{\circ} \)}{}{}}
\LANGsk{
\Question{
Daný je pravouhlý trojuholník \( ABC \) s pravým uhlom pri vrchole \( C \). Vypočítajte veľkosť uhla \( CAB \), ak strana \( b=9\,\mathrm{cm} \) a polomer kružnice opísanej danému trojuholníku \( r=6\,\mathrm{cm} \). Výsledok uveďte s presnosťou na jedno desatinné miesto.}
{\( 41{,}4^{\circ} \)}{\( 48{,}6^{\circ} \)}{\( 36{,}9^{\circ} \)}{\( 48{,}2^{\circ} \)}{}{}}
\LANGes{
\Question{
Dado el triángulo rectángulo \( ABC \) siendo \( C \) el vértice del ángulo recto. Calcula la medida del ángulo \( CAB \), suponiendo que \( b=9\,\mathrm{cm} \) y el radio de la circunferencia circunscrita \( r=6\,\mathrm{cm} \). Redondea el resultado a un decimal.}
{\( 41.4^{\circ} \)}{\( 48.6^{\circ} \)}{\( 36.9^{\circ} \)}{\( 48.2^{\circ} \)}{}{}}
\End

\Data\ID{1003077112}
\Updated{2020-07-15 03:07:21}
\Level{B}
\SubArea{Circles}
\LANGen{
\Question{
The length of a circular arc with the central angle of \( 3.5 \) radians is \( 82\,\mathrm{cm} \). Calculate the radius of the corresponding circle. Round the result to two decimal places.}
{\( 23.43\,\mathrm{cm} \)}{\( 287.00\,\mathrm{cm} \)}{\( 1.59\,\mathrm{cm} \)}{\( 4217.40\,\mathrm{cm} \)}{}{}}
\LANGpl{
\Question{
Długość łuku kołowego o kącie środkowym równym \( 3{,}5 \) radianów wynosi \( 82\,\mathrm{cm} \). Oblicz promień okręgu, z którego pochodzi łuk. Zaokrąglij wynik do dwóch miejsc po przecinku.}
{\( 23{,}43\,\mathrm{cm} \)}{\( 287{,}00\,\mathrm{cm} \)}{\( 1{,}59\,\mathrm{cm} \)}{\( 4217{,}40\,\mathrm{cm} \)}{}{}}
\LANGcs{
\Question{
Délka kruhového oblouku se středovým úhlem \( 3{,}5 \) radiánu je \( 82\,\mathrm{cm} \). Vypočítejte poloměr kruhu, ze kterého oblouk pochází. Výsledek zaokrouhlete na dvě desetinná místa.}
{\( 23{,}43\,\mathrm{cm} \)}{\( 287{,}00\,\mathrm{cm} \)}{\( 1{,}59\,\mathrm{cm} \)}{\( 4217{,}40\,\mathrm{cm} \)}{}{}}
\LANGsk{
\Question{
Dĺžka kruhového oblúka so stredovým uhlom  \( 3{,}5 \) radiánu je  \( 82\,\mathrm{cm} \). Vypočítajte polomer kruhu, z ktorého oblúk pochádza. Výsledok zaokrúhlite na dve desatinné miesta.}
{\( 23{,}43\,\mathrm{cm} \)}{\( 287{,}00\,\mathrm{cm} \)}{\( 1{,}59\,\mathrm{cm} \)}{\( 4217{,}40\,\mathrm{cm} \)}{}{}}
\LANGes{
\Question{
La longitud de un arco circular con el ańgulo central de \( 3.5 \) radianes es \( 82\,\mathrm{cm} \). Calcula el radio del círculo correspondiente. Redondea el resultado a dos decimales.}
{\( 23.43\,\mathrm{cm} \)}{\( 287.00\,\mathrm{cm} \)}{\( 1.59\,\mathrm{cm} \)}{\( 4217.40\,\mathrm{cm} \)}{}{}}
\End

\Data\ID{1003077113}
\Updated{2020-07-15 03:07:21}
\Level{B}
\SubArea{Circles}
\LANGen{
\Question{
The curved surface of a cone has an area of \( 4.15\,\mathrm{cm}^2 \). If we flatten it into a plain, we get a sector of a circle with the central angle of \( 126^{\circ} \). Calculate the volume of this cone. Round the result to two decimal places.}
{\( 0.88\,\mathrm{cm}^3 \)}{\( 0.62\,\mathrm{cm}^3 \)}{\( 0.15\,\mathrm{cm}^3 \)}{\( 311.00\,\mathrm{cm}^3 \)}{}{}}
\LANGpl{
\Question{
Powierzchnia boczna stożka wynosi \( 4{,}15\,\mathrm{cm}^2 \). Jeśli ją spłaszczymy, otrzymamy wycinek koła, którego kąt środkowy jest \( 126^{\circ} \). Oblicz objętość tego stożka. Zaokrąglij wynik do dwóch miejsc po przecinku.}
{\( 0{,}88\,\mathrm{cm}^3 \)}{\( 0{,}62\,\mathrm{cm}^3 \)}{\( 0{,}15\,\mathrm{cm}^3 \)}{\( 311{,}00\,\mathrm{cm}^3 \)}{}{}}
\LANGcs{
\Question{
Plášť kužele rozvinutý do roviny má tvar kruhového výseku se středovým úhlem \( 126^{\circ} \)  a obsahem \( 4{,}15\,\mathrm{cm}^2 \). Vypočítejte objem tohoto kuželu. Výsledek zaokrouhlete na dvě desetinná místa.}
{\( 0{,}88\,\mathrm{cm}^3 \)}{\( 0{,}62\,\mathrm{cm}^3 \)}{\( 0{,}15\,\mathrm{cm}^3 \)}{\( 311{,}00\,\mathrm{cm}^3 \)}{}{}}
\LANGsk{
\Question{
Plášť kužeľa rozvinutý do roviny má tvar kruhového výseku so stredovým uhlom \( 126^{\circ} \)  a obsahom \( 4{,}15\,\mathrm{cm}^2 \). Vypočítajte objem tohto kužeľa. Výsledok zaokrúhlite na dve desatinné  miesta.}
{\( 0{,}88\,\mathrm{cm}^3 \)}{\( 0{,}62\,\mathrm{cm}^3 \)}{\( 0{,}15\,\mathrm{cm}^3 \)}{\( 311{,}00\,\mathrm{cm}^3 \)}{}{}}
\LANGes{
\Question{
La superficie encorvanda de un cono tiene el área de \( 4.15\,\mathrm{cm}^2 \). Si lo aplanamos, obtenemos un sector circular con el ángulo central de \( 126^{\circ} \). Calcula el volumen del cono. Redondea el resultado a dos decimales.}
{\( 0.88\,\mathrm{cm}^3 \)}{\( 0.62\,\mathrm{cm}^3 \)}{\( 0.15\,\mathrm{cm}^3 \)}{\( 311.00\,\mathrm{cm}^3 \)}{}{}}
\End

\Data\ID{1003077207}
\Updated{2022-01-16 08:01:24}
\Level{B}
\SubArea{Circles}
\LANGen{
\Question{
Given a circle of which the area and the circumference have the same numerical value, determine the diameter of the circle in \( \mathrm{cm} \).}
{\( 4 \)}{\( 2 \)}{\( \sqrt2 \)}{\( \pi \)}{}{}}
\LANGpl{
\Question{
Biorąc pod uwagę okrąg, którego powierzchnia i obwód mają tę samą wartość liczbową, oblicz średnicę okręgu w \( \mathrm{cm} \).}
{\( 4 \)}{\( 2 \)}{\( \sqrt2 \)}{\( \pi \)}{}{}}
\LANGcs{
\Question{
Je daný kruh, jehož obsah i obvod mají stejnou číselnou hodnotu. Určete v \( \mathrm{cm} \) průměr tohoto kruhu.}
{\( 4 \)}{\( 2 \)}{\( \sqrt2 \)}{\( \pi \)}{}{}}
\LANGsk{
\Question{
Je daný kruh, ktorého obsah aj obvod má rovnakú číselnú hodnotu. Určte   v \( \mathrm{cm} \) priemer tohto kruhu.}
{\( 4 \)}{\( 2 \)}{\( \sqrt2 \)}{\( \pi \)}{}{}}
\LANGes{
\Question{
Dado un círculo cuya área tiene el mismo valor numérico que su perímetro. Halla el diámetro del círculo y exprésalo en \( \mathrm{cm} \).}
{\( 4 \)}{\( 2 \)}{\( \sqrt2 \)}{\( \pi \)}{}{}}
\End

\Data\ID{1003077208}
\Updated{2020-07-15 03:07:21}
\Level{B}
\SubArea{Circles}
\LANGen{
\Question{
Find the length of the diagonal of a square that has the same area as the circle with a radius of \( 10\,\mathrm{cm} \).}
{\( 10\sqrt{2\pi}\,\mathrm{cm} \)}{\( 10\sqrt{\pi}\,\mathrm{cm} \)}{\( \sqrt{20\pi}\,\mathrm{cm} \)}{\( 10\,\mathrm{cm} \)}{}{}}
\LANGpl{
\Question{
Oblicz długość przekątnej kwadratu, którego pole powierzchni jest równe polu powierzchni okręgu o promieniu \( 10\,\mathrm{cm} \).}
{\( 10\sqrt{2\pi}\,\mathrm{cm} \)}{\( 10\sqrt{\pi}\,\mathrm{cm} \)}{\( \sqrt{20\pi}\,\mathrm{cm} \)}{\( 10\,\mathrm{cm} \)}{}{}}
\LANGcs{
\Question{
Vypočítejte délku úhlopříčky čtverce, který má stejný obsah jako kruh s poloměrem \( 10\,\mathrm{cm} \).}
{\( 10\sqrt{2\pi}\,\mathrm{cm} \)}{\( 10\sqrt{\pi}\,\mathrm{cm} \)}{\( \sqrt{20\pi}\,\mathrm{cm} \)}{\( 10\,\mathrm{cm} \)}{}{}}
\LANGsk{
\Question{
Vypočítajte dĺžku uhlopriečky štvorca, ktorý má rovnaký obsah, ako kruh s polomerom \( 10\,\mathrm{cm} \).}
{\( 10\sqrt{2\pi}\,\mathrm{cm} \)}{\( 10\sqrt{\pi}\,\mathrm{cm} \)}{\( \sqrt{20\pi}\,\mathrm{cm} \)}{\( 10\,\mathrm{cm} \)}{}{}}
\LANGes{
\Question{
Calcula la longitud de la diagonal de un cuadrado cuya área equivale al área de un círculo con el radio de \( 10\,\mathrm{cm} \).}
{\( 10\sqrt{2\pi}\,\mathrm{cm} \)}{\( 10\sqrt{\pi}\,\mathrm{cm} \)}{\( \sqrt{20\pi}\,\mathrm{cm} \)}{\( 10\,\mathrm{cm} \)}{}{}}
\End

\Data\ID{1103021511}
\Updated{2022-01-16 08:01:33}
\Level{B}
\SubArea{Circles}
\IMG{1103021511.pdf}
\LANGen{
\Question{
An acute triangle \( ABC \) is inscribed in the circle of radius \( r=4\,\mathrm{cm} \). Determine the measure of the angle \( ACB \), if the length of side \( c \) is \( 6\,\mathrm{cm} \). Round the result to two decimal places. (See the picture.)}
{\( 48.59^{\circ} \)}{\( 97.18^{\circ} \)}{\( 24.30^{\circ} \)}{\( 41.41^{\circ} \)}{}{}}
\LANGpl{
\Question{
Trójkąt ostrokątny\( ABC \) jest wpisany w okrąg o promieniu \( r=4\,\mathrm{cm} \). Oblicz miary kąta \( ACB \), jeśli długość boku \( c \) to \( 6\,\mathrm{cm} \). Zaokrągli wynik do dwóch miejsc po przecinku. (Spójrz na rysunek.)}
{\( 48{,}59^{\circ} \)}{\( 97{,}18^{\circ} \)}{\( 24{,}30^{\circ} \)}{\( 41{,}41^{\circ} \)}{}{}}
\LANGcs{
\Question{
Ostroúhlý trojúhelník \( ABC \) je vepsaný do kružnice s poloměrem \( r=4\,\mathrm{cm} \). Jakou velikost má úhel  \( ACB \), pokud je délka strany \( c \)  \( 6\,\mathrm{cm} \). Výsledek zaokrouhlete na dvě desetinná místa. (viz obrázek)}
{\( 48{,}59^{\circ} \)}{\( 97{,}18^{\circ} \)}{\( 24{,}30^{\circ} \)}{\( 41{,}41^{\circ} \)}{}{}}
\LANGsk{
\Question{
Ostrouhlý trojuholník  \( ABC \) je vpísaný do kružnice s polomerom \( r=4\,\mathrm{cm} \). Akú veľkosť má uhol  \( ACB \), ak dĺžka strany  \( c \) je \( 6\,\mathrm{cm} \). Výsledok zaokrúhlite na dve desatinné miesta. (Pozrite obrázok.)}
{\( 48{,}59^{\circ} \)}{\( 97{,}18^{\circ} \)}{\( 24{,}30^{\circ} \)}{\( 41{,}41^{\circ} \)}{}{}}
\LANGes{
\Question{
Un triángulo acutángulo \( ABC \) se inscribe en una circunferencia con radio \( r=4\,\mathrm{cm} \). Calcula la medida del ángulo \( ACB \), suponiendo que el lado \( c \) mide \( 6\,\mathrm{cm} \). Redondea el resultado a dos decimales.}
{\( 48.59^{\circ} \)}{\( 97.18^{\circ} \)}{\( 24.30^{\circ} \)}{\( 41.41^{\circ} \)}{}{}}
\End

\Data\ID{1103021513}
\Updated{2020-07-15 03:07:21}
\Level{B}
\SubArea{Circles}
\IMG{1103021513.pdf}
\LANGen{
\Question{
The distance of the chord \( AB \) from the centre of the circle is equal to \( 2/3 \) of its radius. Find the measure of the angle \(  SAB \). (See the picture.) Round the result to two decimal places.}
{\( 41.81^{\circ} \)}{\( 48.19^{\circ} \)}{\( 33.69^{\circ} \)}{\( 56.31^{\circ} \)}{}{}}
\LANGpl{
\Question{
Odległość cięciwy \( AB \) od środka okręgu jest równa \( 2/3 \) jego promienia. Oblicz miarę kąta \(  SAB \). (Spójrz na rysunek.) Zaokrągli wynik do dwóch miejsc po przecinku.}
{\( 41{,}81^{\circ} \)}{\( 48{,}19^{\circ} \)}{\( 33{,}69^{\circ} \)}{\( 56{,}31^{\circ} \)}{}{}}
\LANGcs{
\Question{
Vzdálenost tětivy \( AB \) od středu kružnice se rovná \( 2/3 \) poloměru dané kružnice. Vypočítejte velikost úhlu \(  SAB \) (viz obrázek). Výsledek zaokrouhlete na dvě desetinná místa.}
{\( 41{,}81^{\circ} \)}{\( 48{,}19^{\circ} \)}{\( 33{,}69^{\circ} \)}{\( 56{,}31^{\circ} \)}{}{}}
\LANGsk{
\Question{
Vzdialenosť tetivy  \( AB \) od stredu kružnice sa rovná  \( 2/3 \) polomeru danej kružnice.  Vypočítajte veľkosť uhla  \(  SAB \). (Pozrite obrázok.) Výsledok zaokrúhlite na dve desatinné miesta.}
{\( 41{,}81^{\circ} \)}{\( 48{,}19^{\circ} \)}{\( 33{,}69^{\circ} \)}{\( 56{,}31^{\circ} \)}{}{}}
\LANGes{
\Question{
La distancia de la cuerda \( AB \) desde el centro de la circunferencia equivale a \( 2/3 \) de su radio. Calcula la medida del ángulo \(  SAB \). (Mira la imagen). Redondea el resultado a dos decimales.}
{\( 41.81^{\circ} \)}{\( 48.19^{\circ} \)}{\( 33.69^{\circ} \)}{\( 56.31^{\circ} \)}{}{}}
\End

\Data\ID{1103021601}
\Updated{2020-07-15 03:07:21}
\Level{B}
\SubArea{Circles}
\IMG{1103021601.pdf}
\LANGen{
\Question{
The distance from the point \( V \) to the centre \( S \) of the circle \( k \) is \( 30\,\mathrm{cm} \). The radius of the circle is \( 15\,\mathrm{cm} \). From the point \( V \) two tangent lines to the circle \( k \) can be drawn. What is the measure of the angle between them? (See the picture.)}
{\( 60^{\circ} \)}{\( 30^{\circ} \)}{\( 90^{\circ} \)}{\( 45^{\circ} \)}{}{}}
\LANGpl{
\Question{
Odległość od punktu \( V \) do środka \( S \) okręgu \( k \) wynosi \( 30\,\mathrm{cm} \). Promień okręgu jest równy  \( 15\,\mathrm{cm} \). Od punktu \( V \) dwie styczne okręgu \( k \) mogą być narysowane. Jaka jest miara kąta pomiędzy nimi? (Spójrz na rysunek.)}
{\( 60^{\circ} \)}{\( 30^{\circ} \)}{\( 90^{\circ} \)}{\( 45^{\circ} \)}{}{}}
\LANGcs{
\Question{
Vzdálenost bodu \( V \) od středu \( S \) kružnice  \( k \) je \( 30\,\mathrm{cm} \). Poloměr kružnice je \( 15\,\mathrm{cm} \).  Bodem \( V \) vedou dvě tečny ke kružnici \( k \). Jakou velikost má úhel, který svírají tyto tečny? (viz obrázek)}
{\( 60^{\circ} \)}{\( 30^{\circ} \)}{\( 90^{\circ} \)}{\( 45^{\circ} \)}{}{}}
\LANGsk{
\Question{
Vzdialenosť bodu \( V \) od stredu \( S \) kružnice  \( k \) je \( 30\,\mathrm{cm} \). Polomer kružnice je  \( 15\,\mathrm{cm} \).  Bodom \( V \) vedú dve dotyčnice  ku kružnici \( k \). Akú veľkosť má uhol, ktorý zvierajú tieto dotyčnice? (Pozri obrázok.)}
{\( 60^{\circ} \)}{\( 30^{\circ} \)}{\( 90^{\circ} \)}{\( 45^{\circ} \)}{}{}}
\LANGes{
\Question{
La distancia desde el punto \( V \) hacia el centro \( S \) de la circunferencia \( k \) es \( 30\,\mathrm{cm} \). El radio de la circunferencia mide \( 15\,\mathrm{cm} \). Desde el punto \( V \) se pueden dibujar dos tangentes hacia la circunferencia \( k \). ¿Cuál es la medida del ángulo formado por las dos tangentes?}
{\( 60^{\circ} \)}{\( 30^{\circ} \)}{\( 90^{\circ} \)}{\( 45^{\circ} \)}{}{}}
\End

\Data\ID{1103021602}
\Updated{2020-07-15 03:07:21}
\Level{B}
\SubArea{Circles}
\IMG{1103021602.pdf}
\LANGen{
\Question{
The side of an equilateral triangle is \( 6\,\mathrm{cm} \) long. Find the area of the annulus between the incircle and circumcircle of the given triangle. (See the picture.)}
{\( 9\pi\,\mathrm{cm}^2 \)}{\( 6\pi\,\mathrm{cm}^2 \)}{\( 12\pi\,\mathrm{cm}^2 \)}{\( 8\pi\,\mathrm{cm}^2 \)}{}{}}
\LANGpl{
\Question{
Bok trójkąta równobocznego jest równy \( 6\,\mathrm{cm} \). Oblicz pole powierzchni pierścienia pomiędzy okręgiem wpisany i opisanym na trójkącie. (Spójrz na rysunek.)}
{\( 9\pi\,\mathrm{cm}^2 \)}{\( 6\pi\,\mathrm{cm}^2 \)}{\( 12\pi\,\mathrm{cm}^2 \)}{\( 8\pi\,\mathrm{cm}^2 \)}{}{}}
\LANGcs{
\Question{
Strana rovnostranného trojúhelníku je dlouhá \( 6\,\mathrm{cm} \). Určete obsah mezikruží ohraničeného vepsanou a opsanou kružnicí daného trojúhelníku (viz obrázek).}
{\( 9\pi\,\mathrm{cm}^2 \)}{\( 6\pi\,\mathrm{cm}^2 \)}{\( 12\pi\,\mathrm{cm}^2 \)}{\( 8\pi\,\mathrm{cm}^2 \)}{}{}}
\LANGsk{
\Question{
Strana rovnostranného trojuholníka je dlhá \( 6\,\mathrm{cm} \). Určte obsah medzikružia ohraničeného vpísanou a opísanou kružnicou daného trojuholníka.  (Pozri obrázok.)}
{\( 9\pi\,\mathrm{cm}^2 \)}{\( 6\pi\,\mathrm{cm}^2 \)}{\( 12\pi\,\mathrm{cm}^2 \)}{\( 8\pi\,\mathrm{cm}^2 \)}{}{}}
\LANGes{
\Question{
El lado de un triángulo equilátero mide \( 6\,\mathrm{cm} \). Calcula el área de la corona circular entre las circunferencias inscrita y circunscrita del triángulo dado. (Mira la imagen).}
{\( 9\pi\,\mathrm{cm}^2 \)}{\( 6\pi\,\mathrm{cm}^2 \)}{\( 12\pi\,\mathrm{cm}^2 \)}{\( 8\pi\,\mathrm{cm}^2 \)}{}{}}
\End

\Data\ID{1103021604}
\Updated{2020-07-15 03:07:21}
\Level{B}
\SubArea{Circles}
\IMG{1103021604.pdf}
\LANGen{
\Question{
Calculate the radius of a circle inscribed into the rhombus \( ABCD \) if the length of its side is \( 10\,\mathrm{cm} \) and the measure of the angle \( DAB \) is \( 40^{\circ} \). (See the picture.) Round the result to two decimal places.}
{\( 3.21\,\mathrm{cm} \)}{\( 1.71\,\mathrm{cm} \)}{\( 3.83\,\mathrm{cm} \)}{\( 6.42\,\mathrm{cm} \)}{}{}}
\LANGpl{
\Question{
Oblicz promień okręgu wpisanego w romb \( ABCD \), długość jego boku wynosi \( 10\,\mathrm{cm} \), a miara kąta  \( DAB \) to \( 40^{\circ} \). (Spójrz na rysunek.) Zaokrągli wynik do dwóch miejsc po przecinku.}
{\( 3{,}21\,\mathrm{cm} \)}{\( 1{,}71\,\mathrm{cm} \)}{\( 3{,}83\,\mathrm{cm} \)}{\( 6{,}42\,\mathrm{cm} \)}{}{}}
\LANGcs{
\Question{
Vypočítejte poloměr kružnice vepsané do kosočtverce \( ABCD \). Délka strany kosočtverce je \( 10\,\mathrm{cm} \)  a velikost úhlu \( DAB \) je \( 40^{\circ} \) (viz obrázek). Výsledek zaokrouhlete na dvě desetinná místa.}
{\( 3{,}21\,\mathrm{cm} \)}{\( 1{,}71\,\mathrm{cm} \)}{\( 3{,}83\,\mathrm{cm} \)}{\( 6{,}42\,\mathrm{cm} \)}{}{}}
\LANGsk{
\Question{
Vypočítajte polomer vpísanej kružnice do kosoštvorca s \( ABCD \). Dĺžka strany kosoštvorca je \( 10\,\mathrm{cm} \)  a veľkosť uhla \( DAB \) is \( 40^{\circ} \). (Pozri obrázok.) Výsledok zaokrúhlite na dve desatinné miesta.}
{\( 3{,}21\,\mathrm{cm} \)}{\( 1{,}71\,\mathrm{cm} \)}{\( 3{,}83\,\mathrm{cm} \)}{\( 6{,}42\,\mathrm{cm} \)}{}{}}
\LANGes{
\Question{
Calcula el radio de una circunferencia inscrita en el rombo \( ABCD \) cuyo lado mide \( 10\,\mathrm{cm} \) y la medida del ángulo \( DAB \) es \( 40^{\circ} \). (Mira la imagen). Redondea el resultado a dos decimales.}
{\( 3.21\,\mathrm{cm} \)}{\( 1.71\,\mathrm{cm} \)}{\( 3.83\,\mathrm{cm} \)}{\( 6.42\,\mathrm{cm} \)}{}{}}
\End

\Data\ID{1103021605}
\Updated{2022-01-16 09:01:17}
\Level{B}
\SubArea{Circles}
\IMG{1103021605.pdf}
\LANGen{
\Question{
A circle of radius \( 22\,\mathrm{cm} \) is inscribed into the rhombus \( ABCD \). Calculate the measure of the angle \( CAB \) if the length of the rhombus side is \( 90\,\mathrm{cm} \). (See the picture.) Round the result to two decimal places.}
{\( 14.63^{\circ} \)}{\( 29.27^{\circ} \)}{\( 30.37^{\circ} \)}{\( 28.30^{\circ} \)}{}{}}
\LANGpl{
\Question{
Okrąg o promieniu równym \( 22\,\mathrm{cm} \) jest wpisany w romb \( ABCD \). Oblicz miarę kąta \( CAB \) jeśli romb ma bok o długości  \( 90\,\mathrm{cm} \).  (Spójrz na rysunek.) Zaokrągli wynik do dwóch miejsc po przecinku.}
{\( 14{,}63^{\circ} \)}{\( 29{,}27^{\circ} \)}{\( 30{,}37^{\circ} \)}{\( 28{,}30^{\circ} \)}{}{}}
\LANGcs{
\Question{
Do kosočtverce \( ABCD \) je vepsaná kružnice s poloměrem \( 22\,\mathrm{cm} \). Vypočítejte velikost úhlu \( CAB \), jestliže je délka strany kosočtverce \( 90\,\mathrm{cm} \) (viz obrázek). Výsledek zaokrouhlete na dvě desetinná místa.}
{\( 14{,}63^{\circ} \)}{\( 29{,}27^{\circ} \)}{\( 30{,}37^{\circ} \)}{\( 28{,}30^{\circ} \)}{}{}}
\LANGsk{
\Question{
Do kosoštvorca \( ABCD \) je vpísaná kružnica s polomerom \( 22\,\mathrm{cm} \). Vypočítajte veľkosť uhla \( CAB \), ak dĺžka strany kosoštvorca je \( 90\,\mathrm{cm} \). (Pozri obrázok.) Výsledok zaokrúhlite na dve desatinné miesta.}
{\( 14{,}63^{\circ} \)}{\( 29{,}27^{\circ} \)}{\( 30{,}37^{\circ} \)}{\( 28{,}30^{\circ} \)}{}{}}
\LANGes{
\Question{
Una circunferencia con un radio de \( 22\,\mathrm{cm} \) se inscribe en el rombo \( ABCD \). Calcula la medida del ángulo \( CAB \), suponiendo que el lado del rombo mide \( 90\,\mathrm{cm} \). Redondea el resultado a dos decimales.}
{\( 14.63^{\circ} \)}{\( 29.27^{\circ} \)}{\( 30.37^{\circ} \)}{\( 28.30^{\circ} \)}{}{}}
\End

\Data\ID{1103021606}
\Updated{2020-07-15 03:07:21}
\Level{B}
\SubArea{Circles}
\IMG{1103021606.pdf}
\LANGen{
\Question{
In the rectangle \( ABCD \), \( a=6\,\mathrm{cm} \) and the radius of the circumcircle \( r=4\,\mathrm{cm} \) (see the picture). Find the measure of the angle between the diagonals of the rectangle. Round the result to two decimal places.}
{\( 82.82^{\circ} \)}{\( 48.59^{\circ} \)}{\( 97.18^{\circ} \)}{\( 36.12^{\circ} \)}{}{}}
\LANGpl{
\Question{
Dany jest prostokąt \( ABCD \), \( a=6\,\mathrm{cm} \) a promień okręgu opisanego jest równy \( r=4\,\mathrm{cm} \) (spójrz na rysunek). Oblicz miarę kąta pomiędzy przekątnymi prostokąta. Zaokrągli wynik do dwóch miejsc po przecinku.}
{\( 82{,}82^{\circ} \)}{\( 48{,}59^{\circ} \)}{\( 97{,}18^{\circ} \)}{\( 36{,}12^{\circ} \)}{}{}}
\LANGcs{
\Question{
Je dán obdélník \( ABCD \), jehož strana \( a=6\,\mathrm{cm} \). Obdélníku je opsaná kružnice s poloměrem \( r=4\,\mathrm{cm} \) (viz obrázek). Vypočítejte velikost úhlu, který svírají úhlopříčky obdélníka. Výsledek zaokrouhlete na dvě desetinná místa.}
{\( 82{,}82^{\circ} \)}{\( 48{,}59^{\circ} \)}{\( 97{,}18^{\circ} \)}{\( 36{,}12^{\circ} \)}{}{}}
\LANGsk{
\Question{
Daný je obdĺžnik \( ABCD \), ktorého strana \( a=6\,\mathrm{cm} \). Obdĺžniku je opísaná kružnica s polomerom \( r=4\,\mathrm{cm} \) (pozri obrázok). Vypočítajte veľkosť uhla, ktorý zvierajú uhlopriečky obdĺžnika. Výsledok zaokrúhlite na dve desatinné miesta.}
{\( 82{,}82^{\circ} \)}{\( 48{,}59^{\circ} \)}{\( 97{,}18^{\circ} \)}{\( 36{,}12^{\circ} \)}{}{}}
\LANGes{
\Question{
Dado el rectángulo \( ABCD \), con el lado \( a=6\,\mathrm{cm} \) y el radio de la circunferencia circunscrita \( r=4\,\mathrm{cm} \) (mira la imagen). Calcula la medida del ángulo formado por las diagonales del rectángulo. Redondea el resultado a dos decimales.}
{\( 82.82^{\circ} \)}{\( 48.59^{\circ} \)}{\( 97.18^{\circ} \)}{\( 36.12^{\circ} \)}{}{}}
\End

\Data\ID{1103021608}
\Updated{2022-04-23 05:04:54}
\Level{B}
\SubArea{Circles}
\IMG{1103021608.pdf}
\LANGen{
\Question{
Consider a circle \( k \) with radius \( 2.5\,\mathrm{cm} \). In the circle is inscribed a convex quadrilateral \( ABCD \) so that the diagonal \( AC \) is the diameter of the circle, the length of \( BC \) is \( \sqrt{21}\,\mathrm{cm} \), and the length of \( DC \) is \( 4\,\mathrm{cm} \). What is the length of the shortest side of this quadrilateral? (See the picture.)}
{\( 2\,\mathrm{cm} \)}{\( 3\,\mathrm{cm} \)}{\( \sqrt5\,\mathrm{cm} \)}{\( 2.5\,\mathrm{cm} \)}{}{}}
\LANGpl{
\Question{
Dany jest okrąg  \( k \) o promieniu  \( 2{,}5\,\mathrm{cm} \). Czworobok \( ABCD \) jest wpisany w okrąg tak, aby przekątna \( AC \) była  średnicą okręgu, długość  \( BC \) to \( \sqrt{21}\,\mathrm{cm} \), długość \( DC \) to \( 4\,\mathrm{cm} \). Jaka jest długość najkrótszego boku danego czworoboku? (Spójrz na rysunek.)}
{\( 2\,\mathrm{cm} \)}{\( 3\,\mathrm{cm} \)}{\( \sqrt5\,\mathrm{cm} \)}{\( 2{,}5\,\mathrm{cm} \)}{}{}}
\LANGcs{
\Question{
Je dána kružnice \( k \) s poloměrem \( 2{,}5\,\mathrm{cm} \). Do této kružnice je vepsaný konvexní čtyřúhelník \( ABCD \). Úhlopříčka \( AC \) je průměrem kružnice, \( BC =\)  \( \sqrt{21}\,\mathrm{cm} \) a \( DC = \) \( 4\,\mathrm{cm} \). Jakou délku má nejkratší strana tohoto čtyřúhelníka? (Viz obrázek.)}
{\( 2\,\mathrm{cm} \)}{\( 3\,\mathrm{cm} \)}{\( \sqrt5\,\mathrm{cm} \)}{\( 2{,}5\,\mathrm{cm} \)}{}{}}
\LANGsk{
\Question{
Daná je kružnica \( k \) s polomerom \( 2{,}5\,\mathrm{cm} \). Do tejto kružnice je vpísaný konvexný štvoruholník \( ABCD \). Uhlopriečka \( AC \) je priemerom kružnice, \( BC =\)  \( \sqrt{21}\,\mathrm{cm} \)  a \( DC = \) \( 4\,\mathrm{cm} \).  Akú dĺžku má najkratšia strana tohto štvoruholníka?  (Pozri obrázok.)}
{\( 2\,\mathrm{cm} \)}{\( 3\,\mathrm{cm} \)}{\( \sqrt5\,\mathrm{cm} \)}{\( 2{,}5\,\mathrm{cm} \)}{}{}}
\LANGes{
\Question{
Dada la circunferencia \( k \) cuyo radio mide \( 2.5\,\mathrm{cm} \). En la circunferencia se inscribe el cuadrilátero convexo \( ABCD \) cuya diagonal \( AC \) es el diámetro de la circunferencia. La longitud del \( BC \) es \( \sqrt{21}\,\mathrm{cm} \), y la longitud del \( DC \) es \( 4\,\mathrm{cm} \). ¿Cuánto mide el lado más corto del cuadrilátero?}
{\( 2\,\mathrm{cm} \)}{\( 3\,\mathrm{cm} \)}{\( \sqrt5\,\mathrm{cm} \)}{\( 2.5\,\mathrm{cm} \)}{}{}}
\End

\Data\ID{1103021609}
\Updated{2022-01-16 09:01:20}
\Level{B}
\SubArea{Circles}
\IMG{1103021609.pdf}
\LANGen{
\Question{
Points \( A \), \( B \) and \( C \) lie on the circle \( k \). The line segment \( AC \) is the diameter of the circle  and the lines \( AC \) and \( BC \) contain the angle of \( 60^{\circ} \). Calculate the length of \( AC \) if the length of \( BC \) is \( 10\,\mathrm{cm} \). (See the picture.)}
{\( 20\,\mathrm{cm} \)}{\( 5\sqrt3\,\mathrm{cm} \)}{\( 5\,\mathrm{cm} \)}{\( 2\sqrt3\,\mathrm{cm} \)}{}{}}
\LANGpl{
\Question{
Punkty \( A \), \( B \) i \( C \) leżą na okręgu \( k \). Odcinek \( AC \) to średnica okręgu, odcinki \( AC \) i \( BC \) zawierają kąt  \( 60^{\circ} \). Oblicz długość\( AC \) jeśli długość \( BC \) to \( 10\,\mathrm{cm} \). (Spójrz na rysunek.)}
{\( 20\,\mathrm{cm} \)}{\( 5\sqrt3\,\mathrm{cm} \)}{\( 5\,\mathrm{cm} \)}{\( 2\sqrt3\,\mathrm{cm} \)}{}{}}
\LANGcs{
\Question{
Na kružnici \( k \) jsou dány body \( A \), \( B \) a \( C \). Úsečka \( AC \) je průměrem kružnice \( k \) a úsečky \( AC \) a \( BC \) svírají úhel \( 60^{\circ} \). Vypočítejte délku úsečky \( AC \), jestliže délka úsečky \( BC \) je  \( 10\,\mathrm{cm} \).}
{\( 20\,\mathrm{cm} \)}{\( 5\sqrt3\,\mathrm{cm} \)}{\( 5\,\mathrm{cm} \)}{\( 2\sqrt3\,\mathrm{cm} \)}{}{}}
\LANGsk{
\Question{
Na kružnici \( k \) sú dané body \( A \), \( B \) a \( C \). Úsečka \( AC \) je priemerom kružnice \( k \) a úsečky \( AC \) a \( BC \) zvierajú uhol \( 60^{\circ} \). Vypočítajte dĺžku úsečky  \( AC \), ak dĺžka úsečky \( BC \) je  \( 10\,\mathrm{cm} \).}
{\( 20\,\mathrm{cm} \)}{\( 5\sqrt3\,\mathrm{cm} \)}{\( 5\,\mathrm{cm} \)}{\( 2\sqrt3\,\mathrm{cm} \)}{}{}}
\LANGes{
\Question{
Los puntos \( A \), \( B \) y \( C \) pertenecen a la circunferencia \( k \). El segmento \( AC \) es el diámetro de la circunferencia y los segmentos \( AC \) y \( BC \) forman un ángulo de \( 60^{\circ} \). Calcula la longitud del \( AC \) si la longitud del \( BC \) es \( 10\,\mathrm{cm} \).}
{\( 20\,\mathrm{cm} \)}{\( 5\sqrt3\,\mathrm{cm} \)}{\( 5\,\mathrm{cm} \)}{\( 2\sqrt3\,\mathrm{cm} \)}{}{}}
\End

\Data\ID{1103021610}
\Updated{2020-07-15 03:07:21}
\Level{B}
\SubArea{Circles}
\IMG{1103021610.pdf}
\LANGen{
\Question{
A regular hexagon \( ABCDEF \) is inscribed in a circle with radius of \( 12\,\mathrm{cm} \). Calculate the length of its diagonal \( EC \). (See the picture.)}
{\( 12\sqrt3\,\mathrm{cm} \)}{\( 12\,\mathrm{cm} \)}{\( 6\sqrt3\,\mathrm{cm} \)}{\( 24\,\mathrm{cm} \)}{}{}}
\LANGpl{
\Question{
Sześciokąt foremny \( ABCDEF \) został wpisany w okrąg o promieniu \( 12\,\mathrm{cm} \). Oszacuj długość jego przekątnej  \( EC \). (Patrz rysunek.)}
{\( 12\sqrt3\,\mathrm{cm} \)}{\( 12\,\mathrm{cm} \)}{\( 6\sqrt3\,\mathrm{cm} \)}{\( 24\,\mathrm{cm} \)}{}{}}
\LANGcs{
\Question{
Pravidelný šestiúhelník \( ABCDEF \) je vepsaný do kružnice s poloměrem \( 12\,\mathrm{cm} \). Vypočítejte délku jeho úhlopříčky \( EC \) (viz obrázek).}
{\( 12\sqrt3\,\mathrm{cm} \)}{\( 12\,\mathrm{cm} \)}{\( 6\sqrt3\,\mathrm{cm} \)}{\( 24\,\mathrm{cm} \)}{}{}}
\LANGsk{
\Question{
Pravidelný šesťuholník \( ABCDEF \) je vpísaný do kružnice s polomerom \( 12\,\mathrm{cm} \). Vypočítajte dĺžku jeho uhlopriečky \( EC \). (Pozri obrázok.)}
{\( 12\sqrt3\,\mathrm{cm} \)}{\( 12\,\mathrm{cm} \)}{\( 6\sqrt3\,\mathrm{cm} \)}{\( 24\,\mathrm{cm} \)}{}{}}
\LANGes{
\Question{
Un hexágono regular \( ABCDEF \) se inscribe en una circunferencia cuyo radio mide \( 12\,\mathrm{cm} \). Calcula la longitud de su diagonal \( EC \).}
{\( 12\sqrt3\,\mathrm{cm} \)}{\( 12\,\mathrm{cm} \)}{\( 6\sqrt3\,\mathrm{cm} \)}{\( 24\,\mathrm{cm} \)}{}{}}
\End

\Data\ID{1103021611}
\Updated{2022-01-16 09:01:59}
\Level{B}
\SubArea{Circles}
\IMG{1103021611.pdf}
\LANGen{
\Question{
What is the length of the side of a regular pentagon circumscribed to a circle with radius of \( 9\,\mathrm{cm} \)?  (See the picture.) Round the result to two decimal places.}
{\( 13.08\,\mathrm{cm} \)}{\( 55.39\,\mathrm{cm} \)}{\( 6.54\,\mathrm{cm} \)}{\( 10.58\,\mathrm{cm} \)}{}{}}
\LANGpl{
\Question{
Jaka jest długość boku pięciokąta foremnego opisanego na kole o promieniu \( 9\,\mathrm{cm} \)?  (Patrz rysunek.) Zaokrąglij do dwóch miejsc dziesiętnych.}
{\( 13{,}08\,\mathrm{cm} \)}{\( 55{,}39\,\mathrm{cm} \)}{\( 6{,}54\,\mathrm{cm} \)}{\( 10{,}58\,\mathrm{cm} \)}{}{}}
\LANGcs{
\Question{
Jaká je délka strany pravidelného pětiúhelníku, do kterého je vepsaná kružnice s poloměrem \( 9\,\mathrm{cm} \) (viz obrázek)? Výsledek zaokrouhlete na dvě desetinná místa.}
{\( 13{,}08\,\mathrm{cm} \)}{\( 55{,}39\,\mathrm{cm} \)}{\( 6{,}54\,\mathrm{cm} \)}{\( 10{,}58\,\mathrm{cm} \)}{}{}}
\LANGsk{
\Question{
Aká je dĺžka strany pravidelného päťuholníka, do ktorého je vpísaná kružnica s polomerom \( 9\,\mathrm{cm} \)? (Pozri obrázok.) Výsledok zaokrúhlite na dve desatinné miesta.}
{\( 13{,}08\,\mathrm{cm} \)}{\( 55{,}39\,\mathrm{cm} \)}{\( 6{,}54\,\mathrm{cm} \)}{\( 10{,}58\,\mathrm{cm} \)}{}{}}
\LANGes{
\Question{
Dada una circunferencia, cuyo radio mide  \( 9\,\mathrm{cm} \), que se inscribe en un pentágono regular. ¿Cuánto mide su lado? Redondea el resultado a dos decimales.}
{\( 13.08\,\mathrm{cm} \)}{\( 55.39\,\mathrm{cm} \)}{\( 6.54\,\mathrm{cm} \)}{\( 10.58\,\mathrm{cm} \)}{}{}}
\End

\Data\ID{1103021612}
\Updated{2022-01-16 09:01:25}
\Level{B}
\SubArea{Circles}
\IMG{1103021612.pdf}
\LANGen{
\Question{
Consider two circles: the circle \( k \) with centre \( S_1 \) and radius \( 3\,\mathrm{cm} \), and the circle \( n \) with  centre \( S_2 \) and radius \( 8\,\mathrm{cm} \). The distance between \( S_1 \) and \( S_2 \) is \( 22\,\mathrm{cm} \). Common internal tangents of the circles intersect at point \( A \). Calculate the distance of the point \( A \) from the centre \( S_1 \). (See the picture.)}
{\( 6\,\mathrm{cm} \)}{\( 16\,\mathrm{cm} \)}{\( 11\,\mathrm{cm} \)}{\( 5\,\mathrm{cm} \)}{}{}}
\LANGpl{
\Question{
Rozważ dwa okręgi: okrąg \( k \) ze środkiem \( S_1 \) i promieniem  \( 3\,\mathrm{cm} \), oraz okrąg \( n \) ze środkiem  \( S_2 \) i promieniem \( 8\,\mathrm{cm} \). Odległość pomiędzy \( S_1 \) i \( S_2 \) wynosi \( 22\,\mathrm{cm} \). Wspólne wewnętrzne styczne przecinają się w punkcie \( A \). Oszacuj odległość punktu \( A \) od środka \( S_1 \). (Patrz rysunek.)}
{\( 6\,\mathrm{cm} \)}{\( 16\,\mathrm{cm} \)}{\( 11\,\mathrm{cm} \)}{\( 5\,\mathrm{cm} \)}{}{}}
\LANGcs{
\Question{
Jsou dány dvě kružnice: \( k \) se středem \( S_1 \) a poloměrem \( 3\,\mathrm{cm} \) a kružnice \( n \) se středem \( S_2 \) a poloměrem \( 8\,\mathrm{cm} \). Vzdálenost \( S_1 \) a \( S_2 \) je \( 22\,\mathrm{cm} \). Společné vnitřní tečny těchto kružnic se protínají v bodě \( A \). Vypočítejte vzdálenost bodu \( A \) od středu \( S_1 \) (viz obrázek).}
{\( 6\,\mathrm{cm} \)}{\( 16\,\mathrm{cm} \)}{\( 11\,\mathrm{cm} \)}{\( 5\,\mathrm{cm} \)}{}{}}
\LANGsk{
\Question{
Dané sú dve kružnice: \( k \)  so stredom \( S_1 \) a polomerom \( 3\,\mathrm{cm} \) a kružnica \( n \) so stredom \( S_2 \) a polomerom \( 8\,\mathrm{cm} \). Vzdialenosť \( S_1 \) a \( S_2 \) je \( 22\,\mathrm{cm} \). Spoločné vnútorné dotyčnice týchto kružníc sa pretínajú v bode \( A \). Vypočítajte vzdialenosť bodu \( A \) od stredu \( S_1 \). (Pozri obrázok.)}
{\( 6\,\mathrm{cm} \)}{\( 16\,\mathrm{cm} \)}{\( 11\,\mathrm{cm} \)}{\( 5\,\mathrm{cm} \)}{}{}}
\LANGes{
\Question{
Dadas dos circunferencias: la circunferencia \( k \) con centro en \( S_1 \) y  radio de \( 3\,\mathrm{cm} \), y la circunferencia \( n \) con centro en \( S_2 \) y radio de \( 8\,\mathrm{cm} \). La distancia entre \( S_1 \) y \( S_2 \) es \( 22\,\mathrm{cm} \). Las tangentes comunes se cortan en el punto \( A \). Calcula la distancia desde el punto \( A \) hasta el centro \( S_1 \).}
{\( 6\,\mathrm{cm} \)}{\( 16\,\mathrm{cm} \)}{\( 11\,\mathrm{cm} \)}{\( 5\,\mathrm{cm} \)}{}{}}
\End

\Data\ID{1103021613}
\Updated{2020-07-15 03:07:21}
\Level{B}
\SubArea{Circles}
\IMG{1103021613.pdf}
\LANGen{
\Question{
A circle is inscribed in a rhombus \( ABCD \). The touching points of the circle and the rhombus divide each side into two parts that are \( 12\,\mathrm{dm} \) and \( 25\,\mathrm{dm} \) long. (See the picture.) Find the measure of the angle \( CAB \). Round the result to two decimal places.}
{\( 34.72^{\circ} \)}{\( 43.85^{\circ} \)}{\( 46.15^{\circ} \)}{\( 23.14^{\circ} \)}{}{}}
\LANGpl{
\Question{
Okrąg jest wpisany w romb \( ABCD \). Punkty styczne okręgu i rombu dzielą każdy bok na dwie części o długości \( 12\,\mathrm{dm} \) i \( 25\,\mathrm{dm} \). (Patrz rysunek.) Wyznacz miarę kąta \( CAB \).  Zaokrąglij wynik do dwóch miejsc dziesiętnych.}
{\( 34{,}72^{\circ} \)}{\( 43{,}85^{\circ} \)}{\( 46{,}15^{\circ} \)}{\( 23{,}14^{\circ} \)}{}{}}
\LANGcs{
\Question{
Do kosočtverce \( ABCD \) je vepsaná kružnice. Body dotyku kružnice a kosočtverce rozdělují jeho strany na části dlouhé \( 12\,\mathrm{dm} \) a \( 25\,\mathrm{dm} \) (viz obrázek). Vypočítejte velikost úhlu \( CAB \). Výsledek zaokrouhlete na dvě desetinná místa.}
{\( 34{,}72^{\circ} \)}{\( 43{,}85^{\circ} \)}{\( 46{,}15^{\circ} \)}{\( 23{,}14^{\circ} \)}{}{}}
\LANGsk{
\Question{
Do kosoštvorca \( ABCD \) je vpísaná kružnica. Body dotyku kružnice a kosoštvorca rozdeľujú jeho strany na časti dlhé \( 12\,\mathrm{dm} \) a \( 25\,\mathrm{dm} \). (Pozri obrázok.) Vypočítajte veľkosť uhla \( CAB \). Výsledok zaokrúhlite na dve desatinné miesta.}
{\( 34{,}72^{\circ} \)}{\( 43{,}85^{\circ} \)}{\( 46{,}15^{\circ} \)}{\( 23{,}14^{\circ} \)}{}{}}
\LANGes{
\Question{
Una circunferencia se inscribe en un rombo \( ABCD \). Los puntos de tangencia de la circunferencia y del rombo dividen a cada lado en dos partes que miden \( 12\,\mathrm{dm} \) y \( 25\,\mathrm{dm} \). (Mira la imagen). Calcula la medida del ángulo \( CAB \). Redondea el resultado a dos decimales.}
{\( 34.72^{\circ} \)}{\( 43.85^{\circ} \)}{\( 46.15^{\circ} \)}{\( 23.14^{\circ} \)}{}{}}
\End

\Data\ID{1103077101}
\Updated{2022-04-23 05:04:00}
\Level{B}
\SubArea{Circles}
\IMG{1103077101.pdf}
\LANGen{
\Question{
Derive the formula for calculating the perimeter of an equilateral triangle inscribed in a circle of radius \( r \).}
{\( 3\sqrt3 r \)}{\( \sqrt3 r \)}{\( 3 r \)}{\( \frac{3\sqrt3 r}2 \)}{}{}}
\LANGpl{
\Question{
Wyznacz wzór, który pozwoli obliczyć obwód trójkąta równobocznego wpisanego w okrąg o promieniu \( r \).}
{\( 3\sqrt3 r \)}{\( \sqrt3 r \)}{\( 3 r \)}{\( \frac{3\sqrt3 r}2 \)}{}{}}
\LANGcs{
\Question{
Odvoďte vztah pro výpočet obvodu rovnostranného trojúhelníku vepsaného do kružnice s poloměrem \( r \).}
{\( 3\sqrt3 r \)}{\( \sqrt3 r \)}{\( 3 r \)}{\( \frac{3\sqrt3 r}2 \)}{}{}}
\LANGsk{
\Question{
Odvoďte vzťah  pre výpočet obvodu rovnostranného trojuholníka vpísaného do kružnice s polomerom  \( r \).}
{\( 3\sqrt3 r \)}{\( \sqrt3 r \)}{\( 3 r \)}{\( \frac{3\sqrt3 r}2 \)}{}{}}
\LANGes{
\Question{
Averigua la fórmula para calcular el perímetro de un triángulo equilátero inscrito en una circunferencia con el radio \( r \).}
{\( 3\sqrt3 r \)}{\( \sqrt3 r \)}{\( 3 r \)}{\( \frac{3\sqrt3 r}2 \)}{}{}}
\End

\Data\ID{1103077102}
\Updated{2022-01-16 09:01:46}
\Level{B}
\SubArea{Circles}
\IMG{1103077102.pdf}
\LANGen{
\Question{
Derive the formula for calculating the perimeter of an equilateral triangle circumscribed about a circle of radius \( r \).}
{\( 6\sqrt3 r \)}{\( 2\sqrt3 r \)}{\( \frac{2r}{\sqrt3} \)}{\( \frac{6r}{\sqrt3} \)}{}{}}
\LANGpl{
\Question{
Wyprowadź wzór, który pozwoli obliczyć obwód trójkąta równobocznego opisanego na okręgu o promieniu \( r \).}
{\( 6\sqrt3 r \)}{\( 2\sqrt3 r \)}{\( \frac{2r}{\sqrt3} \)}{\( \frac{6r}{\sqrt3} \)}{}{}}
\LANGcs{
\Question{
Odvoďte vztah pro výpočet obvodu rovnostranného trojúhelníku, do kterého je vepsaná kružnice s poloměrem  \( r \).}
{\( 6\sqrt3 r \)}{\( 2\sqrt3 r \)}{\( \frac{2r}{\sqrt3} \)}{\( \frac{6r}{\sqrt3} \)}{}{}}
\LANGsk{
\Question{
Odvoďte vzťah  pre výpočet obvodu rovnostranného trojuholníka, do ktorého je vpísaná kružnica  s polomerom  \( r \).}
{\( 6\sqrt3 r \)}{\( 2\sqrt3 r \)}{\( \frac{2r}{\sqrt3} \)}{\( \frac{6r}{\sqrt3} \)}{}{}}
\LANGes{
\Question{
Averigua la fórmula para calcular el perímetro de un triángulo equilátero circunscrito alrededor de una circunferencia con el radio \( r \).}
{\( 6\sqrt3 r \)}{\( 2\sqrt3 r \)}{\( \frac{2r}{\sqrt3} \)}{\( \frac{6r}{\sqrt3} \)}{}{}}
\End

\Data\ID{1103077103}
\Updated{2022-01-16 09:01:27}
\Level{B}
\SubArea{Circles}
\IMG{1103077103.pdf}
\LANGen{
\Question{
The length of the shortest diagonal in a regular polygon is \( 8\,\mathrm{cm} \). The measure of the angle between this diagonal and the side of the polygon is \( 20^{\circ} \). Calculate the radius of a circle circumscribed about this polygon. Round the result to two decimal places.}
{\( 6.22\,\mathrm{cm} \)}{\( 5.22\,\mathrm{cm} \)}{\( 4.26\,\mathrm{cm} \)}{\( 11.69\,\mathrm{cm} \)}{}{}}
\LANGpl{
\Question{
Długość najkrótszej przekątnej wielokąta foremnego wynosi \( 8\,\mathrm{cm} \).  Miara kąta pomiędzy tą przekątną a bokiem tego wielokąta wynosi  \( 20^{\circ} \). Oblicz promień okręgu opisanego na tym wielokącie. Zaokrąglij do dwóch miejsc dziesiętnych.}
{\( 6{,}22\,\mathrm{cm} \)}{\( 5{,}22\,\mathrm{cm} \)}{\( 4{,}26\,\mathrm{cm} \)}{\( 11{,}69\,\mathrm{cm} \)}{}{}}
\LANGcs{
\Question{
V pravidelném mnohoúhelníku má nejkratší úhlopříčka délku \( 8\,\mathrm{cm} \). Velikost úhlu, který svírá tato úhlopříčka se stranou mnohoúhelníku, je \( 20^{\circ} \). Vypočítejte poloměr kružnice, která je tomuto mnohoúhelníku opsaná. Výsledek zaokrouhlete na dvě desetinná místa.}
{\( 6{,}22\,\mathrm{cm} \)}{\( 5{,}22\,\mathrm{cm} \)}{\( 4{,}26\,\mathrm{cm} \)}{\( 11{,}69\,\mathrm{cm} \)}{}{}}
\LANGsk{
\Question{
V pravidelnom mnohouholníku má najkratšia uhlopriečka dĺžku  \( 8\,\mathrm{cm} \). Veľkosť uhla, ktorý zviera táto uhlopriečka so stranou mnohouholníka je \( 20^{\circ} \). Vypočítajte polomer kružnice, ktorá je tomuto mnohouholníku opísaná. Výsledok zaokrúhlite na dve desatinné miesta.}
{\( 6{,}22\,\mathrm{cm} \)}{\( 5{,}22\,\mathrm{cm} \)}{\( 4{,}26\,\mathrm{cm} \)}{\( 11{,}69\,\mathrm{cm} \)}{}{}}
\LANGes{
\Question{
La longitud de la diagonal más corta en un polígono regular es \( 8\,\mathrm{cm} \). El ángulo entre esta diagonal y el lado del polígono mide \( 20^{\circ} \). Calcula el radio de la circunferencia circunscrita en este polígono. Redondea el resultado a dos decimales.}
{\( 6.22\,\mathrm{cm} \)}{\( 5.22\,\mathrm{cm} \)}{\( 4.26\,\mathrm{cm} \)}{\( 11.69\,\mathrm{cm} \)}{}{}}
\End

\Data\ID{1103077104}
\Updated{2022-01-16 09:01:49}
\Level{B}
\SubArea{Circles}
\IMG{1103077104.pdf}
\LANGen{
\Question{
Three equal circles, each of radius \( 6\,\mathrm{cm} \), touch each other as shown in the figure. Find the area of the region bounded by the circles. Round the result to one decimal place.}
{\( 5.8\,\mathrm{cm}^2 \)}{\( 62.3\,\mathrm{cm}^2 \)}{\( 6.2\,\mathrm{cm}^2 \)}{\( 8.4\,\mathrm{cm}^2 \)}{}{}}
\LANGpl{
\Question{
Trzy jednakowe okręgi, każdy o promieniu \( 6\,\mathrm{cm} \), stykają się tak jak przedstawiono na rysunku. Wyznacz powierzchnię obszaru objętego przez te okręgi. Zaokrąglij do jednego miejsca po przecinku.}
{\( 5{,}8\,\mathrm{cm}^2 \)}{\( 62{,}3\,\mathrm{cm}^2 \)}{\( 6{,}2\,\mathrm{cm}^2 \)}{\( 8{,}4\,\mathrm{cm}^2 \)}{}{}}
\LANGcs{
\Question{
Tři shodné kružnice s poloměrem \( 6\,\mathrm{cm} \) se navzájem dotýkají. Určete obsah plochy ležící mezi kružnicemi. Výsledek zaokrouhlete na jedno desetinné místo.}
{\( 5{,}8\,\mathrm{cm}^2 \)}{\( 62{,}3\,\mathrm{cm}^2 \)}{\( 6{,}2\,\mathrm{cm}^2 \)}{\( 8{,}4\,\mathrm{cm}^2 \)}{}{}}
\LANGsk{
\Question{
Tri rovnaké kružnice s polomerom  \( 6\,\mathrm{cm} \) sa navzájom dotýkajú. Určite obsah plochy ležiacej medzi kružnicami. Výsledok zaokrúhlite na jedno desatinné miesto.}
{\( 5{,}8\,\mathrm{cm}^2 \)}{\( 62{,}3\,\mathrm{cm}^2 \)}{\( 6{,}2\,\mathrm{cm}^2 \)}{\( 8{,}4\,\mathrm{cm}^2 \)}{}{}}
\LANGes{
\Question{
Tres circunferencias iguales con radio de \( 6\,\mathrm{cm} \) se tocan. Calcula el área de la superficie delimitada por las circunferencias. (Mira la imagen). Redondea el resultado a un decimal.}
{\( 5.8\,\mathrm{cm}^2 \)}{\( 62.3\,\mathrm{cm}^2 \)}{\( 6.2\,\mathrm{cm}^2 \)}{\( 8.4\,\mathrm{cm}^2 \)}{}{}}
\End

\Data\ID{1103077105}
\Updated{2020-07-15 03:07:21}
\Level{B}
\SubArea{Circles}
\IMG{1103077105_0.pdf}
\LANGen{
\Question{
In a triangle \( ABC \), \( a=7\,\mathrm{cm} \), \( b=8\,\mathrm{cm} \), \( c=11\,\mathrm{cm} \). What is the radius of a circle circumscribed about this triangle? Round the result to two decimal places.}
{\( 5.51\,\mathrm{cm} \)}{\( 6.11\,\mathrm{cm} \)}{\( 4.92\,\mathrm{cm} \)}{\( 6.52\,\mathrm{cm} \)}{}{}}
\LANGpl{
\Question{
W trójkącie  \( ABC \), \( a=7\,\mathrm{cm} \), \( b=8\,\mathrm{cm} \), \( c=11\,\mathrm{cm} \). Jaki jest promień okręgu opisanego na tym trójkącie? Zaokrąglij do dwóch miejsc po przecinku.}
{\( 5{,}51\,\mathrm{cm} \)}{\( 6{,}11\,\mathrm{cm} \)}{\( 4{,}92\,\mathrm{cm} \)}{\( 6{,}52\,\mathrm{cm} \)}{}{}}
\LANGcs{
\Question{
V trojúhelníku \( ABC \), \( a=7\,\mathrm{cm} \), \( b=8\,\mathrm{cm} \), \( c=11\,\mathrm{cm} \). Jaký poloměr má kružnice opsaná tomuto trojúhelníku? Výsledek zaokrouhlete na dvě desetinná místa.}
{\( 5{,}51\,\mathrm{cm} \)}{\( 6{,}11\,\mathrm{cm} \)}{\( 4{,}92\,\mathrm{cm} \)}{\( 6{,}52\,\mathrm{cm} \)}{}{}}
\LANGsk{
\Question{
V trojuholníku \( ABC \), \( a=7\,\mathrm{cm} \), \( b=8\,\mathrm{cm} \), \( c=11\,\mathrm{cm} \). Aký polomer má kružnica opísaná tomuto trojuholníku? Výsledok zaokrúhlite na dve desatinné miesta.}
{\( 5{,}51\,\mathrm{cm} \)}{\( 6{,}11\,\mathrm{cm} \)}{\( 4{,}92\,\mathrm{cm} \)}{\( 6{,}52\,\mathrm{cm} \)}{}{}}
\LANGes{
\Question{
Dado el triángulo \( ABC \), \( a=7\,\mathrm{cm} \), \( b=8\,\mathrm{cm} \), \( c=11\,\mathrm{cm} \). ¿Cuánto mide el radio de la circunferencia circunscrita al triángulo? Redondea el resultado a dos decimales.}
{\( 5.51\,\mathrm{cm} \)}{\( 6.11\,\mathrm{cm} \)}{\( 4.92\,\mathrm{cm} \)}{\( 6.52\,\mathrm{cm} \)}{}{}}
\End

\Data\ID{1103077106}
\Updated{2020-07-15 03:07:21}
\Level{B}
\SubArea{Circles}
\IMG{1103077106.pdf}
\LANGen{
\Question{
Let an equilateral triangle have a side of length of \( 10\,\mathrm{cm} \). Suppose there is a circular sector inside the triangle that has the centre at one of the vertices of the triangle, and the arc touches the opposite side (see the picture). Calculate the length of the arc of the sector. Round the result to two decimal places.}
{\( 9.07\,\mathrm{cm} \)}{\( 8.62\,\mathrm{cm} \)}{\( 8.93\,\mathrm{cm} \)}{\( 9.05\,\mathrm{cm} \)}{}{}}
\LANGpl{
\Question{
Dany jest trójkąt równoboczny o boku długości \( 10\,\mathrm{cm} \). Załóżmy, że wewnatrz trójkąta jest obszar wycinek koła, którego środek znajduje się na jednym z wierzchołków tego trójkąta, a łuk dotyka się przeciwnego boku (rysunek poniżej). Oblicz długość łuku tego wycinka koła. Zaokrąglij do dwóch miejsc dziesiętnych.}
{\( 9{,}07\,\mathrm{cm} \)}{\( 8{,}62\,\mathrm{cm} \)}{\( 8{,}93\,\mathrm{cm} \)}{\( 9{,}05\,\mathrm{cm} \)}{}{}}
\LANGcs{
\Question{
Je dán rovnostranný trojúhelník s délkou strany \( 10\,\mathrm{cm} \). Do trojúhelníku je vepsaná kruhová výseč, jejíž střed je v jednom z vrcholů trojúhelníku a oblouk se dotýká protilehlé strany. Vypočítejte délku oblouku dané výseče. Výsledek zaokrouhlete na dvě desetinná místa.}
{\( 9{,}07\,\mathrm{cm} \)}{\( 8{,}62\,\mathrm{cm} \)}{\( 8{,}93\,\mathrm{cm} \)}{\( 9{,}05\,\mathrm{cm} \)}{}{}}
\LANGsk{
\Question{
Daný je rovnostranný trojuholník s dĺžkou strany  \( 10\,\mathrm{cm} \). Do trojuholníka je vpísaný kruhový výsek, ktorého stred je v jednom z vrcholov trojuholníka a oblúk sa dotýka protiľahlej strany.  Vypočítajte dĺžku oblúka daného výseku. Výsledok zaokrúhlite na dve desatinné miesta.}
{\( 9{,}07\,\mathrm{cm} \)}{\( 8{,}62\,\mathrm{cm} \)}{\( 8{,}93\,\mathrm{cm} \)}{\( 9{,}05\,\mathrm{cm} \)}{}{}}
\LANGes{
\Question{
Dado un triángulo equilátero cuyo lado mide \( 10\,\mathrm{cm} \). En el triángulo se inscribe un sector circular cuyo centro es uno de los vértices del triángulo y el arco toca el lado opuesto (mira la imagen). Calcula la longitud del arco del sector circular. Redondea el resultado a dos decimales.}
{\( 9.07\,\mathrm{cm} \)}{\( 8.62\,\mathrm{cm} \)}{\( 8.93\,\mathrm{cm} \)}{\( 9.05\,\mathrm{cm} \)}{}{}}
\End

\Data\ID{1103077107}
\Updated{2022-01-16 09:01:10}
\Level{B}
\SubArea{Circles}
\IMG{1103077107.pdf}
\LANGen{
\Question{
The figure shows an equilateral triangle whose side is \( 10\,\mathrm{cm} \) long. The circular sector inside the triangle has the centre at one of the vertices of the triangle, and the arc touches the opposite side. Find the ratio of the circumference of the sector to the perimeter of the triangle. Round the result to one decimal place.}
{\( 0.9 \)}{\( 0.5 \)}{\( 0.8 \)}{\( 1.5 \)}{}{}}
\LANGpl{
\Question{
Wycinek koła wewnątrz trójkąta ma swój środek na jednym z wierzchołków tego trójkąta, a łuk dotyka się przeciwnego boku. Wyznacz stosunek obwodu tego wycinka koła do obwodu trójkąta. Zaokrąglij do jednego miejsca po przecinku.}
{\( 0{,}9 \)}{\( 0{,}5 \)}{\( 0{,}8 \)}{\( 1{,}5 \)}{}{}}
\LANGcs{
\Question{
Je dán rovnostranný trojúhelník s délkou strany \( 10\,\mathrm{cm} \). Do trojúhelníku je vepsaná kruhová výseč, jejíž střed je v jednom z vrcholů trojúhelníku a oblouk se dotýká protilehlé strany. Vypočítejte poměr obvodu výseče k obvodu trojúhelníku. Výsledek zaokrouhlete na jedno desetinné místo.}
{\( 0{,}9 \)}{\( 0{,}5 \)}{\( 0{,}8 \)}{\( 1{,}5 \)}{}{}}
\LANGsk{
\Question{
Daný je rovnostranný trojuholník s dĺžkou strany  \( 10\,\mathrm{cm} \). Do trojuholníka je vpísaný kruhový výsek, ktorého stred je v jednom z vrcholov trojuholníka a oblúk sa dotýka protiľahlej strany.  Vypočítajte pomer obvodu výseku ku obvodu trojuholníka. Výsledok zaokrúhlite na jedno desatinné miesto.}
{\( 0{,}9 \)}{\( 0{,}5 \)}{\( 0{,}8 \)}{\( 1{,}5 \)}{}{}}
\LANGes{
\Question{
Dado un triángulo equilátero cuyo lado mide \( 10\,\mathrm{cm} \). En el triángulo se inscribe un sector circular cuyo centro es uno de los vértices del triángulo y el arco toca el lado opuesto (mira la imagen). Halla la razón entre el perímetro del sector circular y el perímetro del triángulo. Redondea el resultado a un decimal.}
{\( 0.9 \)}{\( 0.5 \)}{\( 0.8 \)}{\( 1.5 \)}{}{}}
\End

\Data\ID{1103077108}
\Updated{2020-07-15 03:07:21}
\Level{B}
\SubArea{Circles}
\IMG{1103077108.pdf}
\LANGen{
\Question{
The figure shows an equilateral triangle whose side is \( 10\,\mathrm{cm} \) long. The circular sector inside the triangle has the centre at one of the vertices of the triangle, and the arc touches the opposite side. Calculate the area of the sector. Round the result to one decimal place.}
{\( 39.3\,\mathrm{cm}^2 \)}{\( 37.5\,\mathrm{cm}^2 \)}{\( 14.4\,\mathrm{cm}^2 \)}{\( 3.75\,\mathrm{cm}^2 \)}{}{}}
\LANGpl{
\Question{
Rysunek przedstawia trójkąt równoboczny którego bok ma długość  \( 10\,\mathrm{cm} \). Wycinek koła wewnątrz tego trójkąta ma swój środek na jednym z wierzchołków trójkąta, a łuk dotyka się przeciwległego boku. Oblicz powierzchnię tego obszaru. Zaokrąglij do jednego miejsca po przecinku.}
{\( 39{,}3\,\mathrm{cm}^2 \)}{\( 37{,}5\,\mathrm{cm}^2 \)}{\( 14{,}4\,\mathrm{cm}^2 \)}{\( 3{,}75\,\mathrm{cm}^2 \)}{}{}}
\LANGcs{
\Question{
Je dán rovnostranný trojúhelník s délkou strany \( 10\,\mathrm{cm} \). Do trojúhelníku je vepsaná kruhová výseč, jejíž střed je v jednom z vrcholů trojúhelníku a oblouk se dotýká protilehlé strany. Vypočítejte obsah dané výseče. Výsledek zaokrouhlete na jedno desetinné místo.}
{\( 39{,}3\,\mathrm{cm}^2 \)}{\( 37{,}5\,\mathrm{cm}^2 \)}{\( 14{,}4\,\mathrm{cm}^2 \)}{\( 3{,}75\,\mathrm{cm}^2 \)}{}{}}
\LANGsk{
\Question{
Daný je rovnostranný trojuholník s dĺžkou strany  \( 10\,\mathrm{cm} \). Do trojuholníka je vpísaný kruhový výsek, ktorého stred je v jednom z vrcholov trojuholníka a oblúk sa dotýka protiľahlej strany.  Vypočítajte obsah daného výseku. Výsledok zaokrúhlite na jedno desatinné miesto.}
{\( 39{,}3\,\mathrm{cm}^2 \)}{\( 37{,}5\,\mathrm{cm}^2 \)}{\( 14{,}4\,\mathrm{cm}^2 \)}{\( 3{,}75\,\mathrm{cm}^2 \)}{}{}}
\LANGes{
\Question{
Dado un triángulo equilátero cuyo lado mide \( 10\,\mathrm{cm} \). En el triángulo se inscribe un sector circular cuyo centro es uno de los vértices del triángulo y el arco toca el lado opuesto (mira la imagen). Calcula el área del sector circular. Redondea el resultado a un decimal.}
{\( 39.3\,\mathrm{cm}^2 \)}{\( 37.5\,\mathrm{cm}^2 \)}{\( 14.4\,\mathrm{cm}^2 \)}{\( 3.75\,\mathrm{cm}^2 \)}{}{}}
\End

\Data\ID{1103077109}
\Updated{2020-07-15 03:07:21}
\Level{B}
\SubArea{Circles}
\IMG{1103077109.pdf}
\LANGen{
\Question{
Two quarter circles are inscribed in the square with sides of \( 2\,\mathrm{dm} \). The centres of the quarter circles are at the opposite vertices of the square (see the picture). Calculate the area of the region between the quarter circles. Round the result to two decimal places.}
{\( 2.28\,\mathrm{dm}^2 \)}{\( 3.14\,\mathrm{dm}^2 \)}{\( 21.12\,\mathrm{dm}^2 \)}{\( 1.72\,\mathrm{dm}^2 \)}{}{}}
\LANGpl{
\Question{
Dwie ćwiartki zostały wpisane w kwadrat o boku \( 2\,\mathrm{dm} \).  Środki tych ćwiartek znajdują się na przeciwległych wierzchołkach kwadratu (patrz rysunek). Oblicz powierzchnię obszaru pomiędzy ćwiartkami. Zaokrąglij do dwóch miejsc po przecinku.}
{\( 2{,}28\,\mathrm{dm}^2 \)}{\( 3{,}14\,\mathrm{dm}^2 \)}{\( 21{,}12\,\mathrm{dm}^2 \)}{\( 1{,}72\,\mathrm{dm}^2 \)}{}{}}
\LANGcs{
\Question{
Do čtverce o délce strany \( 2\,\mathrm{dm} \) jsou vepsané dvě čtvrtkružnice se středy v protilehlých vrcholech čtverce. Vypočítejte obsah vyznačené části čtverce ohraničené dvěma čtvrtkružnicemi. Výsledek uveďte s přesností na dvě desetinná místa.}
{\( 2{,}28\,\mathrm{dm}^2 \)}{\( 3{,}14\,\mathrm{dm}^2 \)}{\( 21{,}12\,\mathrm{dm}^2 \)}{\( 1{,}72\,\mathrm{dm}^2 \)}{}{}}
\LANGsk{
\Question{
Do štvorca so stranou dlhou  \( 2\,\mathrm{dm} \). sú vpísané dve štvrťkružnice so stredmi  v protiľahlých vrcholoch štvorca. Vypočítajte obsah vyznačenej časti štvorca, ohraničenej dvoma štvrťkružnicami. Výsledok uveďte s presnosťou na dve desatinné miesta.}
{\( 2{,}28\,\mathrm{dm}^2 \)}{\( 3{,}14\,\mathrm{dm}^2 \)}{\( 21{,}12\,\mathrm{dm}^2 \)}{\( 1{,}72\,\mathrm{dm}^2 \)}{}{}}
\LANGes{
\Question{
En un cuadrado, cuyo lado mide \( 2\,\mathrm{dm} \), se inscriben dos cuartos de circunferencias. Sus centros están en los vértices opuestos del cuadrado (mira la imagen). Calcula el área de la superficie entre los cuartos de las circunferencias. Redondea el resultado a dos decimales.}
{\( 2.28\,\mathrm{dm}^2 \)}{\( 3.14\,\mathrm{dm}^2 \)}{\( 21.12\,\mathrm{dm}^2 \)}{\( 1.72\,\mathrm{dm}^2 \)}{}{}}
\End

\Data\ID{1103077110}
\Updated{2022-01-16 09:01:22}
\Level{B}
\SubArea{Circles}
\IMG{1103077110.pdf}
\LANGen{
\Question{
A circular sector with the central angle of \( 60^{\circ} \) has an area of \( 201\,\mathrm{cm}^2 \). Find its radius \( r \). Round the result to one decimal place.}
{\( 19.6\,\mathrm{cm} \)}{\( 384.1\,\mathrm{cm} \)}{\( 22.5\,\mathrm{cm} \)}{\( 123.7\,\mathrm{cm} \)}{}{}}
\LANGpl{
\Question{
Wycinek koła o kącie środkowym \( 60^{\circ} \) ma obszar \( 201\,\mathrm{cm}^2 \). Oblicz jego promień  \( r \). Zaokrąglij do jednego miejsca po przecinku.}
{\( 19{,}6\,\mathrm{cm} \)}{\( 384{,}1\,\mathrm{cm} \)}{\( 22{,}5\,\mathrm{cm} \)}{\( 123{,}7\,\mathrm{cm} \)}{}{}}
\LANGcs{
\Question{
Kruhová výseč se středovým úhlem \( 60^{\circ} \) má obsah  \( 201\,\mathrm{cm}^2 \). Najděte její poloměr \( r \). Výsledek zaokrouhlete na jedno desetinné místo.}
{\( 19{,}6\,\mathrm{cm} \)}{\( 384{,}1\,\mathrm{cm} \)}{\( 22{,}5\,\mathrm{cm} \)}{\( 123{,}7\,\mathrm{cm} \)}{}{}}
\LANGsk{
\Question{
Kruhový výsek so stredovým uhlom \( 60^{\circ} \) má obsah  \( 201\,\mathrm{cm}^2 \). Nájdite jeho polomer \( r \). Výsledok zaokrúhlite na jedno desatinné miesto.}
{\( 19{,}6\,\mathrm{cm} \)}{\( 384{,}1\,\mathrm{cm} \)}{\( 22{,}5\,\mathrm{cm} \)}{\( 123{,}7\,\mathrm{cm} \)}{}{}}
\LANGes{
\Question{
Un sector circular con un ángulo central de \( 60^{\circ} \) tiene un área de \( 201\,\mathrm{cm}^2 \). Calcula el radio \( r \). Redondea el resultado a un decimal.}
{\( 19.6\,\mathrm{cm} \)}{\( 384.1\,\mathrm{cm} \)}{\( 22.5\,\mathrm{cm} \)}{\( 123.7\,\mathrm{cm} \)}{}{}}
\End

\Data\ID{1103077111}
\Updated{2022-01-16 09:01:15}
\Level{B}
\SubArea{Circles}
\IMG{1103077111.pdf}
\LANGen{
\Question{
A piece of land in the shape of a circular sector with central angle of \( 60^{\circ} \) needs to be fenced. We have used \( 10 \) metres of wire mesh on the curved part of the fence. How many conventional metres of mesh still have to be purchased? Round the result to the nearest metre.}
{\( 19\,\mathrm{m} \)}{\( 10\,\mathrm{m} \)}{\( 15\,\mathrm{m} \)}{\( 25\,\mathrm{m} \)}{}{}}
\LANGpl{
\Question{
Kawałek ziemi, który ma wycinka koła z kątem środkowym \( 60^{\circ} \) należy ogrodzić. Zużyliśmy \( 10 \) metrów siatki drucianej na zaokrąglonej części płotu. Ile metrów bieżących siatki należy jeszcze kupić?  Zaokrąglij do pełnych metrów.}
{\( 19\,\mathrm{m} \)}{\( 10\,\mathrm{m} \)}{\( 15\,\mathrm{m} \)}{\( 25\,\mathrm{m} \)}{}{}}
\LANGcs{
\Question{
Pozemek tvaru kruhové výseče je třeba oplotit. Na oblou část plotu jsme použili \( 10 \) m pletiva. Kolik běžných metrů pletiva je ještě třeba koupit, jestliže příslušný středový úhel je \( 60^{\circ} \)? Výsledek zaokrouhlete na celé metry.}
{\( 19\,\mathrm{m} \)}{\( 10\,\mathrm{m} \)}{\( 15\,\mathrm{m} \)}{\( 25\,\mathrm{m} \)}{}{}}
\LANGsk{
\Question{
Pozemok tvaru kruhového výseku treba oplotiť. Na oblú časť plota sme použili \( 10 \) m pletiva. Koľko bežných metrov pletiva ešte treba dokúpiť, ak príslušný stredový uhol je \( 60^{\circ} \)? Výsledok zaokrúhlite na celé metre.}
{\( 19\,\mathrm{m} \)}{\( 10\,\mathrm{m} \)}{\( 15\,\mathrm{m} \)}{\( 25\,\mathrm{m} \)}{}{}}
\LANGes{
\Question{
Necesitamos cercar una parcela que tiene forma de sector circular con un ángulo central de \( 60^{\circ} \). Hemos usado \( 10 \) metros de malla de alambre para la parte redonda. ¿Cuántos metros de malla de alambre tenemos aún que comprar? Redondea el resultado a los metros más cercanos.}
{\( 19\,\mathrm{m} \)}{\( 10\,\mathrm{m} \)}{\( 15\,\mathrm{m} \)}{\( 25\,\mathrm{m} \)}{}{}}
\End

\Data\ID{1103077201}
\Updated{2022-01-16 09:01:07}
\Level{B}
\SubArea{Circles}
\IMG{1103077201.pdf}
\LANGen{
\Question{
The flower bed has the shape of a circle sector of radius \( 3\,\mathrm{m} \) with central angle \( 75^{\circ} \). Calculate the area of this flower bed. Round the result to two decimal places.}
{\( 5.89\,\mathrm{m}^2 \)}{\( 1.96\,\mathrm{m}^2 \)}{\( 11.78\,\mathrm{m}^2 \)}{\( 9.34\,\mathrm{m}^2 \)}{}{}}
\LANGpl{
\Question{
Rabata kwiatowa ma kształt zaokrąglonego wycinka o promieniu  \( 3\,\mathrm{m} \) z kątem środkowym \( 75^{\circ} \). Oblicz powierzchnię tej rabaty. Zaokrąglij do dwóch miejsc po przecinku.}
{\( 5{,}89\,\mathrm{m}^2 \)}{\( 1{,}96\,\mathrm{m}^2 \)}{\( 11{,}78\,\mathrm{m}^2 \)}{\( 9{,}34\,\mathrm{m}^2 \)}{}{}}
\LANGcs{
\Question{
Květinový záhon má tvar kruhové výseče s poloměrem \( 3\,\mathrm{m} \) a středovým úhlem \( 75^{\circ} \). Vypočítejte plochu tohoto záhonu. Výsledek zaokrouhlete na dvě desetinná místa.}
{\( 5{,}89\,\mathrm{m}^2 \)}{\( 1{,}96\,\mathrm{m}^2 \)}{\( 11{,}78\,\mathrm{m}^2 \)}{\( 9{,}34\,\mathrm{m}^2 \)}{}{}}
\LANGsk{
\Question{
Kvetinový záhon má tvar kruhového výseku s polomerom  \( 3\,\mathrm{m} \) a stredovým uhlom \( 75^{\circ} \). Vypočítajte plochu tohto záhona. Výsledok zaokrúhlite na dve desatinné miesta.}
{\( 5{,}89\,\mathrm{m}^2 \)}{\( 1{,}96\,\mathrm{m}^2 \)}{\( 11{,}78\,\mathrm{m}^2 \)}{\( 9{,}34\,\mathrm{m}^2 \)}{}{}}
\LANGes{
\Question{
Un macizo de flores tiene forma de sector circular, cuyo radio mide \( 3\,\mathrm{m} \), con un ángulo central de \( 75^{\circ} \). Calcula el área del macizo de flores. Redondea el resultado a dos decimales.}
{\( 5.89\,\mathrm{m}^2 \)}{\( 1.96\,\mathrm{m}^2 \)}{\( 11.78\,\mathrm{m}^2 \)}{\( 9.34\,\mathrm{m}^2 \)}{}{}}
\End

\Data\ID{1103077202}
\Updated{2020-07-15 03:07:21}
\Level{B}
\SubArea{Circles}
\IMG{1103077202.pdf}
\LANGen{
\Question{
Let \( ABCDEF \) be a regular hexagon. Six circles of equal radii are drawn touching each other with their centres at the hexagon vertices (see the picture). Calculate the area of the coloured region inside the hexagon if you know that the perimeter of the hexagon \( ABCDEF \) is \( 36\,\mathrm{cm} \). Round the result to two decimal places.}
{\( 36.98\,\mathrm{cm}^2 \)}{\( 93.53\,\mathrm{cm}^2 \)}{\( 65.26\,\mathrm{cm}^2 \)}{\( 25.37\,\mathrm{cm}^2 \)}{}{}}
\LANGpl{
\Question{
Niech \( ABCDEF \) będzie sześciokątem foremnym. Sześć kół o jednakowych promieniach narysowano tak, że stykają się z ich środkami na wierzchołkach tego sześciokąta (patrz rysunek). Oblicz powierzchnię  pomalowanego obszaru w środku sześciokąta jeśli wiesz, że obwód sześciokąta  \( ABCDEF \) wynosi \( 36\,\mathrm{cm} \). Zaokrąglij do dwóch miejsc po przecinku.}
{\( 36{,}98\,\mathrm{cm}^2 \)}{\( 93{,}53\,\mathrm{cm}^2 \)}{\( 65{,}26\,\mathrm{cm}^2 \)}{\( 25{,}37\,\mathrm{cm}^2 \)}{}{}}
\LANGcs{
\Question{
Na obrázku je pravidelný šestiúhelník \( ABCDEF \).  Okolo všech jeho vrcholů jsou sestrojené navzájem se dotýkající kružnicové oblouky se shodnými poloměry. Obvod šestiúhelníku \( ABCDEF \) je \( 36\,\mathrm{cm} \). Jaký je obsah vzniklého vybarveného vnitřního útvaru? Výsledek zaokrouhlete na dvě desetinná místa.}
{\( 36{,}98\,\mathrm{cm}^2 \)}{\( 93{,}53\,\mathrm{cm}^2 \)}{\( 65{,}26\,\mathrm{cm}^2 \)}{\( 25{,}37\,\mathrm{cm}^2 \)}{}{}}
\LANGsk{
\Question{
Na obrázku je pravidelný šesťuholník  \( ABCDEF \).  Okolo všetkých jeho vrcholov sú zostrojené navzájom sa dotýkajúce kružnicové oblúky s rovnakými polomermi. Obvod šesťuholníka \( ABCDEF \) je \( 36\,\mathrm{cm} \). Aký je obsah vzniknutého vyfarbeného vnútorného útvaru? Výsledok zaokrúhlite na dve desatinné miesta.}
{\( 36{,}98\,\mathrm{cm}^2 \)}{\( 93{,}53\,\mathrm{cm}^2 \)}{\( 65{,}26\,\mathrm{cm}^2 \)}{\( 25{,}37\,\mathrm{cm}^2 \)}{}{}}
\LANGes{
\Question{
Dado el hexágono regular \( ABCDEF \). Se dibujan seis circunferencias de radios iguales que se tocan entre sí con sus centros en los vértices del hexágono (mira la imagen). Calcula el área de la superficie coloreada dentro del hexágono si sabes que el perímetro del hexágono \( ABCDEF \) es \( 36\,\mathrm{cm} \). Redondea el resultado a dos decimales.}
{\( 36.98\,\mathrm{cm}^2 \)}{\( 93.53\,\mathrm{cm}^2 \)}{\( 65.26\,\mathrm{cm}^2 \)}{\( 25.37\,\mathrm{cm}^2 \)}{}{}}
\End

\Data\ID{1103077203}
\Updated{2022-01-16 09:01:23}
\Level{B}
\SubArea{Circles}
\IMG{1103077203.pdf}
\LANGen{
\Question{
The tip of a minute hand is at distance of \( 15\,\mathrm{mm} \) from the clock centre. Calculate the length of the path the tip travels in \( 42 \) minutes. Round the result to two decimal places.}
{\( 65.97\,\mathrm{mm} \)}{\( 94.20\,\mathrm{mm} \)}{\( 35.27\,\mathrm{mm} \)}{\( 72.12\,\mathrm{mm} \)}{}{}}
\LANGpl{
\Question{
Koniec minutowej wskazówki zegara znajduje się w odległości \( 15\,\mathrm{mm} \) od środka zegara. Oblicz długość ścieżki jaką przebędzie końcówka w ciągu  \( 42 \) minut. Zaokrąglij do dwóch miejsc po przecinku.}
{\( 65{,}97\,\mathrm{mm} \)}{\( 94{,}20\,\mathrm{mm} \)}{\( 35{,}27\,\mathrm{mm} \)}{\( 72{,}12\,\mathrm{mm} \)}{}{}}
\LANGcs{
\Question{
Vzdálenost hrotu minutové ručičky od středu ciferníku je  \( 15\,\mathrm{mm} \). Vypočítejte délku dráhy, kterou urazí hrot ručičky za \( 42 \) minut. Výsledek zaokrouhlete na dvě desetinná místa.}
{\( 65{,}97\,\mathrm{mm} \)}{\( 94{,}20\,\mathrm{mm} \)}{\( 35{,}27\,\mathrm{mm} \)}{\( 72{,}12\,\mathrm{mm} \)}{}{}}
\LANGsk{
\Question{
Vzdialenosť hrotu minútovej ručičky od stredu ciferníku je  \( 15\,\mathrm{mm} \). Vypočítajte dĺžku dráhy, ktorú prejde hrot ručičky za \( 42 \) minút. Výsledok zaokrúhlite na dve desatinné miesta.}
{\( 65{,}97\,\mathrm{mm} \)}{\( 94{,}20\,\mathrm{mm} \)}{\( 35{,}27\,\mathrm{mm} \)}{\( 72{,}12\,\mathrm{mm} \)}{}{}}
\LANGes{
\Question{
La punta de un minutero está a una distancia de \( 15\,\mathrm{mm} \) del centro del reloj. Calcula la longitud del camino que recorre la punta durante \( 42 \) minutos. Redondea el resultado a dos decimales.}
{\( 65.97\,\mathrm{mm} \)}{\( 94.20\,\mathrm{mm} \)}{\( 35.27\,\mathrm{mm} \)}{\( 72.12\,\mathrm{mm} \)}{}{}}
\End

\Data\ID{1103077204}
\Updated{2022-01-16 09:01:00}
\Level{B}
\SubArea{Circles}
\IMG{1103077204.pdf}
\LANGen{
\Question{
Given a circle, the length of the chord \( AB \) is \( 16\,\mathrm{cm} \) and the height \( v \) of the corresponding circular segment is \( 5\,\mathrm{cm} \) (see the picture). Calculate the area of the segment. Round the result to two decimal places.}
{\( 57.29\,\mathrm{cm}^2 \)}{\( 55.12\,\mathrm{cm}^2 \)}{\( 47.12\,\mathrm{cm}^2 \)}{\( 63.12\,\mathrm{cm}^2 \)}{}{}}
\LANGpl{
\Question{
Dany jest okrąg, którego cięciwa \( AB \) równa jest \( 16\,\mathrm{cm} \), a wysokość \( v \) odpowiadającego okrągłego wycinka wynosi \( 5\,\mathrm{cm} \) (patrz rysunek). Oszacuj powierzchnię wycinka. Zaokrąglij wynik do dwóch miejsc po przecinku.}
{\( 57{,}29\,\mathrm{cm}^2 \)}{\( 55{,}12\,\mathrm{cm}^2 \)}{\( 47{,}12\,\mathrm{cm}^2 \)}{\( 63{,}12\,\mathrm{cm}^2 \)}{}{}}
\LANGcs{
\Question{
Je dána kružnice k. Délka její tětivy \( AB \) je \( 16\,\mathrm{cm} \) a výška příslušné kruhové výseče \( v \) je \( 5\,\mathrm{cm} \) (viz obrázek). Vypočítejte obsah dané výseče. Výsledek zaokrouhlete na dvě desetinná místa.}
{\( 57{,}29\,\mathrm{cm}^2 \)}{\( 55{,}12\,\mathrm{cm}^2 \)}{\( 47{,}12\,\mathrm{cm}^2 \)}{\( 63{,}12\,\mathrm{cm}^2 \)}{}{}}
\LANGsk{
\Question{
Daná je kružnica k. Dĺžka jej tetivy \( AB \) je \( 16\,\mathrm{cm} \) a výška príslušného kruhového odseku \( v \) je \( 5\,\mathrm{cm} \) (pozri obrázok). Vypočítajte obsah daného odseku. Výsledok zaokrúhlite na dve desatinné miesta.}
{\( 57{,}29\,\mathrm{cm}^2 \)}{\( 55{,}12\,\mathrm{cm}^2 \)}{\( 47{,}12\,\mathrm{cm}^2 \)}{\( 63{,}12\,\mathrm{cm}^2 \)}{}{}}
\LANGes{
\Question{
Dada una circunferencia cuya cuerda \( AB \) mide \( 16\,\mathrm{cm} \) y la altura \( v \) del sector circular correspondiente mide \( 5\,\mathrm{cm} \) (mira la imagen). Calcula el área del sector circular. Redondea el resultado a dos decimales.}
{\( 57.29\,\mathrm{cm}^2 \)}{\( 55.12\,\mathrm{cm}^2 \)}{\( 47.12\,\mathrm{cm}^2 \)}{\( 63.12\,\mathrm{cm}^2 \)}{}{}}
\End

\Data\ID{1103077205}
\Updated{2022-01-16 09:01:39}
\Level{B}
\SubArea{Circles}
\IMG{1103077205.pdf}
\LANGen{
\Question{
A farmer has a fenced rhombus shaped garden with a side length of \( 4\,\mathrm{m} \). At one corner where the angle between the sides is \( 60^{\circ} \) the farmer tied a goat (see the picture). Of what length has to be the rope  so that the goat  grazes down exactly half the area of the garden? Round the result to one decimal place.}
{\( 3.6\,\mathrm{m} \)}{\( 3.2\,\mathrm{m} \)}{\( 4.1\,\mathrm{m} \)}{\( 2.9\,\mathrm{m} \)}{}{}}
\LANGpl{
\Question{
Farmer posiada ogród w kształcie rombu o boku długości \( 4\,\mathrm{m} \). W jednym rogu rombu, gdzie kąt pomiędzy bokami jest równy  \( 60^{\circ} \), farmer przywiązał kozę (patrz rysunek). Jaka musi być długość liny, by koza zjadła dokładnie połowę trawy w tym ogrodzie. Zaokrąglij do dwóch miejsc po przecinku.}
{\( 3{,}6\,\mathrm{m} \)}{\( 3{,}2\,\mathrm{m} \)}{\( 4{,}1\,\mathrm{m} \)}{\( 2{,}9\,\mathrm{m} \)}{}{}}
\LANGcs{
\Question{
Farmář má oplocenou kosočtvercovou zahradu s délkou strany \( 4\,\mathrm{m} \). Do rohu zahrady, kde strany svírají úhel \( 60^{\circ} \), uvázal kozu (viz obrázek). Jak dlouhý provaz potřeboval na uvázání kozy, aby koza mohla spást přesně polovinu zahrady? Výsledek uveďte s přesností na desetiny.}
{\( 3{,}6\,\mathrm{m} \)}{\( 3{,}2\,\mathrm{m} \)}{\( 4{,}1\,\mathrm{m} \)}{\( 2{,}9\,\mathrm{m} \)}{}{}}
\LANGsk{
\Question{
Farmár má oplotenú kosoštvorcovú  záhradu s dĺžkou strany \( 4\,\mathrm{m} \). Do rohu záhrady, kde strany zvierajú uhol \( 60^{\circ} \)  uviazal kozu (pozri obrázok). Aký dlhý povraz je potrebný na uviazanie kozy, aby koza mohla spásť presne polovicu záhrady? Výsledok uveďte s presnosťou na desatiny.}
{\( 3{,}6\,\mathrm{m} \)}{\( 3{,}2\,\mathrm{m} \)}{\( 4{,}1\,\mathrm{m} \)}{\( 2{,}9\,\mathrm{m} \)}{}{}}
\LANGes{
\Question{
Un agricultor tiene un jardín vallado en forma de rombo cuyo lado mide \( 4\,\mathrm{m} \). En una esquina, donde el ángulo entre los lados mide \( 60^{\circ} \), el agricultor ató una cabra (mira la imagen). ¿Cuánto tiene que medir la cuerda para que la cabra pueda pastar exactamente la mitad del área del jardín? Redondea el resultado a un decimal.}
{\( 3.6\,\mathrm{m} \)}{\( 3.2\,\mathrm{m} \)}{\( 4.1\,\mathrm{m} \)}{\( 2.9\,\mathrm{m} \)}{}{}}
\End

\Data\ID{1103077206}
\Updated{2020-07-15 03:07:21}
\Level{B}
\SubArea{Circles}
\IMG{1103077206.pdf}
\LANGen{
\Question{
The circumference of a semicircle is \( 10\,\mathrm{cm} \). Find its radius.}
{\( \frac{10}{\pi+2}\,\mathrm{cm} \)}{\( \frac{10}{\pi}\,\mathrm{cm} \)}{\( \sqrt{\frac{10}{\pi}}\,\mathrm{cm} \)}{\( \frac{10}{\pi+1}\,\mathrm{cm} \)}{}{}}
\LANGpl{
\Question{
Obwód półokręgu jest równy  \( 10\,\mathrm{cm} \). Wyznacz jego promień.}
{\( \frac{10}{\pi+2}\,\mathrm{cm} \)}{\( \frac{10}{\pi}\,\mathrm{cm} \)}{\( \sqrt{\frac{10}{\pi}}\,\mathrm{cm} \)}{\( \frac{10}{\pi+1}\,\mathrm{cm} \)}{}{}}
\LANGcs{
\Question{
Obvod polokružnice je \( 10\,\mathrm{cm} \). Zjistěte poloměr této polokružnice.}
{\( \frac{10}{\pi+2}\,\mathrm{cm} \)}{\( \frac{10}{\pi}\,\mathrm{cm} \)}{\( \sqrt{\frac{10}{\pi}}\,\mathrm{cm} \)}{\( \frac{10}{\pi+1}\,\mathrm{cm} \)}{}{}}
\LANGsk{
\Question{
Obvod polkruhu je \( 10\,\mathrm{cm} \). Nájdite polomer tohto polkruhu.}
{\( \frac{10}{\pi+2}\,\mathrm{cm} \)}{\( \frac{10}{\pi}\,\mathrm{cm} \)}{\( \sqrt{\frac{10}{\pi}}\,\mathrm{cm} \)}{\( \frac{10}{\pi+1}\,\mathrm{cm} \)}{}{}}
\LANGes{
\Question{
La longitud de una semicircunferencia es \( 10\,\mathrm{cm} \). Calcula el radio.}
{\( \frac{10}{\pi+2}\,\mathrm{cm} \)}{\( \frac{10}{\pi}\,\mathrm{cm} \)}{\( \sqrt{\frac{10}{\pi}}\,\mathrm{cm} \)}{\( \frac{10}{\pi+1}\,\mathrm{cm} \)}{}{}}
\End

\Data\ID{1103077209}
\Updated{2020-07-15 03:07:21}
\Level{B}
\SubArea{Circles}
\IMG{1103077209.pdf}
\LANGen{
\Question{
A semicircle is inscribed in a triangle \( KLM \) so that the diameter of the semicircle is parallel to side \( KL \) (see the picture). The length of \( KL \) is \( 8\,\mathrm{cm} \) and the height to side \( KL \) is \( 4\,\mathrm{cm} \) long. Determine the radius of the semicircle.}
{\( 2\,\mathrm{cm} \)}{\( 4\,\mathrm{cm} \)}{\( 6\,\mathrm{cm} \)}{\( 8\,\mathrm{cm} \)}{}{}}
\LANGpl{
\Question{
Półkole jest wpisane w trójkąt \( KLM \) tak, że średnica jest równoległa do boku  \( KL \) (patrz rysunek). Długość \( KL \) wynosi \( 8\,\mathrm{cm} \), a wysokość do boku  \( KL \) jest równa \( 4\,\mathrm{cm} \). Wyznacz promień półkola.}
{\( 2\,\mathrm{cm} \)}{\( 4\,\mathrm{cm} \)}{\( 6\,\mathrm{cm} \)}{\( 8\,\mathrm{cm} \)}{}{}}
\LANGcs{
\Question{
Do trojúhelníku \( KLM \) je vepsaná polokružnice, přičemž její průměr je rovnoběžný se stranou \( KL \) (viz obrázek).  Délka strany  \( KL \) je \( 8\,\mathrm{cm} \) a výška na stranu  \( KL \) je \( 4\,\mathrm{cm} \). Vypočítejte poloměr polokružnice.}
{\( 2\,\mathrm{cm} \)}{\( 4\,\mathrm{cm} \)}{\( 6\,\mathrm{cm} \)}{\( 8\,\mathrm{cm} \)}{}{}}
\LANGsk{
\Question{
Do trojuholníka  \( KLM \) je vpísaný polkruh, pričom priemer polkruhu je rovnobežný so stranou \( KL \) (pozri obrázok).  Dĺžka strany  \( KL \) je \( 8\,\mathrm{cm} \) a výška na stranu  \( KL \) je \( 4\,\mathrm{cm} \). Vypočítajte polomer polkruhu.}
{\( 2\,\mathrm{cm} \)}{\( 4\,\mathrm{cm} \)}{\( 6\,\mathrm{cm} \)}{\( 8\,\mathrm{cm} \)}{}{}}
\LANGes{
\Question{
Una semicircunferencia se inscribe en el triángulo \( KLM \), suponiendo que el diámetro de la semicircunferencia es paralelo al lado \( KL \) (mira la imagen). La longitud del \( KL \) es \( 8\,\mathrm{cm} \) y la altura sobre el lado \( KL \) mide \( 4\,\mathrm{cm} \). Calcula el radio de la semicircunferencia.}
{\( 2\,\mathrm{cm} \)}{\( 4\,\mathrm{cm} \)}{\( 6\,\mathrm{cm} \)}{\( 8\,\mathrm{cm} \)}{}{}}
\End

\Data\ID{1103077210}
\Updated{2020-07-15 03:07:21}
\Level{B}
\SubArea{Circles}
\IMG{1103077210.pdf}
\LANGen{
\Question{
The picture shows a roundabout with the radius of \( 6\,\mathrm{m} \). Inside the roundabout there is a flower bed, which has the shape of an equilateral triangle inscribed in it. The remaining part inside the roundabout is the lawn. Calculate the area of the lawn.}
{\( 66.33\,\mathrm{cm}^2 \)}{\( 46.77\,\mathrm{cm}^2 \)}{\( 113.10\,\mathrm{cm}^2 \)}{\( 24.66\,\mathrm{cm}^2 \)}{}{}}
\LANGpl{
\Question{
Rysunek przedstawia rondo o promieniu  \( 6\,\mathrm{m} \). W środku ronda znajduje się rabata kwiatowa w kształcie trójkąta równobocznego. Pozostałą część w środku ronda stanowi trawnik. Oblicz powierzchnię trawnika.}
{\( 66{,}33\,\mathrm{cm}^2 \)}{\( 46{,}77\,\mathrm{cm}^2 \)}{\( 113{,}10\,\mathrm{cm}^2 \)}{\( 24{,}66\,\mathrm{cm}^2 \)}{}{}}
\LANGcs{
\Question{
Ve středu kruhového objezdu je ostrůvek ve tvaru kruhu, do kterého je vepsaný rovnostranný trojúhelník. V trojúhelníku jsou vysazené květiny a zbytek ostrůvku tvoří trávník (viz obrázek). Vypočítejte obsah plochy trávníku, jestliže poloměr ostrůvku je \( 6\,\mathrm{m} \).}
{\( 66{,}33\,\mathrm{cm}^2 \)}{\( 46{,}77\,\mathrm{cm}^2 \)}{\( 113{,}10\,\mathrm{cm}^2 \)}{\( 24{,}66\,\mathrm{cm}^2 \)}{}{}}
\LANGsk{
\Question{
V strede kruhového objazdu je ostrovček v tvare kruhu, do ktorého je vpísaný rovnostranný trojuholník. V trojuholníku sú vysadené kvety a zvyšok ostrovčeka tvorí trávnik (pozri obrázok). Vypočítajte obsah plochy trávnika, ak polomer ostrovčeka je  \( 6\,\mathrm{m} \).}
{\( 66{,}33\,\mathrm{cm}^2 \)}{\( 46{,}77\,\mathrm{cm}^2 \)}{\( 113{,}10\,\mathrm{cm}^2 \)}{\( 24{,}66\,\mathrm{cm}^2 \)}{}{}}
\LANGes{
\Question{
La imagen representa una rotonda cuyo radio mide \( 6\,\mathrm{m} \). Dentro de la rotonda hay un macizo de flores que tiene forma de triángulo equilátero inscrito en ella. En la parte restante de la rotonda hay cespéd. Calcula el área del cespéd.}
{\( 66.33\,\mathrm{cm}^2 \)}{\( 46.77\,\mathrm{cm}^2 \)}{\( 113.10\,\mathrm{cm}^2 \)}{\( 24.66\,\mathrm{cm}^2 \)}{}{}}
\End

\Data\ID{2000005904}
\Updated{2022-05-16 11:05:07}
\Level{B}
\SubArea{Circles}
\IMG{2000005901-figure3.pdf}
\LANGen{
\Question{
Find the magnitude of the angle that the diagonals \(DB\) and \(CG\) make in the regular heptagon \(ABCDEFG\). (See the picture.)}
{\( 180^{\circ}-\left(\frac{360^{\circ}}{14} +3\cdot\frac{360^{\circ}}{14}\right)\)}{\( 180^{\circ}-\left(\frac{360^{\circ}}{7} +3\cdot\frac{360^{\circ}}{7}\right)\)}{\( 180^{\circ}-\frac{360^{\circ}}{14} +3\cdot\frac{360^{\circ}}{14}\)}{\( 180^{\circ}-\left(\frac{360^{\circ}}{14} +4\cdot\frac{360^{\circ}}{14}\right)\)}{}{}}
\LANGpl{
\Question{
Wyznacz miarę kąta, jaki tworzą przekątne \(DB\) i \(CG\) w siedmiokącie foremnym \(ABCDEFG\).  (Patrz rysunek.)}
{\( 180^{\circ}-\left(\frac{360^{\circ}}{14} +3\cdot\frac{360^{\circ}}{14}\right)\)}{\( 180^{\circ}-\left(\frac{360^{\circ}}{7} +3\cdot\frac{360^{\circ}}{7}\right)\)}{\( 180^{\circ}-\frac{360^{\circ}}{14} +3\cdot\frac{360^{\circ}}{14}\)}{\( 180^{\circ}-\left(\frac{360^{\circ}}{14} +4\cdot\frac{360^{\circ}}{14}\right)\)}{}{}}
\LANGcs{
\Question{
Vypočtěte velikost úhlu, který svírají úhlopříčky \(DB\) a \(CG\) v pravidelném sedmiúhelníku \(ABCDEFG\), (viz obrázek).}
{\( 180^{\circ}-\left(\frac{360^{\circ}}{14} +3\cdot\frac{360^{\circ}}{14}\right)\)}{\( 180^{\circ}-\left(\frac{360^{\circ}}{7} +3\cdot\frac{360^{\circ}}{7}\right)\)}{\( 180^{\circ}-\frac{360^{\circ}}{14} +3\cdot\frac{360^{\circ}}{14}\)}{\( 180^{\circ}-\left(\frac{360^{\circ}}{14} +4\cdot\frac{360^{\circ}}{14}\right)\)}{}{}}
\LANGsk{
\Question{
Nájdite veľkosť uhla, ktorý zvierajú uhlopriečky  \(DB\) a \(CG\) v pravidelnom šesťuholníku \(ABCDEFG\), (pozri obrázok).}
{\( 180^{\circ}-\left(\frac{360^{\circ}}{14} +3\cdot\frac{360^{\circ}}{14}\right)\)}{\( 180^{\circ}-\left(\frac{360^{\circ}}{7} +3\cdot\frac{360^{\circ}}{7}\right)\)}{\( 180^{\circ}-\frac{360^{\circ}}{14} +3\cdot\frac{360^{\circ}}{14}\)}{\( 180^{\circ}-\left(\frac{360^{\circ}}{14} +4\cdot\frac{360^{\circ}}{14}\right)\)}{}{}}
\LANGes{
\Question{
Dado el heptágono regular \(ABCDEFG\). Calcula la medida del ángulo entre las diagonales \(DB\) y \(CG\).}
{\( 180^{\circ}-\left(\frac{360^{\circ}}{14} +3\cdot\frac{360^{\circ}}{14}\right)\)}{\( 180^{\circ}-\left(\frac{360^{\circ}}{7} +3\cdot\frac{360^{\circ}}{7}\right)\)}{\( 180^{\circ}-\frac{360^{\circ}}{14} +3\cdot\frac{360^{\circ}}{14}\)}{\( 180^{\circ}-\left(\frac{360^{\circ}}{14} +4\cdot\frac{360^{\circ}}{14}\right)\)}{}{}}
\End

\Data\ID{2000005905}
\Updated{2022-05-16 11:05:07}
\Level{B}
\SubArea{Circles}
\IMG{2000005901-figure4.pdf}
\LANGen{
\Question{
Find the magnitude of the angle that the diagonals \(AF\) and \(CG\) make in the regular octagon \(ABCDEFGH\). (See the picture.)}
{\(112.5^{\circ}\)}{\(225^{\circ}\)}{\(45^{\circ}\)}{\(135^{\circ}\)}{}{}}
\LANGpl{
\Question{
Wyznacz miarę kąta, jaki przekątne \(AF\) i \(CG\) tworzą w ośmiokącie foremnym \(ABCDEFGH\). (Patrz rysunek.)}
{\(112{,}5^{\circ}\)}{\(225^{\circ}\)}{\(45^{\circ}\)}{\(135^{\circ}\)}{}{}}
\LANGcs{
\Question{
Vypočtěte velikost úhu, který svírají úhlopříčky \(AF\) a \(CG\) v pravidellném osmiúhelníku \(ABCDEFGH\), (viz obrázek).}
{\(112{,}5^{\circ}\)}{\(225^{\circ}\)}{\(45^{\circ}\)}{\(135^{\circ}\)}{}{}}
\LANGsk{
\Question{
Nájdite veľkosť uhla, ktorý zvierajú uhlopriečky \(AF\) a \(CG\) v pravidelnom osemuholníku \(ABCDEFGH\), (pozri obrázok).}
{\(112{,}5^{\circ}\)}{\(225^{\circ}\)}{\(45^{\circ}\)}{\(135^{\circ}\)}{}{}}
\LANGes{
\Question{
Dado el octógono regular \(ABCDEFGH\). Calcula la medida del ángulo entre las diagonales \(AF\) y \(CG\).}
{\(112.5^{\circ}\)}{\(225^{\circ}\)}{\(45^{\circ}\)}{\(135^{\circ}\)}{}{}}
\End

\Data\ID{2000005906}
\Updated{2022-05-16 11:05:07}
\Level{B}
\SubArea{Circles}
\IMG{2000005901-figure5.pdf}
\LANGen{
\Question{
Find the magnitude of the angle that the diagonals \(AE\) and \(FD\) make in the regular hexagon \(ABCDEF\). (See the picture.)}
{\( 60^{\circ}\)}{\( 30^{\circ}\)}{\( 120^{\circ}\)}{\( 80^{\circ}\)}{}{}}
\LANGpl{
\Question{
Wyznacz miarę kąta, jaki tworzą przekątne \(AE\) i \(FD\) w  sześciokącie foremnym \(ABCDEF\). (Patrz rysunek.)}
{\( 60^{\circ}\)}{\( 30^{\circ}\)}{\( 120^{\circ}\)}{\( 80^{\circ}\)}{}{}}
\LANGcs{
\Question{
Vypočtěte velikost úhlu, který svírají úhlopříčky \(AE\) a \(FD\) v pravidelném šestiúhelníku \(ABCDEF\), (viz obrázek).}
{\( 60^{\circ}\)}{\( 30^{\circ}\)}{\( 120^{\circ}\)}{\( 80^{\circ}\)}{}{}}
\LANGsk{
\Question{
Nájdite veľkosť uhla, ktorý zvierajú uhlopriečky \(AE\) a \(FD\) v pravidelnom šesťuholníku \(ABCDEF\), (pozri obrázok).}
{\( 60^{\circ}\)}{\( 30^{\circ}\)}{\( 120^{\circ}\)}{\( 80^{\circ}\)}{}{}}
\LANGes{
\Question{
Dado el hexágono regular \(ABCDEF\). Calcula la medida del ángulo entre las diagonales \(AE\) y \(FD\).}
{\( 60^{\circ}\)}{\( 30^{\circ}\)}{\( 120^{\circ}\)}{\( 80^{\circ}\)}{}{}}
\End

\Data\ID{2000005907}
\Updated{2022-05-16 11:05:07}
\Level{B}
\SubArea{Circles}
\IMG{2000005901-figure6.pdf}
\LANGen{
\Question{
Find the magnitude of the angle that the diagonals \(BE\) and \(DH\) make in the regular nonagon \(ABCDEFGHI\). (See the picture.)}
{\( 80^{\circ}\)}{\( 60^{\circ}\)}{\( 70^{\circ}\)}{\( 40^{\circ}\)}{}{}}
\LANGpl{
\Question{
Wyznacz miarę kąta, jaki tworzą przekątne \(BE\) i \(DH\) w dziewięnciokącie foremnym \(ABCDEFGHI\). (Patrz rysunek.)}
{\( 80^{\circ}\)}{\( 60^{\circ}\)}{\( 70^{\circ}\)}{\( 40^{\circ}\)}{}{}}
\LANGcs{
\Question{
Vypočtěte velikost úhlu, který svírají úhlopříčky \(BE\) a \(DH\) v pravidelném devítiúhelníku \(ABCDEFGHI\), (viz obrázek).}
{\( 80^{\circ}\)}{\( 60^{\circ}\)}{\( 70^{\circ}\)}{\( 40^{\circ}\)}{}{}}
\LANGsk{
\Question{
Nájdite veľkosť uhla, ktorý zvierajú uhlopriečky \(BE\) a \(DH\) v pravidelnom deväťuholníku \(ABCDEFGHI\), (pozri obrázok).}
{\( 80^{\circ}\)}{\( 60^{\circ}\)}{\( 70^{\circ}\)}{\( 40^{\circ}\)}{}{}}
\LANGes{
\Question{
Dado el eneágono regular \(ABCDEFGHI\). Calcula la medida del ángulo entre las diagonales \(BE\) y \(DH\).}
{\( 80^{\circ}\)}{\( 60^{\circ}\)}{\( 70^{\circ}\)}{\( 40^{\circ}\)}{}{}}
\End

\Data\ID{2000005908}
\Updated{2022-05-23 02:05:13}
\Level{B}
\SubArea{Circles}
\IMG{2000005901-figure7.pdf}
\LANGen{
\Question{
Which of the following formulas gives the area of a regular nonagon inscribed in a circle with radius \(r\)? (See the picture.)}
{\(\frac{9r^2\sin{40^{\circ}}}{2}\)}{\({9r^2\sin{40^{\circ}}}\)}{\(\frac{9r^2\cos{40^{\circ}}}{2}\)}{\(\frac{9r^2\sin{20^{\circ}}}{2}\)}{}{}}
\LANGpl{
\Question{
Który z poniższych wzorów wyznacza pole powierzchni dziewięnciokąta foremnego wpisanego w okrąg o promieniu \(r\)? (Patrz rysunek.)}
{\(\frac{9r^2\sin{40^{\circ}}}{2}\)}{\({9r^2\sin{40^{\circ}}}\)}{\(\frac{9r^2\cos{40^{\circ}}}{2}\)}{\(\frac{9r^2\sin{20^{\circ}}}{2}\)}{}{}}
\LANGcs{
\Question{
Kterou z uvedených rovnic vypočteme obsah pravidelného devítiúhelníku vepsaného do kružnice s poloměrem \(r\)? (Viz obrázek.)}
{\(\frac{9r^2\sin{40^{\circ}}}{2}\)}{\({9r^2\sin{40^{\circ}}}\)}{\(\frac{9r^2\cos{40^{\circ}}}{2}\)}{\(\frac{9r^2\sin{20^{\circ}}}{2}\)}{}{}}
\LANGsk{
\Question{
Ktorý z nasledujúcich vzorcov vyjadruje obsah pravidelného deväťuholníka vpísaného do kružnice s polomerom \(r\), (pozri obrázok)?}
{\(\frac{9r^2\sin{40^{\circ}}}{2}\)}{\({9r^2\sin{40^{\circ}}}\)}{\(\frac{9r^2\cos{40^{\circ}}}{2}\)}{\(\frac{9r^2\sin{20^{\circ}}}{2}\)}{}{}}
\LANGes{
\Question{
¿Cuál de las siguientes fórmulas nos da el área de un eneágono regular inscrito en una circunferencia de radio \(r\)?}
{\(\frac{9r^2\sin{40^{\circ}}}{2}\)}{\({9r^2\sin{40^{\circ}}}\)}{\(\frac{9r^2\cos{40^{\circ}}}{2}\)}{\(\frac{9r^2\sin{20^{\circ}}}{2}\)}{}{}}
\End

\Data\ID{2000005909}
\Updated{2022-08-30 04:08:35}
\Level{B}
\SubArea{Circles}
\IMG{2000005901-figure8.pdf}
\LANGen{
\Question{
The regular octagon \(ABCDEFGH\) is inscribed in the circle. Calculate the magnitudes of the interior angles of the chordal quadrilateral \(HBCF\). (See the picture.)}
{\( \alpha=90^{\circ}\); \( \beta=112.5^{\circ}\); \( \gamma=90^{\circ}\); \( \delta=67.5^{\circ}\)}{\( \alpha=90^{\circ}\); \( \beta=67.5^{\circ}\); \( \gamma=90^{\circ}\); \( \delta=67.5^{\circ}\)}{\( \alpha=90^{\circ}\); \( \beta=122.5^{\circ}\); \( \gamma=80^{\circ}\); \( \delta=67.5^{\circ}\)}{\( \alpha=90^{\circ}\); \( \beta=67.5^{\circ}\); \( \gamma=90^{\circ}\); \( \delta=112.5^{\circ}\)}{}{}}
\LANGpl{
\Question{
Ośmiokąt foremny \(ABCDEFGH\) jest wpisany w okrąg. Oblicz miary kątów wewnętrznych czworokąta \(HBCF\). (Patrz rysunek.)}
{\( \alpha=90^{\circ}\); \( \beta=112{,}5^{\circ}\); \( \gamma=90^{\circ}\); \( \delta=67{,}5^{\circ}\)}{\( \alpha=90^{\circ}\); \( \beta=67{,}5^{\circ}\); \( \gamma=90^{\circ}\); \( \delta=67{,}5^{\circ}\)}{\( \alpha=90^{\circ}\); \( \beta=122{,}5^{\circ}\); \( \gamma=80^{\circ}\); \( \delta=67{,}5^{\circ}\)}{\( \alpha=90^{\circ}\); \( \beta=67{,}5^{\circ}\); \( \gamma=90^{\circ}\); \( \delta=112{,}5^{\circ}\)}{}{}}
\LANGcs{
\Question{
Je dán pravidelný osmiúhelník \(ABCDEFGH\) vepsaný do kružnice. Vypočtěte velikost vnitřních úhlů tětivového čtyřúhelníku \(HBCF\), (viz obrázek).}
{\( \alpha=90^{\circ}\); \( \beta=112{,}5^{\circ}\); \( \gamma=90^{\circ}\); \( \delta=67{,}5^{\circ}\)}{\( \alpha=90^{\circ}\); \( \beta=67{,}5^{\circ}\); \( \gamma=90^{\circ}\); \( \delta=67{,}5^{\circ}\)}{\( \alpha=90^{\circ}\); \( \beta=122{,}5^{\circ}\); \( \gamma=80^{\circ}\); \( \delta=67{,}5^{\circ}\)}{\( \alpha=90^{\circ}\); \( \beta=67{,}5^{\circ}\); \( \gamma=90^{\circ}\); \( \delta=112{,}5^{\circ}\)}{}{}}
\LANGsk{
\Question{
Do kružnice je vpísaný pravidelný 8-uholník \(ABCDEFGH\). Vypočítajte veľkosti vnútorných uhlov tetivového štvoruholníka \(HBCF\), pozri obrázok.}
{\( \alpha=90^{\circ}\); \( \beta=112{,}5^{\circ}\); \( \gamma=90^{\circ}\); \( \delta=67{,}5^{\circ}\)}{\( \alpha=90^{\circ}\); \( \beta=67{,}5^{\circ}\); \( \gamma=90^{\circ}\); \( \delta=67{,}5^{\circ}\)}{\( \alpha=90^{\circ}\); \( \beta=122{,}5^{\circ}\); \( \gamma=80^{\circ}\); \( \delta=67{,}5^{\circ}\)}{\( \alpha=90^{\circ}\); \( \beta=67{,}5^{\circ}\); \( \gamma=90^{\circ}\); \( \delta=112{,}5^{\circ}\)}{}{}}
\LANGes{
\Question{
Dado el octógono regular \(ABCDEFGH\) inscrito en una circunferencia. Calcula las medidas de los ángulos interiores del cuadrilátero inscrito \(HBCF\).}
{\( \alpha=90^{\circ}\); \( \beta=112.5^{\circ}\); \( \gamma=90^{\circ}\); \( \delta=67.5^{\circ}\)}{\( \alpha=90^{\circ}\); \( \beta=67.5^{\circ}\); \( \gamma=90^{\circ}\); \( \delta=67.5^{\circ}\)}{\( \alpha=90^{\circ}\); \( \beta=122.5^{\circ}\); \( \gamma=80^{\circ}\); \( \delta=67.5^{\circ}\)}{\( \alpha=90^{\circ}\); \( \beta=67.5^{\circ}\); \( \gamma=90^{\circ}\); \( \delta=112.5^{\circ}\)}{}{}}
\End

\Data\ID{2000005910}
\Updated{2022-08-30 04:08:59}
\Level{B}
\SubArea{Circles}
\IMG{2000005901-figure9.pdf}
\LANGen{
\Question{
The regular heptagon is inscribed in a circle. Calculate the magnitudes of the interior angles of the chordal quadrilateral \(ACEG\). (See the picture.)}
{\( \alpha=4\cdot\frac{360^{\circ}}{14}\); \( \beta=3\cdot\frac{360^{\circ}}{14}\); \( \gamma=3\cdot\frac{360^{\circ}}{14}\); \( \delta=4\cdot\frac{360^{\circ}}{14}\)}{\( \alpha=4\cdot\frac{360^{\circ}}{7}\); \( \beta=3\cdot\frac{360^{\circ}}{7}\); \( \gamma=3\cdot\frac{360^{\circ}}{7}\); \( \delta=4\cdot\frac{360^{\circ}}{7}\)}{\( \alpha=4\cdot\frac{180^{\circ}}{14}\); \( \beta=3\cdot\frac{180^{\circ}}{14}\); \( \gamma=3\cdot\frac{180^{\circ}}{14}\); \( \delta=4\cdot\frac{180^{\circ}}{14}\)}{\( \alpha=4\cdot\frac{360^{\circ}}{14}\); \( \beta=4\cdot\frac{360^{\circ}}{14}\); \( \gamma=3\cdot\frac{360^{\circ}}{14}\); \( \delta=3\cdot\frac{360^{\circ}}{14}\)}{}{}}
\LANGpl{
\Question{
Siedmiokąt foremny jest wpisany w okrąg. Oblicz wielkości kątów wewnętrznych czworokąta \(ACEG\). (Patrz rysunek.)}
{\( \alpha=4\cdot\frac{360^{\circ}}{14}\); \( \beta=3\cdot\frac{360^{\circ}}{14}\); \( \gamma=3\cdot\frac{360^{\circ}}{14}\); \( \delta=4\cdot\frac{360^{\circ}}{14}\)}{\( \alpha=4\cdot\frac{360^{\circ}}{7}\); \( \beta=3\cdot\frac{360^{\circ}}{7}\); \( \gamma=3\cdot\frac{360^{\circ}}{7}\); \( \delta=4\cdot\frac{360^{\circ}}{7}\)}{\( \alpha=4\cdot\frac{180^{\circ}}{14}\); \( \beta=3\cdot\frac{180^{\circ}}{14}\); \( \gamma=3\cdot\frac{180^{\circ}}{14}\); \( \delta=4\cdot\frac{180^{\circ}}{14}\)}{\( \alpha=4\cdot\frac{360^{\circ}}{14}\); \( \beta=4\cdot\frac{360^{\circ}}{14}\); \( \gamma=3\cdot\frac{360^{\circ}}{14}\); \( \delta=3\cdot\frac{360^{\circ}}{14}\)}{}{}}
\LANGcs{
\Question{
Je dán pravidelný sedmiúhelník vepsaný do kružnice. Vypočtěte velikost vnitřních úhlů tětivového čtyřúhelníku \(ACEG\), (viz obrázek).}
{\( \alpha=4\cdot\frac{360^{\circ}}{14}\); \( \beta=3\cdot\frac{360^{\circ}}{14}\); \( \gamma=3\cdot\frac{360^{\circ}}{14}\); \( \delta=4\cdot\frac{360^{\circ}}{14}\)}{\( \alpha=4\cdot\frac{360^{\circ}}{7}\); \( \beta=3\cdot\frac{360^{\circ}}{7}\); \( \gamma=3\cdot\frac{360^{\circ}}{7}\); \( \delta=4\cdot\frac{360^{\circ}}{7}\)}{\( \alpha=4\cdot\frac{180^{\circ}}{14}\); \( \beta=3\cdot\frac{180^{\circ}}{14}\); \( \gamma=3\cdot\frac{180^{\circ}}{14}\); \( \delta=4\cdot\frac{180^{\circ}}{14}\)}{\( \alpha=4\cdot\frac{360^{\circ}}{14}\); \( \beta=4\cdot\frac{360^{\circ}}{14}\); \( \gamma=3\cdot\frac{360^{\circ}}{14}\); \( \delta=3\cdot\frac{360^{\circ}}{14}\)}{}{}}
\LANGsk{
\Question{
Do kružnice je vpísaný pravidelný sedemuholník \(ABCDEFG\). Vypočítajte veľkosti vnútorných uhlov tetivového štvoruholníka \(ACEG\), pozri obrázok.}
{\( \alpha=4\cdot\frac{360^{\circ}}{14}\); \( \beta=3\cdot\frac{360^{\circ}}{14}\); \( \gamma=3\cdot\frac{360^{\circ}}{14}\); \( \delta=4\cdot\frac{360^{\circ}}{14}\)}{\( \alpha=4\cdot\frac{360^{\circ}}{7}\); \( \beta=3\cdot\frac{360^{\circ}}{7}\); \( \gamma=3\cdot\frac{360^{\circ}}{7}\); \( \delta=4\cdot\frac{360^{\circ}}{7}\)}{\( \alpha=4\cdot\frac{180^{\circ}}{14}\); \( \beta=3\cdot\frac{180^{\circ}}{14}\); \( \gamma=3\cdot\frac{180^{\circ}}{14}\); \( \delta=4\cdot\frac{180^{\circ}}{14}\)}{\( \alpha=4\cdot\frac{360^{\circ}}{14}\); \( \beta=4\cdot\frac{360^{\circ}}{14}\); \( \gamma=3\cdot\frac{360^{\circ}}{14}\); \( \delta=3\cdot\frac{360^{\circ}}{14}\)}{}{}}
\LANGes{
\Question{
Dado el heptágono regular inscrito en una circunferencia. Calcula las medidas de los ángulos interiores del cuadrilátero inscrito \(ACEG\).}
{\( \alpha=4\cdot\frac{360^{\circ}}{14}\); \( \beta=3\cdot\frac{360^{\circ}}{14}\); \( \gamma=3\cdot\frac{360^{\circ}}{14}\); \( \delta=4\cdot\frac{360^{\circ}}{14}\)}{\( \alpha=4\cdot\frac{360^{\circ}}{7}\); \( \beta=3\cdot\frac{360^{\circ}}{7}\); \( \gamma=3\cdot\frac{360^{\circ}}{7}\); \( \delta=4\cdot\frac{360^{\circ}}{7}\)}{\( \alpha=4\cdot\frac{180^{\circ}}{14}\); \( \beta=3\cdot\frac{180^{\circ}}{14}\); \( \gamma=3\cdot\frac{180^{\circ}}{14}\); \( \delta=4\cdot\frac{180^{\circ}}{14}\)}{\( \alpha=4\cdot\frac{360^{\circ}}{14}\); \( \beta=4\cdot\frac{360^{\circ}}{14}\); \( \gamma=3\cdot\frac{360^{\circ}}{14}\); \( \delta=3\cdot\frac{360^{\circ}}{14}\)}{}{}}
\End

\Data\ID{2010012809}
\Updated{2022-08-25 10:08:12}
\Level{B}
\SubArea{Circles}
\IMG{2010012801-figure7.pdf}
\LANGen{
\Question{
Which of the following  formulas expresses the area of a regular pentagon inscribed in a circle of radius \( r \)? (See the picture.)}
{\( \frac{5r^2\sin72^{\circ}}2\)}{\( \frac{10r^2\sin72^{\circ}}2\)}{\( \frac{5r^2\sin36^{\circ}}2\)}{\( \frac{5r \sin36^{\circ}}2\)}{}{}}
\LANGpl{
\Question{
Który z poniższych wzorów wyraża pole pięciokąta foremnego wpisanego w okrąg o promieniu \( r \)? (Patrz rysunek.)}
{\( \frac{5r^2\sin72^{\circ}}2\)}{\( \frac{10r^2\sin72^{\circ}}2\)}{\( \frac{5r^2\sin36^{\circ}}2\)}{\( \frac{5r \sin36^{\circ}}2\)}{}{}}
\LANGcs{
\Question{
Který z následujících výrazů vyjadřuje obsah pravidelného pětiúhelníku vepsaného do kružnice o poloměru \( r \) (Viz obrázek.)?}
{\( \frac{5r^2\sin72^{\circ}}2\)}{\( \frac{10r^2\sin72^{\circ}}2\)}{\( \frac{5r^2\sin36^{\circ}}2\)}{\( \frac{5r \sin36^{\circ}}2\)}{}{}}
\LANGsk{
\Question{
Ktorý z uvedených vzorcov vyjadruje obsah pravidelného päťuholníka vpísaného do kružnice s polomerom \( r \)? (Pozri obrázok.)}
{\( \frac{5r^2\sin72^{\circ}}2\)}{\( \frac{10r^2\sin72^{\circ}}2\)}{\( \frac{5r^2\sin36^{\circ}}2\)}{\( \frac{5r \sin36^{\circ}}2\)}{}{}}
\LANGes{
\Question{
¿Cuál de las siguientes fórmulas expresa el área de un pentágono inscrito en una circunferencia con radio \( r \)?}
{\( \frac{5r^2\sin72^{\circ}}2\)}{\( \frac{10r^2\sin72^{\circ}}2\)}{\( \frac{5r^2\sin36^{\circ}}2\)}{\( \frac{5r \sin36^{\circ}}2\)}{}{}}
\End

\Data\ID{2010012901}
\Updated{2022-08-25 10:08:12}
\Level{B}
\SubArea{Circles}
\IMG{2010012901-figure0.pdf}
\LANGen{
\Question{
Consider a circle \( k \) with radius \( 5\,\mathrm{cm} \). In the circle is inscribed a convex quadrilateral \( ABCD \) so that the diagonal \( AC \) is the diameter of the circle, the length of \( BC \) is \( 8\,\mathrm{cm} \), and the length of \( DC \) is \( 5\,\mathrm{cm} \). Determine the length of side \( AD \). (See the picture.)}
{\(5 \sqrt{3}\,\mathrm{cm} \)}{\( 8\,\mathrm{cm} \)}{\( 10\,\mathrm{cm} \)}{\(3 \sqrt{5}\,\mathrm{cm} \)}{}{}}
\LANGpl{
\Question{
Rozważmy okrąg \( k \) o promieniu \( 5\,\mathrm{cm} \). W okrąg wpisany jest wypukły czworokąt \( ABCD \) tak, że przekątna \( AC \) jest średnicą koła, a długość \( BC \) wynosi \( 8\,\mathrm{cm} \) , zaś długość \( DC \) wynosi \( 5\,\mathrm{cm} \). Jaka jest długość boku \( AD \)? (Patrz rysunek.)}
{\(5 \sqrt{3}\,\mathrm{cm} \)}{\( 8\,\mathrm{cm} \)}{\( 10\,\mathrm{cm} \)}{\(3 \sqrt{5}\,\mathrm{cm} \)}{}{}}
\LANGcs{
\Question{
Je dána kružnice \( k \) s poloměrem \( 5\,\mathrm{cm} \). Do kružnice je vepsaný konvexní čtyřúhelník \( ABCD \) tak, že jeho úhlopříčka \( AC \) tvoří průměr kružnice, velikost strany \( BC \) je \( 8\,\mathrm{cm} \) a velikost \( DC \) je \( 5\,\mathrm{cm} \). Určete velikost strany \( AD \) (Viz obrázek.).}
{\(5 \sqrt{3}\,\mathrm{cm} \)}{\( 8\,\mathrm{cm} \)}{\( 10\,\mathrm{cm} \)}{\(3 \sqrt{5}\,\mathrm{cm} \)}{}{}}
\LANGsk{
\Question{
Daná je kružnica \( k \) s polomerom \( 5\,\mathrm{cm} \). Do tejto kružnice je vpísaný konvexný štvoruholník  \( ABCD \) tak, že uhlopriečkal \( AC \) je priemerom kružnice, dĺžka \( BC \) je \( 8\,\mathrm{cm} \), a dĺžka  \( DC \) je \( 5\,\mathrm{cm} \). Určte dĺžku strany \( AD \). (Pozri obrázok.)}
{\(5 \sqrt{3}\,\mathrm{cm} \)}{\( 8\,\mathrm{cm} \)}{\( 10\,\mathrm{cm} \)}{\(3 \sqrt{5}\,\mathrm{cm} \)}{}{}}
\LANGes{
\Question{
Dada la circunferencia \( k \) cuyo radio mide \( 5\,\mathrm{cm} \). En la circunferencia se inscribe el cuadrilátero convexo \( ABCD \) cuya diagonal \( AC \) es el diámetro de la circunferencia. La longitud del \( BC \) es \( 8\,\mathrm{cm} \), y la longitud del \( DC \) es \( 5\,\mathrm{cm} \). ¿Cuánto mide el lado \( AD \)?}
{\(5 \sqrt{3}\,\mathrm{cm} \)}{\( 8\,\mathrm{cm} \)}{\( 10\,\mathrm{cm} \)}{\(3 \sqrt{5}\,\mathrm{cm} \)}{}{}}
\End

\Data\ID{2010012902}
\Updated{2022-08-25 10:08:12}
\Level{B}
\SubArea{Circles}
\IMG{2010012901-figure1.pdf}
\LANGen{
\Question{
Derive the formula for calculating the area of an equilateral triangle inscribed in a circle of radius \( r \).}
{\(\frac{3\sqrt{3}r^2}{4} \)}{\(\frac{3\sqrt{3}r}{4} \)}{\(\frac{3\sqrt{3}r^2}{2} \)}{\(\frac{\sqrt{3}r^2}{4} \)}{}{}}
\LANGpl{
\Question{
Wyznacz wzór na obliczenie pola trójkąta równobocznego wpisanego w okrąg o promieniu \( r \).}
{\(\frac{3\sqrt{3}r^2}{4} \)}{\(\frac{3\sqrt{3}r}{4} \)}{\(\frac{3\sqrt{3}r^2}{2} \)}{\(\frac{\sqrt{3}r^2}{4} \)}{}{}}
\LANGcs{
\Question{
Odvoďte vztah pro výpočet obsahu rovnostranného trojúhelníku vepsaného do kružnice o poloměru \( r \).}
{\(\frac{3\sqrt{3}r^2}{4} \)}{\(\frac{3\sqrt{3}r}{4} \)}{\(\frac{3\sqrt{3}r^2}{2} \)}{\(\frac{\sqrt{3}r^2}{4} \)}{}{}}
\LANGsk{
\Question{
Odvoďte vzťah pre výpočet obsahu rovnostranného trojuholníka vpísaného do kružnice s polomerom  \( r \).}
{\(\frac{3\sqrt{3}r^2}{4} \)}{\(\frac{3\sqrt{3}r}{4} \)}{\(\frac{3\sqrt{3}r^2}{2} \)}{\(\frac{\sqrt{3}r^2}{4} \)}{}{}}
\LANGes{
\Question{
Averigua la fórmula para calcular el área de un triángulo equilátero inscrito en una circunferencia con el radio \( r \).}
{\(\frac{3\sqrt{3}r^2}{4} \)}{\(\frac{3\sqrt{3}r}{4} \)}{\(\frac{3\sqrt{3}r^2}{2} \)}{\(\frac{\sqrt{3}r^2}{4} \)}{}{}}
\End

\Data\ID{9000035002}
\Updated{2022-01-16 09:01:09}
\Level{B}
\SubArea{Circles}
\LANGen{
\Question{
A line segment of the length \(40\, \mathrm{cm}\) 
joins two points on a circle. The radius of the circle is
\(30\, \mathrm{cm}\). An
angle has the vertex in the center of the circle and the arms on the ends of the line
segment. Find the size of this angle and round the result to the nearest degrees and
minutes.}
{\(83^{\circ }37'\)}{\(97^{\circ }10'\)}{\(41^{\circ }48'\)}{\(96^{\circ }22'\)}{}{}}
\LANGpl{
\Question{
Odcinek prostej o długości  \(40\, \mathrm{cm}\) 
łączy dwa punkty na okręgu. Promień okręgu jest równy
\(30\, \mathrm{cm}\). Wierzchołek kąta znajduje się w środku okręgu, a jego ramiona przechodzą przez  końce odcinka prostej. 
Oblicz  ten kąt i zaokrągli odpowiedź do pełnych stopni i minut.}
{\(83^{\circ }37'\)}{\(97^{\circ }10'\)}{\(41^{\circ }48'\)}{\(96^{\circ }22'\)}{}{}}
\LANGcs{
\Question{
Tětiva v kružnici o poloměru \(30\, \mathrm{cm}\) 
má délku \(40\, \mathrm{cm}\).
Vypočítejte velikost středového úhlu příslušného této tětivě.
(Výsledek zaokrouhlete na celé stupně a minuty.)}
{\(83^{\circ }37'\)}{\(97^{\circ }10'\)}{\(41^{\circ }48'\)}{\(96^{\circ }22'\)}{}{}}
\LANGsk{
\Question{
Tetiva v kružnici s polomerom \(30\, \mathrm{cm}\) 
má dĺžku \(40\, \mathrm{cm}\).
Vypočítajte veľkosť stredového uhla prislúchajúceho tejto tetive.
(Výsledok zaokrúhlite na celé stupne a minúty.)}
{\(83^{\circ }37'\)}{\(97^{\circ }10'\)}{\(41^{\circ }48'\)}{\(96^{\circ }22'\)}{}{}}
\LANGes{
\Question{
Una cuerda en una circunferencia, cuyo radio es \(30\, \mathrm{cm}\), 
mide \(40\, \mathrm{cm}\). Calcula la medida del ángulo central de esta cuerda. Redondea el resultado a los grados y minutos más cercanos.}
{\(83^{\circ }37'\)}{\(97^{\circ }10'\)}{\(41^{\circ }48'\)}{\(96^{\circ }22'\)}{}{}}
\End

\Data\ID{9000045705}
\Updated{2020-07-15 03:07:34}
\Level{B}
\SubArea{Circles}
\IMG{000457_05-figure0.pdf}
\LANGen{
\Question{
Find the formula for the radius \(r\) 
of the circle circumscribed to the right triangle.}
{\(r = \frac{a} 
{2\cdot \sin \alpha } \)}{\(r = \frac{2a} 
 {\sin \alpha } \)}{\(r = \frac{a} 
{\sin 2\alpha } \)}{\(r = \frac{a} 
 {\sin \alpha } \)}{}{}}
\LANGpl{
\Question{
Wskaż wzór na promień \(r\) 
okręgu opisanego na trójkącie prostokątnym.}
{\(r = \frac{a} 
{2\cdot \sin \alpha } \)}{\(r = \frac{2a} 
 {\sin \alpha } \)}{\(r = \frac{a} 
{\sin 2\alpha } \)}{\(r = \frac{a} 
 {\sin \alpha } \)}{}{}}
\LANGcs{
\Question{
Vyberte vztah, který platí pro poloměr
\(r\) kružnice
\(k\) 
opsané pravoúhlému trojúhelníku ABC.}
{\(r = \frac{a} 
{2\cdot \sin \alpha } \)}{\(r = \frac{2a} 
 {\sin \alpha } \)}{\(r = \frac{a} 
{\sin 2\alpha } \)}{\(r = \frac{a} 
 {\sin \alpha } \)}{}{}}
\LANGsk{
\Question{
Vyberte vzťah, ktorý platí pre polomer
\(r\) kružnice
\(k\) 
opísanej pravouhlému trojuholníku ABC.}
{\(r = \frac{a} 
{2\cdot \sin \alpha } \)}{\(r = \frac{2a} 
 {\sin \alpha } \)}{\(r = \frac{a} 
{\sin 2\alpha } \)}{\(r = \frac{a} 
 {\sin \alpha } \)}{}{}}
\LANGes{
\Question{
Halla la fórmula para el radio \(r\) 
de la circunferencia circunscrita al triángulo rectángulo..}
{\(r = \frac{a} 
{2\cdot \sin \alpha } \)}{\(r = \frac{2a} 
 {\sin \alpha } \)}{\(r = \frac{a} 
{\sin 2\alpha } \)}{\(r = \frac{a} 
 {\sin \alpha } \)}{}{}}
\End

\Data\ID{9000045706}
\Updated{2020-07-15 03:07:34}
\Level{B}
\SubArea{Circles}
\LANGen{
\Question{
Given a regular pentagon with the side
\(a\), find the
radius \(r\) of
the circle circumscribed to this pentagon.}
{\(r = \frac{a} 
{2\cdot \cos 54^{\circ }}\)}{\(r = \frac{2a} 
{\cos 72^{\circ }}\)}{\(r = \frac{2a} 
{\cos 54^{\circ }}\)}{\(r = \frac{a} 
{2\cdot \cos 72^{\circ }}\)}{}{}}
\LANGpl{
\Question{
Dany jest pięciokąt równoboczny o boku 
\(a\), wskaż promień  \(r\) okręgu opisanego na tym pięciokącie.}
{\(r = \frac{a} 
{2\cdot \cos 54^{\circ }}\)}{\(r = \frac{2a} 
{\cos 72^{\circ }}\)}{\(r = \frac{2a} 
{\cos 54^{\circ }}\)}{\(r = \frac{a} 
{2\cdot \cos 72^{\circ }}\)}{}{}}
\LANGcs{
\Question{
Vyberte vztah, který platí pro poloměr
\(r\) 
kružnice opsané pravidelnému pětiúhelníku s délkou strany
\(a\).}
{\(r = \frac{a} 
{2\cdot \cos 54^{\circ }}\)}{\(r = \frac{2a} 
{\cos 72^{\circ }}\)}{\(r = \frac{2a} 
{\cos 54^{\circ }}\)}{\(r = \frac{a} 
{2\cdot \cos 72^{\circ }}\)}{}{}}
\LANGsk{
\Question{
Vyberte vzťah, ktorý platí pre polomer
\(r\) 
kružnice opísanej pravidelnému päťuholníku s dĺžkou strany
\(a\).}
{\(r = \frac{a} 
{2\cdot \cos 54^{\circ }}\)}{\(r = \frac{2a} 
{\cos 72^{\circ }}\)}{\(r = \frac{2a} 
{\cos 54^{\circ }}\)}{\(r = \frac{a} 
{2\cdot \cos 72^{\circ }}\)}{}{}}
\LANGes{
\Question{
Dado un pentágono regular con el lado
\(a\), halla el radio \(r\) de la circunferencia circunscrita al pentágono.}
{\(r = \frac{a} 
{2\cdot \cos 54^{\circ }}\)}{\(r = \frac{2a} 
{\cos 72^{\circ }}\)}{\(r = \frac{2a} 
{\cos 54^{\circ }}\)}{\(r = \frac{a} 
{2\cdot \cos 72^{\circ }}\)}{}{}}
\End

\Data\ID{9000045707}
\Updated{2020-07-15 03:07:34}
\Level{B}
\SubArea{Circles}
\LANGen{
\Question{
Given a regular pentagon with the side
\(a\), find the
radius \(\rho \) of
the circle inscribed to this pentagon.}
{\(\rho = \frac{a} 
{2} \cdot \mathop{\mathrm{tg}}\nolimits 54^{\circ }\)}{\(\rho = \frac{2a} 
{\mathop{\mathrm{tg}}\nolimits 54^{\circ }}\)}{\(\rho = \frac{a} 
{2\cdot \mathop{\mathrm{tg}}\nolimits 54^{\circ }}\)}{\(\rho = 2a\cdot \mathop{\mathrm{tg}}\nolimits 54^{\circ }\)}{}{}}
\LANGpl{
\Question{
Dany jest pięciokąt równoboczny o boku
\(a\), wskaż promień \(\rho \) okręgu wpisanego w ten pięciokąt.}
{\(\rho = \frac{a} 
{2} \cdot \mathop{\mathrm{tg}}\nolimits 54^{\circ }\)}{\(\rho = \frac{2a} 
{\mathop{\mathrm{tg}}\nolimits 54^{\circ }}\)}{\(\rho = \frac{a} 
{2\cdot \mathop{\mathrm{tg}}\nolimits 54^{\circ }}\)}{\(\rho = 2a\cdot \mathop{\mathrm{tg}}\nolimits 54^{\circ }\)}{}{}}
\LANGcs{
\Question{
Vyberte vztah, který platí pro poloměr
\(\rho \) 
kružnice vepsané pravidelnému pětiúhelníku s délkou strany
\(a\).}
{\(\rho = \frac{a} 
{2} \cdot \mathop{\mathrm{tg}}\nolimits 54^{\circ }\)}{\(\rho = \frac{2a} 
{\mathop{\mathrm{tg}}\nolimits 54^{\circ }}\)}{\(\rho = \frac{a} 
{2\cdot \mathop{\mathrm{tg}}\nolimits 54^{\circ }}\)}{\(\rho = 2a\cdot \mathop{\mathrm{tg}}\nolimits 54^{\circ }\)}{}{}}
\LANGsk{
\Question{
Vyberte vzťah, ktorý platí pre polomer
\(\rho \) 
kružnice vpísanej do pravidelného päťuholníka s dĺžkou strany
\(a\).}
{\(\rho = \frac{a} 
{2} \cdot \mathop{\mathrm{tg}}\nolimits 54^{\circ }\)}{\(\rho = \frac{2a} 
{\mathop{\mathrm{tg}}\nolimits 54^{\circ }}\)}{\(\rho = \frac{a} 
{2\cdot \mathop{\mathrm{tg}}\nolimits 54^{\circ }}\)}{\(\rho = 2a\cdot \mathop{\mathrm{tg}}\nolimits 54^{\circ }\)}{}{}}
\LANGes{
\Question{
Dado un pentágono regular con el lado
\(a\), halla el radio \(\rho \) de la circunferencia inscrita en el pentágono.}
{\(\rho = \frac{a} 
{2} \cdot \mathop{\mathrm{tg}}\nolimits 54^{\circ }\)}{\(\rho = \frac{2a} 
{\mathop{\mathrm{tg}}\nolimits 54^{\circ }}\)}{\(\rho = \frac{a} 
{2\cdot \mathop{\mathrm{tg}}\nolimits 54^{\circ }}\)}{\(\rho = 2a\cdot \mathop{\mathrm{tg}}\nolimits 54^{\circ }\)}{}{}}
\End

\Data\ID{9000045708}
\Updated{2020-07-15 03:07:34}
\Level{B}
\SubArea{Circles}
\LANGen{
\Question{
Given a regular hexagon with the side
\(a\), find the
radius \(\rho \) of
the circle inscribed to this hexagon.}
{\(\rho = \frac{a} 
{2\cdot \mathop{\mathrm{tg}}\nolimits 30^{\circ }}\)}{\(\rho = 2a\cdot \mathop{\mathrm{tg}}\nolimits 30^{\circ }\)}{\(\rho = \frac{2a} 
{\mathop{\mathrm{tg}}\nolimits 30^{\circ }}\)}{\(\rho = 2a\cdot \mathop{\mathrm{tg}}\nolimits 60^{\circ }\)}{}{}}
\LANGpl{
\Question{
Dany jest sześciokąt równoboczny o boku 
\(a\), wskaż promień  \(\rho \) okręgu  wpisanego w ten sześciokąt.}
{\(\rho = \frac{a} 
{2\cdot \mathop{\mathrm{tg}}\nolimits 30^{\circ }}\)}{\(\rho = 2a\cdot \mathop{\mathrm{tg}}\nolimits 30^{\circ }\)}{\(\rho = \frac{2a} 
{\mathop{\mathrm{tg}}\nolimits 30^{\circ }}\)}{\(\rho = 2a\cdot \mathop{\mathrm{tg}}\nolimits 60^{\circ }\)}{}{}}
\LANGcs{
\Question{
Vyberte vztah, který platí pro poloměr
\(\rho \) 
kružnice vepsané pravidelnému šestiúhelníku s délkou strany
\(a\).}
{\(\rho = \frac{a} 
{2\cdot \mathop{\mathrm{tg}}\nolimits 30^{\circ }}\)}{\(\rho = 2a\cdot \mathop{\mathrm{tg}}\nolimits 30^{\circ }\)}{\(\rho = \frac{2a} 
{\mathop{\mathrm{tg}}\nolimits 30^{\circ }}\)}{\(\rho = 2a\cdot \mathop{\mathrm{tg}}\nolimits 60^{\circ }\)}{}{}}
\LANGsk{
\Question{
Vyberte vzťah, ktorý platí pre polomer
\(\rho \) 
kružnice vpísanej do pravidelného šesťuholníka s dĺžkou strany
\(a\).}
{\(\rho = \frac{a} 
{2\cdot \mathop{\mathrm{tg}}\nolimits 30^{\circ }}\)}{\(\rho = 2a\cdot \mathop{\mathrm{tg}}\nolimits 30^{\circ }\)}{\(\rho = \frac{2a} 
{\mathop{\mathrm{tg}}\nolimits 30^{\circ }}\)}{\(\rho = 2a\cdot \mathop{\mathrm{tg}}\nolimits 60^{\circ }\)}{}{}}
\LANGes{
\Question{
Dado un hexágono regular con el lado
\(a\), halla el radio \(\rho \) de la circunferencia inscrita en el hexágono.}
{\(\rho = \frac{a} 
{2\cdot \mathop{\mathrm{tg}}\nolimits 30^{\circ }}\)}{\(\rho = 2a\cdot \mathop{\mathrm{tg}}\nolimits 30^{\circ }\)}{\(\rho = \frac{2a} 
{\mathop{\mathrm{tg}}\nolimits 30^{\circ }}\)}{\(\rho = 2a\cdot \mathop{\mathrm{tg}}\nolimits 60^{\circ }\)}{}{}}
\End

\Data\ID{9000046405}
\Updated{2022-01-16 09:01:35}
\Level{B}
\SubArea{Circles}
\IMG{000464_05-figure0.pdf}
\LANGen{
\Question{
A circle is circumscribed to the regular octagon. The perimeter of the octagon is
\(16\, \mathrm{cm}\). Find
the radius of the circle and round the result to two decimal places. (The regular
octagon is a polygon which has eight sides of equal length. The perimeter of the
octagon is the sum of the length of all eight sides.) Circle circumscribed to the regular octagon.}
{\(2.61\, \mathrm{cm}\)}{\(1.08\, \mathrm{cm}\)}{\(1.41\, \mathrm{cm}\)}{}{}{}}
\LANGpl{
\Question{
Oblicz promień okręgu opisanego na ośmiokącie równobocznym o obwodzie  
\(16\, \mathrm{cm}\).  (Wynik jest zaokrąglany do 2 miejsc po przecinku.)}
{\(2.61\, \mathrm{cm}\)}{\(1.08\, \mathrm{cm}\)}{\(1.41\, \mathrm{cm}\)}{}{}{}}
\LANGcs{
\Question{
Určete poloměr kružnice opsané pravidelnému osmiúhelníku o obvodu
\(16\, \mathrm{cm}\) (výsledek je
zaokrouhlen na \(2\) 
desetinná místa).}
{\(2{,}61\, \mathrm{cm}\)}{\(1{,}08\, \mathrm{cm}\)}{\(1{,}41\, \mathrm{cm}\)}{}{}{}}
\LANGsk{
\Question{
Určte polomer kružnice opísanej pravidelnému osemuholníku s obvodom
\(16\, \mathrm{cm}\) (výsledok je
zaokrúhlený na \(2\) 
desatinné miesta).}
{\(2{,}61\, \mathrm{cm}\)}{\(1{,}08\, \mathrm{cm}\)}{\(1{,}41\, \mathrm{cm}\)}{}{}{}}
\LANGes{
\Question{
Una circunferencia está circunscrita en un octógono regular. El perímetro del octógono mide
\(16\, \mathrm{cm}\). Calcula el radio de la circunferencia circunscrita y redondea el resultado a dos decimales.}
{\(2.61\, \mathrm{cm}\)}{\(1.08\, \mathrm{cm}\)}{\(1.41\, \mathrm{cm}\)}{}{}{}}
\End

\Data\ID{1103256901}
\Updated{2022-01-16 09:01:44}
\Level{C}
\SubArea{Circles}
\IMG{1103256901.pdf}
\LANGen{
\Question{
The farmer tied two goats on the meadow. The distance of the stakes \( K_1 \), \( K_2 \) to which the goats are tied is \( 5\,\mathrm{m} \) and the ropes have lengths of \( 3\,\mathrm{m} \) and \( 4\,\mathrm{m} \). What is the area of the grassland which is common for both goats? Round the result to two decimal places.}
{\( 6.64\,\mathrm{m}^2 \)}{\( 0.57\,\mathrm{m}^2 \)}{\( 0.35\,\mathrm{m}^2 \)}{\( 1.52\,\mathrm{m}^2 \)}{}{}}
\LANGpl{
\Question{
Farmer przywiązał dwie kozy na łące. Odległość pomiędzy kołkami  \( K_1 \), \( K_2 \), do których przywiązane są kozy wynosi \( 5\,\mathrm{m} \), a długości lin są równe \( 3\,\mathrm{m} \) i \( 4\,\mathrm{m} \). Jaka jest powierzchnia łąki wspólna dla  obydwu kóz? Zaokrąglij do miejsc dziesiętnych.}
{\( 6{,}64\,\mathrm{m}^2 \)}{\( 0{,}57\,\mathrm{m}^2 \)}{\( 0{,}35\,\mathrm{m}^2 \)}{\( 1{,}52\,\mathrm{m}^2 \)}{}{}}
\LANGcs{
\Question{
Farmář uvázal na louku dvě kozy. Vzdálenost kolíků \(K_1\), \(K_2 \), ke kterým jsou kozy uvázané, je \(5 \, \mathrm {m}\) a lana mají délku \(3 \, \mathrm {m}\)  a \(4 \, \mathrm {m}\) . Jakou plochu má pastvina, která je společná pro obě kozy? Zaokrouhlete na dvě desetinná místa.}
{\( 6{,}64\,\mathrm{m}^2 \)}{\( 0{,}57\,\mathrm{m}^2 \)}{\( 0{,}35\,\mathrm{m}^2 \)}{\( 1{,}52\,\mathrm{m}^2 \)}{}{}}
\LANGsk{
\Question{
Farmár uviazal na lúku dve kozy. Vzdialenosť kolíkov \(K_1\), \(K_2 \), ku ktorým sú kozy uviazané, je \(5 \, \mathrm {m}\) a laná majú dĺžky \(3 \, \mathrm {m}\)  a \(4 \, \mathrm {m}\) . Akú plochu má paša, ktorá je spoločná pre obe kozy? Zaokrúhlite na dve desatinné miesta.}
{\( 6{,}64\,\mathrm{m}^2 \)}{\( 0{,}57\,\mathrm{m}^2 \)}{\( 0{,}35\,\mathrm{m}^2 \)}{\( 1{,}52\,\mathrm{m}^2 \)}{}{}}
\LANGes{
\Question{
El agricultor ató dos cabras en un prado. La distancia entre las estacas \( K_1 \), \( K_2 \) es \( 5\,\mathrm{m} \) y las cuerdas miden \( 3\,\mathrm{m} \) y \( 4\,\mathrm{m} \). Calcula el área del pasto que es común para ambas cabras. Redondea el resultado a dos decimales.}
{\( 6.64\,\mathrm{m}^2 \)}{\( 0.57\,\mathrm{m}^2 \)}{\( 0.35\,\mathrm{m}^2 \)}{\( 1.52\,\mathrm{m}^2 \)}{}{}}
\End

\Data\ID{1103256902}
\Updated{2020-07-15 03:07:21}
\Level{C}
\SubArea{Circles}
\IMG{1103256902.pdf}
\LANGen{
\Question{
The cucumber field has the shape of an isosceles right triangle. The length of its legs is \( 12\,\mathrm{m} \). Rotary sprinklers placed in its vertices have a reach of \( 6\,\mathrm{m} \). Find the area of the field that is not sprinkled with water. Round the result  to two decimal places.}
{\( 15.45\,\mathrm{m}^2 \)}{\( 41.10\,\mathrm{m}^2 \)}{\( 16.29\,\mathrm{m}^2 \)}{\( 15.25\,\mathrm{m}^2 \)}{}{}}
\LANGpl{
\Question{
Pole ogórków ma kształt trójkąta równoramiennego. Długość jego ramion wynosi  \( 12\,\mathrm{m} \). Rozpryskiwocze ustawione na jego wierzchołkach mają zasięg \( 6\,\mathrm{m} \). Wyznacz powierzchnię pola, która nie jest spryskiwana. Zaokrąglij do dwóch  miejsc dziesiętnych.}
{\( 15{,}45\,\mathrm{m}^2 \)}{\( 41{,}10\,\mathrm{m}^2 \)}{\( 16{,}29\,\mathrm{m}^2 \)}{\( 15{,}25\,\mathrm{m}^2 \)}{}{}}
\LANGcs{
\Question{
Okurkové pole má tvar pravoúhlého rovnoramenného trojúhelníka s délkou odvěsny \( 12\,\mathrm{m} \). Ve vrcholech trojúhelníka jsou umístěné otáčivé rozprašovače s dosahem \( 6\,\mathrm{m} \). Jak velkou část pole tyto rozstřikovače nezavlažují? Výsledek zaokrouhlete na dvě desetinná místa.}
{\( 15{,}45\,\mathrm{m}^2 \)}{\( 41{,}10\,\mathrm{m}^2 \)}{\( 16{,}29\,\mathrm{m}^2 \)}{\( 15{,}25\,\mathrm{m}^2 \)}{}{}}
\LANGsk{
\Question{
Uhorkové pole má tvar pravouhlého rovnoramenného trojuholníka s dĺžkou odvesny \( 12\,\mathrm{m} \). Vo vrcholoch trojuholníka sú umiestnené otáčacie postrekovače s dosahom \( 6\,\mathrm{m} \). Akú veľkú časť poľa tieto postrekovače nezavlažujú? Výsledok zaokrúhlite na dve desatinné miesta.}
{\( 15{,}45\,\mathrm{m}^2 \)}{\( 41{,}10\,\mathrm{m}^2 \)}{\( 16{,}29\,\mathrm{m}^2 \)}{\( 15{,}25\,\mathrm{m}^2 \)}{}{}}
\LANGes{
\Question{
Un campo, donde se cultivan pepinos, tiene forma de triángulo rectángulo isósceles. Sus catetos miden \( 12\,\mathrm{m} \). En los vértices del triángulo se sitúan aspersores rotativos con el alcance de \( 6\,\mathrm{m} \). Calcula el área del campo que no está rociada con agua. Redondea el resultado a dos decimales.}
{\( 15.45\,\mathrm{m}^2 \)}{\( 41.10\,\mathrm{m}^2 \)}{\( 16.29\,\mathrm{m}^2 \)}{\( 15.25\,\mathrm{m}^2 \)}{}{}}
\End

\Data\ID{1103256903}
\Updated{2022-01-16 09:01:10}
\Level{C}
\SubArea{Circles}
\IMG{1103256903.pdf}
\LANGen{
\Question{
In isosceles triangle \( ABC \), \( |AB| = 8\,\mathrm{cm} \), \( |BC|=|AC| = 6\,\mathrm{cm} \). Determine what percentage of the triangle area is a circle that is inscribed in it. Round the result to full percentages.}
{\( 56\,\% \)}{\( 48\,\% \)}{\( 62\,\% \)}{\( 64\,\% \)}{}{}}
\LANGpl{
\Question{
W trójkącie równoramiennym \( ABC \), \( |AB| = 8\,\mathrm{cm} \), \( |BC|=|AC| = 6\,\mathrm{cm} \). Oblicz jaki procent powierzchni trójkąta stanowi okrąg w niego wpisany. Zaokrąglij do pełnych procentów.}
{\( 56\,\% \)}{\( 48\,\% \)}{\( 62\,\% \)}{\( 64\,\% \)}{}{}}
\LANGcs{
\Question{
V rovnoramenném trojúhelníku \( ABC \), \( |AB| = 8\,\mathrm{cm} \), \( |BC|=|AC| = 6\,\mathrm{cm} \). Do trojúhelníku je vepsaný kruh. Zjistěte, kolik procent z obsahu trojúhelníka tvoří obsah vepsaného kruhu. Výsledek zaokrouhlete na celá procenta.}
{\( 56\,\% \)}{\( 48\,\% \)}{\( 62\,\% \)}{\( 64\,\% \)}{}{}}
\LANGsk{
\Question{
V rovnoramennom trojuholníku  \( ABC \), \( |AB| = 8\,\mathrm{cm} \), \( |BC|=|AC| = 6\,\mathrm{cm} \). Do trojuholníka je vpísaný kruh. Zistite, koľko percent z obsahu trojuholníka tvorí obsah vpísaného kruhu. Výsledok zaokrúhlite na celé percentá.}
{\( 56\,\% \)}{\( 48\,\% \)}{\( 62\,\% \)}{\( 64\,\% \)}{}{}}
\LANGes{
\Question{
Dado el triángulo isósceles \( ABC \), \( |AB| = 8\,\mathrm{cm} \), \( |BC|=|AC| = 6\,\mathrm{cm} \). En el triángulo se inscribe un círculo. Calcula, qué porcentaje del área del triángulo pertenece al círculo inscrito. Redondea el resultado a un porcentaje determinado.}
{\( 56\,\% \)}{\( 48\,\% \)}{\( 62\,\% \)}{\( 64\,\% \)}{}{}}
\End

\Data\ID{9000036104}
\Updated{2020-07-15 03:07:34}
\Level{C}
\SubArea{Circles}
\LANGen{
\Question{
The angles in the triangle \(ABC\) 
are \(\alpha = 100^{\circ }\) and
\(\beta = 50^{\circ }\). The radius of the
circumscribed circle is \(11\, \mathrm{cm}\).
Find the side \(c\).}
{\(11\, \mathrm{cm}\)}{\(8\, \mathrm{cm}\)}{\(9\, \mathrm{cm}\)}{\(10\, \mathrm{cm}\)}{}{}}
\LANGpl{
\Question{
Dany jest trójkąt  \(ABC\)  o kątach 
 \(\alpha = 100^{\circ }\) i
\(\beta = 50^{\circ }\). Promień okręgu opisanego na  trójkącie  wynosi  \(11\, \mathrm{cm}\).
Oblicz bok  \(c\).}
{\(11\, \mathrm{cm}\)}{\(8\, \mathrm{cm}\)}{\(9\, \mathrm{cm}\)}{\(10\, \mathrm{cm}\)}{}{}}
\LANGcs{
\Question{
Vypočítejte délku strany \(c\) 
v trojúhelníku \(ABC\),
je-li úhel \(\alpha = 100^{\circ }\) a úhel
\(\beta = 50^{\circ }\). Poloměr kružnice
opsané trojúhelníku \(ABC\) 
je \(11\, \mathrm{cm}\).}
{\(11\, \mathrm{cm}\)}{\(8\, \mathrm{cm}\)}{\(9\, \mathrm{cm}\)}{\(10\, \mathrm{cm}\)}{}{}}
\LANGsk{
\Question{
Vypočítajte dĺžku strany \(c\) 
v trojuholníku \(ABC\),
ak je uhol \(\alpha = 100^{\circ }\) a uhol
\(\beta = 50^{\circ }\). Polomer kružnice
opísanej trojuholníku \(ABC\) 
je \(11\, \mathrm{cm}\).}
{\(11\, \mathrm{cm}\)}{\(8\, \mathrm{cm}\)}{\(9\, \mathrm{cm}\)}{\(10\, \mathrm{cm}\)}{}{}}
\LANGes{
\Question{
Los ángulos en el triángulo \(ABC\) 
son \(\alpha = 100^{\circ }\) y
\(\beta = 50^{\circ }\). El radio de la circunferencia circunscrita es \(11\, \mathrm{cm}\).
Halla el lado \(c\).}
{\(11\, \mathrm{cm}\)}{\(8\, \mathrm{cm}\)}{\(9\, \mathrm{cm}\)}{\(10\, \mathrm{cm}\)}{}{}}
\End

\Data\ID{9000036105}
\Updated{2020-07-15 03:07:34}
\Level{C}
\SubArea{Circles}
\LANGen{
\Question{
The side \(b\) in
the triangle \(ABC\) 
is \(17\, \mathrm{cm}\) and
the angle \(\beta \) 
is \(58^{\circ }\).
Find the radius of the circle circumscribed to this triangle and round your answer to
the nearest centimeters.}
{\(10\, \mathrm{cm}\)}{\(8\, \mathrm{cm}\)}{\(9\, \mathrm{cm}\)}{\(11\, \mathrm{cm}\)}{}{}}
\LANGpl{
\Question{
Dany jest trójkąt \(ABC\) o boku  \(b=17\, \mathrm{cm}\) i kącie \(\beta= 58^{\circ }\).
Oblicz  promień okręgu opisanego na trójkącie i zaokrąglij odpowiedź na całe centymetry.}
{\(10\, \mathrm{cm}\)}{\(8\, \mathrm{cm}\)}{\(9\, \mathrm{cm}\)}{\(11\, \mathrm{cm}\)}{}{}}
\LANGcs{
\Question{
Určete poloměr kružnice opsané trojúhelníku
\(ABC\), je-li
strana \(b = 17\, \mathrm{cm}\) a
úhel \(\beta = 58^{\circ }\).
Výsledek zaokrouhlete na celé centimetry.}
{\(10\, \mathrm{cm}\)}{\(8\, \mathrm{cm}\)}{\(9\, \mathrm{cm}\)}{\(11\, \mathrm{cm}\)}{}{}}
\LANGsk{
\Question{
Určte polomer kružnice opísanej trojuholníku
\(ABC\), ak
strana \(b = 17\, \mathrm{cm}\) a
uhol \(\beta = 58^{\circ }\).
Výsledok zaokrúhlite na celé centimetre.}
{\(10\, \mathrm{cm}\)}{\(8\, \mathrm{cm}\)}{\(9\, \mathrm{cm}\)}{\(11\, \mathrm{cm}\)}{}{}}
\LANGes{
\Question{
El lado \(b\) en el triángulo \(ABC\) 
mide \(17\, \mathrm{cm}\) y
el ángulo \(\beta \) 
mide \(58^{\circ }\).
Calcula el radio de la circunferencia circunscrita al triángulo.}
{\(10\, \mathrm{cm}\)}{\(8\, \mathrm{cm}\)}{\(9\, \mathrm{cm}\)}{\(11\, \mathrm{cm}\)}{}{}}
\End
